\chapter{Rozmaitości różniczkowe}
\label{ch:rozmaitosci}
\index{rozmaitosc@Rozmaitość}

\noindent\textbf{Jak czytać skrypt.}
Wystarczą podstawy algebry liniowej, rachunku różniczkowego
wielu zmiennych i topologii. Każdy nowy obiekt najpierw
opisujemy intuicyjnie, potem definiujemy i używamy na
przykładach. W dowodach rozpisujemy kroki specyficzne dla
rozmaitości; gdy potrzebne jest duże twierdzenie z analizy
lub topologii (np.\ twierdzenie o funkcji odwrotnej albo
niezmienniczość obszaru), wskazujemy je wprost.

Niniejszy rozdział stanowi wprowadzenie do teorii rozmaitości gładkich — podstawowego
obiektu geometrii i topologii różniczkowej. Zaczynamy od pojęcia rozmaitości topologicznej,
następnie wprowadzamy strukturę różniczkową za pomocą atlasów, i budujemy krok po kroku
aparat niezbędny do uprawiania analizy na rozmaitościach: przestrzeń styczną i kostyczną,
wiązki wektorowe oraz różniczkowanie odwzorowań między rozmaitościami wraz z klasyfikacją
ich lokalnego zachowania (immersje, submersje, zanurzenia) i teorią podrozmaitości, na
której ten rozdział się kończy. Zbudowany tu aparat stanowi fundament, na którym w
kolejnych rozdziałach nadbudujemy dalszą teorię geometrii i topologii różniczkowej.



\section{Rozmaitości topologiczne}

Rozmaitość topologiczna to przestrzeń topologiczna, która lokalnie ,,wygląda'' jak
przestrzeń euklidesowa $\R^n$. Precyzujemy to poniżej.

\begin{definicja}
Przestrzeń topologiczna $M$ jest \emph{lokalnie euklidesowa wymiaru $n$}, jeżeli każdy
punkt $p\in M$ posiada otoczenie otwarte $U$ homeomorficzne z otwartym podzbiorem $\R^n$.
\end{definicja}

\begin{definicja}
\emph{Rozmaitością topologiczną} wymiaru $n$ nazywamy przestrzeń topologiczną $M$, która:
\begin{enumerate}[label=(\roman*)]
\item jest przestrzenią Hausdorffa,
\item spełnia drugi aksjomat przeliczalności (posiada przeliczalną bazę topologii),
\item jest lokalnie euklidesowa wymiaru $n$.
\end{enumerate}
\end{definicja}

Przyjmujemy $n\geq0$. Jeśli wyraźnie nie zaznaczono inaczej,
rozmaitości rozważamy bez brzegu. Rozmaitości z brzegiem
wprowadzimy osobno w Sekcji~\ref{sec:rozmaitosci-z-brzegiem-1};
ta domyślna konwencja obowiązuje także po ich zdefiniowaniu.
Przypomnijmy znaczenie dwóch warunków z definicji.
\emph{Hausdorffa} znaczy, że dla dowolnych różnych punktów
$p,q$ można znaleźć rozłączne otwarte otoczenia $U\ni p$
i $V\ni q$. \emph{Przeliczalna baza topologii} to rodzina
$B_1,B_2,\ldots$ zbiorów otwartych taka, że każdy zbiór
otwarty jest sumą pewnych $B_i$. Lokalna euklidesowość
dotyczy \emph{małego otoczenia każdego punktu}; nie twierdzi,
że całą przestrzeń można opisać jedną mapą do $\R^n$.

Warunki (i) i (ii) mogą wydawać się techniczne, jednak są istotne i nie wynikają
automatycznie z lokalnej euklidesowości — poniżej pokazujemy stosowne kontrprzykłady.

\begin{uwaga}[Niezmienniczość wymiaru]
Wymiar $n$ w powyższej definicji jest wyznaczony jednoznacznie przez $M$: jeśli niepusty
zbiór otwarty $U\subseteq\R^n$ jest homeomorficzny ze zbiorem otwartym $V\subseteq\R^m$, to
$n=m$. Jest to \emph{niezmienniczość wymiaru}, będąca konsekwencją
\emph{twierdzenia Brouwera o niezmienniczości obszaru}: ciągła injekcja
z otwartego podzbioru $\R^d$ do $\R^d$ ma otwarty obraz.
Samo twierdzenie Brouwera przyjmujemy bez dowodu; wymaga narzędzi
topologii algebraicznej. Wyprowadźmy jednak potrzebny wniosek.
Gdyby $n<m$ i $h:U\to V$ było homeomorfizmem, to złożenie
$V\xrightarrow{h^{-1}}U\hookrightarrow\R^m$, gdzie ostatnie
odwzorowanie dopisuje $m-n$ zer, byłoby ciągłą injekcją.
Jej obraz musiałby być otwarty w $\R^m$, lecz leży w podprzestrzeni
$\R^n\times\{0\}$ o pustym wnętrzu: sprzeczność.
Przypadek $m<n$ wykluczamy, zamieniając role $U,V$.
Wymiar jest więc jednoznaczny dla niepustej rozmaitości.
Pusty zbiór dopuszczamy jako rozmaitość dowolnego ustalonego wymiaru.
\end{uwaga}

\begin{definicja}[Nasycenie zbioru]\label{def:nasycenie}
\index{nasycenie zbioru@Nasycenie zbioru}
Niech $\sim$ będzie relacją równoważności na zbiorze $X$,
a $\pi\colon X\to X/{\sim}$ rzutowaniem na zbiór klas.
\emph{Nasyceniem} zbioru $A\subseteq X$ względem tej relacji
nazywamy
\[
 \operatorname{Sat}_{\sim}(A)
 :=\pi^{-1}(\pi(A))
 =\bigcup_{x\in A}[x].
\]
Zbiór $A$ jest \emph{nasycony}, jeśli
$\operatorname{Sat}_{\sim}(A)=A$, czyli razem z każdym
$x\in A$ zawiera całą jego klasę $[x]$.
Równoważnie, $A=\pi^{-1}(B)$ dla pewnego
$B\subseteq X/{\sim}$. Nasycenie jest najmniejszym
zbiorem nasyconym zawierającym $A$.
\end{definicja}

Ta sama definicja działa dla dowolnej surjekcji
$q\colon X\to Y$: nasycenie względem jej włókien to
$q^{-1}(q(A))$. Dla przykładu w $\R/{\Z}$, gdzie
$x\sim y$ oznacza $x-y\in\Z$, nasycenie
$A=(0,\tfrac12)$ jest sumą
$\bigcup_{k\in\Z}(k,k+\tfrac12)$.

\begin{definicja}[Topologia ilorazowa]\label{def:topologia-ilorazowa}
\index{topologia ilorazowa@Topologia ilorazowa}
Niech $X$ będzie przestrzenią topologiczną, $\sim$ relacją równoważności
na $X$, a $\pi\colon X\to X/{\sim}$ rzutowaniem
$x\mapsto[x]$. \emph{Topologia ilorazowa} na zbiorze klas $X/{\sim}$
jest zadana warunkiem
\[
 V\subseteq X/{\sim}\ \text{jest otwarty}
 \quad\Longleftrightarrow\quad
 \pi^{-1}(V)\subseteq X\ \text{jest otwarty}.
\]
Przeciwobraz dowolnego $V\subseteq X/{\sim}$ jest
automatycznie nasycony w sensie
Definicji~\ref{def:nasycenie}. Zatem otwarte zbiory
ilorazu odpowiadają dokładnie otwartym
i nasyconym podzbiorom $X$.
Rzutowanie $\pi$ jest ciągłe i surjektywne.
\end{definicja}

Intuicyjnie sklejone punkty stają się jednym punktem ilorazu.
Dlatego przy sprawdzaniu otwartości musimy uwzględnić
\emph{wszystkie} punkty, które w ilorazie zostają utożsamione:
temu służy nasycenie. Sam obraz zbioru otwartego
$A\subseteq X$ nie musi być otwarty w ilorazie,
ponieważ $\pi^{-1}(\pi(A))$ może zawierać dodatkowe
punkty spoza $A$ i nie być zbiorem otwartym.
Na przykład, jeśli w $\R$ skleimy tylko $0$ z $1$,
to dla $A=(-\tfrac14,\tfrac14)$ nasycenie ma postać
$A\cup\{1\}$, która nie jest otwarta; $\pi(A)$ też
nie jest otwarty w ilorazie.

W szczególności dla dowolnego $A\subseteq X$ otrzymujemy
\[
 \pi(A)\text{ jest otwarty w }X/{\sim}
 \quad\Longleftrightarrow\quad
 \operatorname{Sat}_{\sim}(A)\text{ jest otwarty w }X.
\]
To pozwala sprawdzać otwartość zbiorów po sklejeniu.
Warunek z definicji rzeczywiście wyznacza topologię: przeciwobraz
zachowuje sumy dowolnych rodzin oraz skończone przecięcia.
Jest to najbogatsza topologia na $X/{\sim}$, dla której $\pi$ jest
ciągłe. Ma też użyteczną własność: dla dowolnej przestrzeni
topologicznej $Z$ odwzorowanie $g\colon X/{\sim}\to Z$ jest ciągłe
wtedy i tylko wtedy, gdy $g\circ\pi\colon X\to Z$ jest ciągłe.
Istotnie, dla każdego otwartego $W\subseteq Z$ zachodzi
$(g\circ\pi)^{-1}(W)=\pi^{-1}(g^{-1}(W))$, więc zastosowanie definicji
do $g^{-1}(W)$ daje obie implikacje. Będziemy z tej topologii
korzystać przy sklejaniu przestrzeni oraz przy przestrzeniach
rzutowych.

\begin{przyklad}[Prosta z podwojonym początkiem]
Pokażemy, że warunku Hausdorffa nie można pominąć w definicji rozmaitości topologicznej. Istnieją bowiem przestrzenie, które są lokalnie euklidesowe i spełniają drugi aksjomat przeliczalności, ale nie są przestrzeniami Hausdorffa.

Rozważmy sumę rozłączną dwóch kopii prostej rzeczywistej $X=(\mathbb{R}\times\{0\})\sqcup(\mathbb{R}\times\{1\})$. Punkty pierwszej kopii będziemy oznaczać przez $(x,0)$, a drugiej przez $(x,1)$.

Na $X$ wprowadzamy relację równoważności $\sim$, która utożsamia odpowiadające sobie punkty obu prostych z wyjątkiem zera. Definiujemy ją przez
$(x,i)\sim(y,j)\Longleftrightarrow x=y\ \text{oraz}\ \bigl(i=j\ \text{lub}\ x\neq0\bigr)$.
Innymi słowy, dla każdego $x\neq0$ zachodzi $(x,0)\sim(x,1)$, natomiast punkty $(0,0)$ oraz $(0,1)$ nie są równoważne.

Przestrzeń ilorazową oznaczamy przez $M=X/\sim$, a przez $\pi\colon X\to M$ naturalną projekcję ilorazową. Dla $x\neq0$ klasa równoważności ma postać $[x]=\{(x,0),(x,1)\}$. Ponadto $[(0,0)]=\{(0,0)\}$ oraz $[(0,1)]=\{(0,1)\}$. Oznaczmy $0_a=[(0,0)]$ oraz $0_b=[(0,1)]$. Są to dwa różne punkty przestrzeni $M$.

Można więc intuicyjnie myśleć o $M$ jako o zwykłej prostej rzeczywistej, w której punkt $0$ został zastąpiony przez dwa różne punkty $0_a$ i $0_b$. Poza zerem obie kopie prostej zostały całkowicie sklejone.

Opiszmy dokładniej otoczenia punktów $0_a$ i $0_b$. Dla każdego otwartego przedziału $I\subseteq\mathbb{R}$ zawierającego $0$ definiujemy
$U_a(I)=\{0_a\}\cup\{[x]:x\in I,\ x\neq0\}$ oraz
$U_b(I)=\{0_b\}\cup\{[x]:x\in I,\ x\neq0\}$.
Oba zbiory zawierają te same punkty odpowiadające niezerowym liczbom rzeczywistym z $I$, ale różnią się punktem nad zerem: $0_a\in U_a(I)$ i $0_b\notin U_a(I)$, natomiast $0_b\in U_b(I)$ i $0_a\notin U_b(I)$.

Sprawdźmy, że zbiory te są otwarte w topologii ilorazowej. Z definicji zbiór $U\subseteq M$ jest otwarty wtedy i tylko wtedy, gdy $\pi^{-1}(U)$ jest otwarty w $X$. Dla $U_a(I)$ mamy
$\pi^{-1}(U_a(I))=(I\times\{0\})\cup((I\setminus\{0\})\times\{1\})$.
Pierwszy składnik jest otwarty w pierwszej kopii $\mathbb{R}$, a drugi w drugiej kopii, ponieważ $I\setminus\{0\}$ jest otwarty w $\mathbb{R}$. Zatem $\pi^{-1}(U_a(I))$ jest otwarty w sumie rozłącznej $X$, a więc $U_a(I)$ jest otwarty w $M$. Analogicznie
$\pi^{-1}(U_b(I))=((I\setminus\{0\})\times\{0\})\cup(I\times\{1\})$,
więc $U_b(I)$ również jest otwarty.

Zauważmy teraz, że $U_a=M\setminus\{0_b\}$ oraz $U_b=M\setminus\{0_a\}$ są otwarte i każdy z tych zbiorów jest homeomorficzny z $\mathbb{R}$. Rzeczywiście, definiujemy $\varphi_a\colon U_a\to\mathbb{R}$ przez $\varphi_a(0_a)=0$ oraz $\varphi_a([x])=x$ dla $x\neq0$. Ciągła funkcja $(x,i)\mapsto x$ jest stała na klasach równoważności,
więc z własności ilorazu indukuje ciągłą funkcję na $M$;
jej ograniczeniem jest $\varphi_a$. Odwrotność $\varphi_a^{-1}$ ma postać
$x\mapsto\pi(x,0)$ i również jest ciągła. Analogicznie, używając
$x\mapsto\pi(x,1)$, otrzymujemy homeomorfizm
$\varphi_b\colon U_b\to\mathbb{R}$.

A zatem każdy punkt przestrzeni $M$ posiada otwarte otoczenie homeomorficzne z otwartym podzbiorem $\mathbb{R}$. Przestrzeń $M$ jest więc lokalnie euklidesowa wymiaru $1$.

Ponadto drugi aksjomat przeliczalności jest spełniony. Możemy wziąć przeliczalną bazę przedziałów o końcach wymiernych na $\mathbb{R}$, a następnie przenieść ją przez homeomorfizmy $\varphi_a$ i $\varphi_b$ na $U_a$ i $U_b$. Suma dwóch przeliczalnych rodzin jest nadal przeliczalna, więc otrzymujemy przeliczalną bazę topologii na całym $M$.

Pozostaje sprawdzić własność Hausdorffa. Rozważmy dwa różne punkty $0_a\neq0_b$. Niech $U$ będzie dowolnym otwartym otoczeniem $0_a$, a $V$ dowolnym otwartym otoczeniem $0_b$. Ponieważ $U$ jest otoczeniem $0_a$, istnieje $\varepsilon>0$ takie, że $U_a((-\varepsilon,\varepsilon))\subseteq U$. Podobnie istnieje $\delta>0$ takie, że $U_b((-\delta,\delta))\subseteq V$.

Wybierzmy $x\neq0$ takie, że $|x|<\min\{\varepsilon,\delta\}$. Wtedy $[x]\in U_a((-\varepsilon,\varepsilon))$ oraz $[x]\in U_b((-\delta,\delta))$, a więc $[x]\in U\cap V$. Zatem $U\cap V\neq\varnothing$.

Nie istnieją więc rozłączne otoczenia punktów $0_a$ i $0_b$. Przestrzeń $M$ nie jest Hausdorffa.

Podsumowując, $M$ jest lokalnie euklidesowa wymiaru $1$ i spełnia drugi aksjomat przeliczalności, ale nie jest przestrzenią Hausdorffa. Pokazuje to, że warunek Hausdorffa w definicji rozmaitości topologicznej nie wynika z pozostałych założeń i musi być wymagany osobno.
\end{przyklad}

\begin{uwaga}
Symetrycznie, istnieją przestrzenie Hausdorffa i lokalnie euklidesowe, które nie spełniają
drugiego aksjomatu przeliczalności. Prostym przykładem jest suma
rozłączna nieprzeliczalnie wielu kopii $\R$: każda kopia jest otwarta,
co daje lokalną euklidesowość, a różne kopie rozdzielają punkty.
Gdyby istniała przeliczalna baza, w każdej kopii można byłoby wybrać
niepusty element bazy zawarty w tej kopii. Różne kopie wymagałyby
różnych elementów, co jest niemożliwe.
Klasycznym spójnym przykładem jest tzw.\ \emph{długa prosta}
(ang.\ \emph{long line}), skonstruowana przy użyciu pierwszej nieprzeliczalnej liczby
porządkowej. Nie przytaczamy tu tej konstrukcji; wspominamy o niej jedynie, aby podkreślić,
że oba warunki (i) oraz (ii) są od siebie niezależne i żadnego z nich nie można pominąć.
\end{uwaga}

\begin{fakt}
Każda rozmaitość topologiczna jest przestrzenią parazwartą i metryzowalną.
\end{fakt}

\begin{proof}
Każdy punkt ma otoczenie, którego domknięcie jest zwarte: w mapie
wybieramy kulę domkniętą zawartą w jej obrazie i przenosimy ją do $M$.
Jej obraz jest zwarty, a więc domknięty w przestrzeni Hausdorffa $M$.
Korzystamy teraz z trzech standardowych twierdzeń topologii ogólnej:
lokalnie zwarta przestrzeń Hausdorffa jest regularna;
regularna przestrzeń Hausdorffa z przeliczalną bazą jest metryzowalna
(twierdzenie Urysohna); każda przestrzeń metryczna jest parazwarta.
To daje obie tezy.
\end{proof}
Parazwartość oznacza, że każde pokrycie otwarte ma lokalnie skończone
otwarte pokrycie wpisane. Na rozmaitości gładkiej pozwala ona,
w połączeniu z lokalnymi funkcjami odcinającymi, konstruować gładkie
rozkłady jedności. Sama lokalna funkcja odcinająca z
Sekcji~\ref{sec:przestrzen-styczna} nie wymaga jeszcze tego globalnego narzędzia.

\section{Atlasy i rozmaitości gładkie}\label{sec:atlas}

Struktura rozmaitości topologicznej pozwala mówić o ciągłości, ale nie wystarcza, by mówić
o różniczkowalności — pojęcie to nie jest niezmiennicze względem dowolnych homeomorfizmów.
Aby zdefiniować rozmaitość \emph{gładką}, musimy wyposażyć $M$ w dodatkową strukturę: atlas
map, których wzajemne przejścia są dyfeomorfizmami.

\begin{definicja}
Niech $M$ będzie rozmaitością topologiczną wymiaru $n$. \emph{Mapą} na $M$
nazywamy parę $(U,\varphi)$, gdzie $U\subseteq M$ jest zbiorem otwartym, a
$\varphi\colon U\to\varphi(U)\subseteq\R^n$ jest homeomorfizmem na otwarty podzbiór $\R^n$.
\end{definicja}

\noindent\textbf{Ustalenie terminologii.}
W całym skrypcie \emph{mapa} oznacza angielskie \emph{chart},
czyli lokalny układ współrzędnych $(U,\varphi)$.
Angielskie \emph{map} tłumaczymy jako \emph{odwzorowanie}.
Sama funkcja współrzędnych $\varphi$ jest więc odwzorowaniem
wchodzącym w skład mapy.

\begin{uwaga}
Mapa jest lokalnym opisem rozmaitości: jej dziedziną jest otwarty zbiór $U\subset M$,
a obrazem otwarty zbiór $\varphi(U)\subset\mathbb R^n$. Dzięki temu lokalne rachunki na $M$
możemy wykonywać w zwykłych współrzędnych euklidesowych.
\end{uwaga}

\begin{figure}[H]
\centering
\begin{tikzpicture}[>=Stealth,scale=1.04,line cap=round,line join=round]
% Fragment M oraz wyróżniony otwarty zbiór U.
\filldraw[fill=manifoldblue!10,draw=manifoldblue!75!black,very thick]
 (-5.05,-.55) .. controls (-5.15,.55) and (-4.15,1.34) .. (-2.65,1.2)
 .. controls (-1.43,1.14) and (-.67,.48) .. (-.67,-.35)
 .. controls (-.7,-1.12) and (-2.18,-1.43) .. (-3.5,-1.18)
 .. controls (-4.5,-1.1) and (-5.0,-.94) .. cycle;
\filldraw[fill=manifoldred!12,draw=manifoldred,very thick]
 (-4.08,-.28) .. controls (-4.1,.34) and (-3.4,.73) .. (-2.72,.64)
 .. controls (-2.03,.63) and (-1.56,.29) .. (-1.56,-.25)
 .. controls (-1.57,-.79) and (-2.23,-.96) .. (-2.97,-.85)
 .. controls (-3.65,-.82) and (-4.07,-.69) .. cycle;
\node[manifoldblue!75!black,font=\small\bfseries] at (-4.71,1.35) {$M$};
\node[manifoldred!85!black,font=\small] at (-3.58,.17) {$U$};
\fill[manifoldred] (-2.87,-.1) circle (2.5pt);
\node[font=\small] at (-2.88,-.44) {$p$};
% Obraz mapy w przestrzeni współrzędnych, bez siatki.
\filldraw[fill=manifoldgreen!6,draw=manifoldgreen!70!black,thick]
 (1.45,-1.19)--(4.99,-1.19)--(5.48,1.0)--(1.94,1.0)--cycle;
\filldraw[fill=manifoldred!12,draw=manifoldred,very thick]
 (2.1,-.31) .. controls (2.14,.23) and (2.78,.6) .. (3.45,.59)
 .. controls (4.16,.58) and (4.57,.26) .. (4.5,-.27)
 .. controls (4.43,-.75) and (3.76,-.86) .. (3.04,-.79)
 .. controls (2.43,-.73) and (2.08,-.62) .. cycle;
\node[manifoldgreen!65!black,font=\small\bfseries] at (5.2,1.32) {$\R^n$};
\node[manifoldred!85!black,font=\small] at (3.5,.29) {$\varphi(U)$};
\fill[manifoldred] (3.2,-.07) circle (2.5pt);
\node[font=\small] at (3.23,-.44) {$\varphi(p)$};
\draw[->,very thick,manifoldred] (-.39,.02)--(1.12,.02)
 node[midway,above,font=\small] {$\varphi$};
\end{tikzpicture}
\caption{Mapa $(U,\varphi)$ przenosi otwarty fragment $U\subset M$ na obszar
$\varphi(U)\subset\R^n$; punktowi $p$ przyporządkowuje współrzędne $\varphi(p)$.}
\label{fig:mapa-rozmaitosc}
\end{figure}


\begin{definicja}
Dwie mapy $(U,\varphi)$ i $(V,\psi)$ na $M$ są \emph{zgodne} ($C^\infty$-zgodne), jeżeli
albo $U\cap V=\emptyset$, albo odwzorowania przejścia
\[
\psi\circ\varphi^{-1}\colon\varphi(U\cap V)\to\psi(U\cap V), \qquad
\varphi\circ\psi^{-1}\colon\psi(U\cap V)\to\varphi(U\cap V)
\]
są klasy $C^\infty$ (są to wtedy automatycznie odwzorowania wzajemnie odwrotne między
otwartymi podzbiorami $\R^n$, a więc dyfeomorfizmy).
\end{definicja}



Na Rysunku~\ref{fig:atlas-przejscie-dwie-mapy} pokazano
dwie zgodne mapy i współrzędne punktu
leżącego w ich części wspólnej. Odwzorowanie przejścia
działa \emph{między obrazami tej części wspólnej}, a nie
między całymi $\varphi_\alpha(U_\alpha)$ i
$\varphi_\beta(U_\beta)$, jeśli dziedziny map się różnią.

\begin{figure}[H]
\centering
\begin{tikzpicture}[>=Latex,line cap=round,line join=round,
                     font=\small]
  % Fragment rozmaitości i dwa nakładające się obszary map.
  \filldraw[fill=manifoldblue!5,draw=manifoldblue!75!black,thick]
   (-3.35,1.25) .. controls (-3.32,2.76) and (-2.20,3.55)
   .. (-.25,3.38) .. controls (1.72,3.57) and (3.36,2.88)
   .. (3.40,1.59) .. controls (3.48,.77) and (1.68,.63)
   .. (.18,.85) .. controls (-1.90,.53) and (-3.39,.67)
   .. cycle;
  \fill[manifoldblue!22,opacity=.85]
       (-.88,2.12) ellipse (1.71 and .94);
  \fill[manifoldgold!26,opacity=.86]
       (.88,2.12) ellipse (1.71 and .94);
  \draw[manifoldblue!78!black,thick]
       (-.88,2.12) ellipse (1.71 and .94);
  \draw[manifoldgold!82!black,thick]
       (.88,2.12) ellipse (1.71 and .94);
  \node[manifoldblue!75!black,font=\bfseries] at (-3.05,3.13) {$M$};
  \node[manifoldblue!80!black] at (-1.71,2.32) {$U_\alpha$};
  \node[manifoldgold!70!black] at (1.72,2.32) {$U_\beta$};
  \node[font=\scriptsize] at (0,2.44) {$U_\alpha\cap U_\beta$};
  \fill[manifoldred] (0,1.95) circle (2.1pt);
  \node[below,manifoldred!85!black] at (0,1.88) {$p$};
  % Obrazy w dwóch osobnych przestrzeniach współrzędnych.
  \filldraw[fill=manifoldblue!5,draw=manifoldblue!55!black,thick]
       (-4.82,-1.57)--(-1.79,-1.57)--(-1.42,-.03)
       --(-4.45,-.03)--cycle;
  \fill[manifoldblue!15] (-3.08,-.77) ellipse (1.05 and .52);
  \draw[manifoldblue!78!black,thick]
       (-3.08,-.77) ellipse (1.05 and .52);
  \draw[manifoldred!70!black,densely dashed]
       (-2.80,-.83) ellipse (.40 and .28);
  \fill[manifoldred] (-2.80,-.83) circle (1.9pt);
  \node[manifoldred!85!black] at (-2.80,-1.32)
       {$\varphi_\alpha(p)$};
  \node[manifoldblue!70!black] at (-4.54,.13) {$\R^n$};

  \filldraw[fill=manifoldgold!5,draw=manifoldgold!68!black,thick]
       (1.42,-1.57)--(4.45,-1.57)--(4.82,-.03)
       --(1.79,-.03)--cycle;
  \fill[manifoldgold!17] (3.08,-.77) ellipse (1.05 and .52);
  \draw[manifoldgold!82!black,thick]
       (3.08,-.77) ellipse (1.05 and .52);
  \draw[manifoldred!70!black,densely dashed]
       (2.80,-.83) ellipse (.40 and .28);
  \fill[manifoldred] (2.80,-.83) circle (1.9pt);
  \node[manifoldred!85!black] at (2.80,-1.32)
       {$\varphi_\beta(p)$};
  \node[manifoldgold!72!black] at (4.50,.13) {$\R^n$};

  \draw[->,manifoldblue!80!black,very thick]
       (-1.91,1.05)--(-3.37,.16)
       node[midway,left] {$\varphi_\alpha$};
  \draw[->,manifoldgold!78!black,very thick]
       (1.91,1.05)--(3.37,.16)
       node[midway,right] {$\varphi_\beta$};
  \node[manifoldblue!80!black] at (-3.08,-1.90)
       {$\varphi_\alpha(U_\alpha)$};
  \node[manifoldgold!78!black] at (3.08,-1.90)
       {$\varphi_\beta(U_\beta)$};
  \draw[->,manifoldgreen!75!black,very thick]
       (-2.11,-2.41)--(2.11,-2.41)
       node[midway,below] {$\varphi_\beta\circ\varphi_\alpha^{-1}$};
\end{tikzpicture}
\caption{Dwie nakładające się mapy. Przerywane obszary
oznaczają obrazy części wspólnej. Punkt $p$ ma
dwa zapisy współrzędnych, a przejście
$\varphi_\beta\circ\varphi_\alpha^{-1}$ prowadzi
z $\varphi_\alpha(U_\alpha\cap U_\beta)$ do
$\varphi_\beta(U_\alpha\cap U_\beta)$. W atlasie
pozostałe mapy pokrywają ewentualną resztę $M$.}
\label{fig:atlas-przejscie-dwie-mapy}
\end{figure}

\begin{definicja}
\emph{Atlasem} (gładkim, $C^\infty$) na rozmaitości topologicznej $M$ nazywamy rodzinę map
$\mathcal{A}=\{(U_\alpha,\varphi_\alpha)\}_{\alpha\in I}$ taką, że:
\begin{enumerate}[label=(\roman*)]
\item $\bigcup_{\alpha\in I}U_\alpha=M$,
\item mapy $(U_\alpha,\varphi_\alpha)$ oraz $(U_\beta,\varphi_\beta)$ są zgodne dla
wszystkich $\alpha,\beta\in I$.
\end{enumerate}
\end{definicja}



Chcemy teraz porównywać różne atlasy na tej samej przestrzeni topologicznej $M$ — dwa
atlasy powinny wyznaczać tę samą strukturę gładką dokładnie wtedy, gdy razem tworzą
(większy) atlas.

\begin{definicja}
Dwa atlasy $\mathcal{A}_1$, $\mathcal{A}_2$ na $M$ są \emph{równoważne}, $\mathcal A_1\sim
\mathcal A_2$, jeżeli ich suma $\mathcal{A}_1\cup\mathcal{A}_2$ jest atlasem, tzn.\ każda
mapa z $\mathcal{A}_1$ jest zgodna z każdą mapą z $\mathcal{A}_2$.
\end{definicja}

Poniższy lemat techniczny jest kluczowym narzędziem rachunkowym tego paragrafu; wykorzystamy
go dwukrotnie: do dowodu, że $\sim$ jest relacją równoważności, oraz do dowodu istnienia i
jednoznaczności atlasu maksymalnego.

\begin{lemat}\label{lem:zgodnosc-przechodnia}
Niech $(U,\varphi)$, $(W,\eta)$, $(V,\psi)$ będą mapami na $M$ takimi, że $(U,\varphi)$
jest zgodna z $(W,\eta)$, a $(W,\eta)$ jest zgodna z $(V,\psi)$. Wówczas dla dowolnego
punktu $p\in U\cap V\cap W$, odwzorowania $\psi\circ\varphi^{-1}$ i $\varphi\circ\psi^{-1}$
są klasy $C^\infty$ w (odpowiednich) otoczeniach punktów $\varphi(p)$ i $\psi(p)$.
\end{lemat}

\begin{proof}
Na otoczeniu $\varphi(U\cap V\cap W)$ punktu $\varphi(p)$ mamy równość
$\psi\circ\varphi^{-1}=(\psi\circ\eta^{-1})\circ(\eta\circ\varphi^{-1})$; odwzorowanie
$\eta\circ\varphi^{-1}$ jest gładkie na $\varphi(U\cap W)$ (zgodność $(U,\varphi)$ z
$(W,\eta)$), a $\psi\circ\eta^{-1}$ jest gładkie na $\eta(W\cap V)$ (zgodność $(W,\eta)$ z
$(V,\psi)$); złożenie odwzorowań gładkich jest gładkie, więc $\psi\circ\varphi^{-1}$ jest
gładkie w otoczeniu $\varphi(p)$. Analogicznie, na otoczeniu $\psi(p)$:
$\varphi\circ\psi^{-1}=(\varphi\circ\eta^{-1})\circ(\eta\circ\psi^{-1})$ jest złożeniem
odwzorowań gładkich ($\varphi\circ\eta^{-1}$ i $\eta\circ\psi^{-1}$ gładkie, z tych samych
dwóch założeń zgodności), a więc gładkie.
\end{proof}

\begin{stwierdzenie}
Równoważność atlasów jest relacją równoważności na zbiorze wszystkich atlasów na $M$.
\end{stwierdzenie}

\begin{proof}
Zwrotność i symetria są oczywiste z definicji. Sprawdzimy przechodniość. Niech
$\mathcal{A}_1\sim\mathcal{A}_2$ i $\mathcal{A}_2\sim\mathcal{A}_3$; pokażemy, że
$\mathcal{A}_1\sim\mathcal{A}_3$, czyli że dowolna mapa $(U,\varphi)\in\mathcal{A}_1$ jest
zgodna z dowolną mapą $(V,\psi)\in\mathcal{A}_3$. Jeśli $U\cap V=\emptyset$, zgodność jest
trywialna. W przeciwnym razie ustalmy $p\in U\cap V$. Ponieważ $\mathcal{A}_2$ jest atlasem
pokrywającym $M$, istnieje mapa $(W,\eta)\in\mathcal{A}_2$ z $p\in W$. Z założenia
$\mathcal{A}_1\sim\mathcal{A}_2$, mapa $(U,\varphi)$ jest zgodna z $(W,\eta)$; z założenia
$\mathcal{A}_2\sim\mathcal{A}_3$, mapa $(W,\eta)$ jest zgodna z $(V,\psi)$. Na mocy
Lematu~\ref{lem:zgodnosc-przechodnia} (zastosowanego w punkcie $p\in U\cap V\cap W$),
odwzorowania $\psi\circ\varphi^{-1}$ i $\varphi\circ\psi^{-1}$ są gładkie w otoczeniach
odpowiednio $\varphi(p)$ i $\psi(p)$. Ponieważ punkt $p\in U\cap V$ był dowolny, oba te
odwzorowania są gładkie na całych swoich dziedzinach $\varphi(U\cap V)$, $\psi(U\cap V)$
(gładkość jest własnością lokalną). Zatem $(U,\varphi)$ i $(V,\psi)$ są zgodne, co dowodzi
$\mathcal{A}_1\sim\mathcal{A}_3$.
\end{proof}

W praktyce nie chcemy jednak utożsamiać struktury gładkiej z pojedynczym atlasem — dwa
równoważne atlasy powinny wyznaczać \emph{tę samą} strukturę. Wygodnie jest wybrać
kanonicznego reprezentanta każdej klasy równoważności: atlas maksymalny.

\begin{definicja}
Atlas $\mathcal{A}$ na $M$ jest \emph{maksymalny}, jeśli nie jest właściwym podzbiorem
żadnego innego atlasu na $M$, tzn.\ jeśli mapa $(U,\varphi)$ jest zgodna z każdą mapą
należącą do $\mathcal{A}$, to $(U,\varphi)\in\mathcal{A}$.
\end{definicja}

\begin{twierdzenie}\label{tw:atlas-maksymalny}
Każdy atlas $\mathcal{A}$ na $M$ jest zawarty w dokładnie jednym atlasie maksymalnym
$\mathcal{A}^+$, zadanym wzorem
\[
\mathcal{A}^+=\{(U,\varphi)\text{ -- mapa na }M : (U,\varphi)\text{ jest zgodna z każdą
mapą z }\mathcal{A}\}.
\]
Ponadto dwa atlasy $\mathcal{A}_1$, $\mathcal{A}_2$ są równoważne wtedy i tylko wtedy, gdy
$\mathcal{A}_1^+=\mathcal{A}_2^+$.
\end{twierdzenie}

\begin{proof}
\emph{Krok 1: $\mathcal{A}^+$ jest atlasem.} Ponieważ $\mathcal{A}\subseteq\mathcal{A}^+$
(każda mapa z $\mathcal{A}$ jest zgodna z każdą mapą z $\mathcal{A}$, gdyż $\mathcal{A}$
jest atlasem), zbiory dziedzin map z $\mathcal{A}^+$ pokrywają $M$. Pozostaje sprawdzić, że
dowolne dwie mapy $(U,\varphi),(V,\psi)\in\mathcal{A}^+$ są ze sobą zgodne. Ustalmy
$p\in U\cap V$ (dla $U\cap V=\emptyset$ zgodność jest trywialna); ponieważ $\mathcal{A}$
pokrywa $M$, istnieje $(W,\eta)\in\mathcal{A}$ z $p\in W$. Mapa $(U,\varphi)$ jest zgodna z
$(W,\eta)$ (bo $(U,\varphi)\in\mathcal{A}^+$ jest zgodna z każdą mapą z $\mathcal A$), a
$(W,\eta)$ jest zgodna z $(V,\psi)$ (analogicznie, bo $(V,\psi)\in\mathcal A^+$). Na mocy
Lematu~\ref{lem:zgodnosc-przechodnia}, odwzorowania $\psi\circ\varphi^{-1}$,
$\varphi\circ\psi^{-1}$ są gładkie w otoczeniach $\varphi(p)$, $\psi(p)$; wobec dowolności
$p\in U\cap V$, są gładkie na całych dziedzinach $\varphi(U\cap V)$, $\psi(U\cap V)$. Zatem
$\mathcal{A}^+$ jest atlasem.

\emph{Krok 2: $\mathcal{A}^+$ jest maksymalny.} Niech mapa $(U,\varphi)$ będzie zgodna z
każdą mapą należącą do $\mathcal{A}^+$. Ponieważ $\mathcal{A}\subseteq\mathcal{A}^+$, mapa
$(U,\varphi)$ jest w szczególności zgodna z każdą mapą z $\mathcal{A}$, a więc z definicji
$\mathcal{A}^+$ mamy $(U,\varphi)\in\mathcal{A}^+$.

\emph{Krok 3: jednoznaczność.} Niech $\mathcal B\supseteq\mathcal A$ będzie atlasem
maksymalnym. Ponieważ $\mathcal B$ jest atlasem i $\mathcal A\subseteq\mathcal B$, każda
mapa z $\mathcal B$ jest zgodna z każdą mapą z $\mathcal A$ (obie należą do atlasu
$\mathcal B$) — więc $\mathcal B\subseteq\mathcal A^+$ z definicji $\mathcal A^+$.

Odwrotnie, niech $(U,\varphi)\in\mathcal A^+$; pokażemy $(U,\varphi)\in\mathcal B$, co —
wobec dowolności $(U,\varphi)\in\mathcal A^+$ — da $\mathcal A^+\subseteq\mathcal B$, a
wraz z poprzednią inkluzją: $\mathcal B=\mathcal A^+$. Z maksymalności $\mathcal B$
wystarczy pokazać, że $(U,\varphi)$ jest zgodna z każdą mapą $(V,\psi)\in\mathcal B$.
Ustalmy taką mapę i punkt $p\in U\cap V$ (dla $U\cap V=\emptyset$ zgodność jest trywialna).
Ponieważ $\mathcal A$ pokrywa $M$, istnieje $(W,\eta)\in\mathcal A$ z $p\in W$. Mapa
$(U,\varphi)$ jest zgodna z $(W,\eta)$ (bo $(U,\varphi)\in\mathcal A^+$), a $(W,\eta)\in
\mathcal A\subseteq\mathcal B$ jest zgodna z $(V,\psi)\in\mathcal B$ (bo obie należą do
atlasu $\mathcal B$). Na mocy Lematu~\ref{lem:zgodnosc-przechodnia}, $(U,\varphi)$ i
$(V,\psi)$ są zgodne w otoczeniu $p$; wobec dowolności $p\in U\cap V$ — zgodne na całych
dziedzinach. Zatem $(U,\varphi)$ jest zgodna z każdą mapą z $\mathcal B$, więc — z
maksymalności $\mathcal B$ — $(U,\varphi)\in\mathcal B$.

\emph{Krok 4: $\mathcal A_1\sim\mathcal A_2\iff\mathcal A_1^+=\mathcal A_2^+$.}
($\Leftarrow$) Jeśli $\mathcal A_1^+=\mathcal A_2^+=:\mathcal B$, to $\mathcal A_1\subseteq
\mathcal B$ i $\mathcal A_2\subseteq\mathcal B$, a skoro $\mathcal B$ jest atlasem, każda
mapa z $\mathcal A_1$ jest zgodna z każdą mapą z $\mathcal B$, w szczególności z każdą mapą
z $\mathcal A_2$ — czyli $\mathcal A_1\sim\mathcal A_2$.

($\Rightarrow$) Załóżmy $\mathcal A_1\sim\mathcal A_2$. Pokażemy $\mathcal A_1^+\subseteq
\mathcal A_2^+$ (inkluzja odwrotna analogicznie, skąd równość). Niech $(U,\varphi)\in
\mathcal A_1^+$. Ustalmy dowolną mapę $(V,\psi)\in\mathcal A_2$ i (dla $U\cap V\neq
\emptyset$, inaczej trywialnie) punkt $p\in U\cap V$. Ponieważ $\mathcal A_1$ pokrywa $M$,
istnieje $(W,\eta)\in\mathcal A_1$ z $p\in W$. Mapa $(U,\varphi)$ jest zgodna z $(W,\eta)\in
\mathcal A_1$ (bo $(U,\varphi)\in\mathcal A_1^+$), a $(W,\eta)\in\mathcal A_1$ jest zgodna z
$(V,\psi)\in\mathcal A_2$ (bo $\mathcal A_1\sim\mathcal A_2$). Na mocy
Lematu~\ref{lem:zgodnosc-przechodnia}, $(U,\varphi)$ i $(V,\psi)$ są zgodne w otoczeniu
$p$; wobec dowolności $p\in U\cap V$ — na całych dziedzinach. Zatem $(U,\varphi)$ jest
zgodna z $(V,\psi)$; wobec dowolności $(V,\psi)\in\mathcal A_2$, mamy $(U,\varphi)\in
\mathcal A_2^+$.
\end{proof}

Powyższe twierdzenie pozwala oddzielić topologię przestrzeni od
wyboru współrzędnych, w których będziemy różniczkować.

\begin{definicja}[Struktura różniczkowa]\label{def:struktura-rozniczkowa}
\index{struktura rozniczkowa@Struktura różniczkowa}
\emph{Strukturą różniczkową} klasy $C^\infty$ na rozmaitości
topologicznej $M$ nazywamy klasę równoważności atlasów gładkich
na $M$. Równoważnie, na mocy Twierdzenia~\ref{tw:atlas-maksymalny},
jest nią jedyny atlas maksymalny $\mathcal A=\mathcal A_0^+$
zawierający dowolny atlas $\mathcal A_0$ z tej klasy. Mapy
należące do $\mathcal A$ nazywamy \emph{dopuszczalnymi}.
\end{definicja}

\begin{definicja}[Rozmaitość gładka]\label{def:rozmaitosc-gladka}
\index{rozmaitosc gladka@Rozmaitość gładka}
\emph{Rozmaitością gładką} (klasy $C^\infty$) wymiaru $n$
nazywamy parę $(M,\mathcal A)$, gdzie $M$ jest rozmaitością
topologiczną wymiaru $n$, a $\mathcal A$ jest strukturą
różniczkową zapisaną jako atlas maksymalny.
\end{definicja}

W praktyce wystarczy podać mały atlas $\mathcal A_0$:
wyznacza on jednoznacznie strukturę $\mathcal A_0^+$.
Topologia określa, które mapy współrzędnych są homeomorfizmami;
struktura różniczkowa określa dodatkowo, które przejścia
między nimi są gładkie.

\begin{przyklad}[Różne struktury na tej samej przestrzeni]
Na topologicznej prostej $\R$ atlas $\mathcal A_{\mathrm{std}}$
z mapą $\id\colon\R\to\R$ i atlas $\mathcal A_h$ z mapą
$h\colon\R\to\R$, $h(x)=x^3$, są osobno atlasami gładkimi.
Nie są równoważne: w jednym kierunku
przejście jest dane przez $h^{-1}(u)=\sqrt[3]{u}$,
które nie jest różniczkowalne w zerze.
Definiują zatem różne struktury na \emph{ustalonym} zbiorze
topologicznym $\R$, choć otrzymane rozmaitości są
dyfeomorficzne (w sensie definicji~\ref{def dyfeo}): dyfeomorfizm
$h\colon(\R,\mathcal A_h^+)\to(\R,\mathcal A_{\mathrm{std}}^+)$
ma w odpowiednich mapach współrzędnych postać identyczności.
\end{przyklad}

\begin{uwaga}[Intuicja: mapa, atlas i rozmaitość]
Najprostszym sposobem zrozumienia pojęć \emph{mapy}, \emph{atlasu} i
\emph{rozmaitości} jest analogia z powierzchnią Ziemi.

Powierzchnia Ziemi jest obiektem zakrzywionym i nie możemy przedstawić jej
w całości na jednej płaskiej kartce bez zniekształceń. Możemy jednak wybrać
niewielki fragment powierzchni, na przykład okolice Krakowa, i przedstawić go
na płaskiej mapie. Matematycznie dokładnie taką rolę pełni \emph{mapa}
rozmaitości.

Jeżeli $M$ jest rozmaitością, to mapa jest parą $(U,\varphi)$, gdzie
$U\subset M$ jest otwartym fragmentem rozmaitości, a
\[
    \varphi\colon U\longrightarrow \varphi(U)\subset\mathbb{R}^n
\]
jest homeomorfizmem. Oznacza to, że punkty znajdujące się w $U$ możemy opisać
za pomocą zwykłych współrzędnych
\[
    (x^1,\ldots,x^n)\in\mathbb{R}^n.
\]
Można więc myśleć o $\varphi$ jako o sposobie ``rozłożenia'' niewielkiego
fragmentu rozmaitości na płaskiej przestrzeni $\mathbb{R}^n$.

Jedna mapa zazwyczaj nie wystarcza do opisania całej przestrzeni. Podobnie
jedna mapa geograficzna przedstawiająca okolice Krakowa nie opisuje całej
powierzchni Ziemi. Potrzebujemy wielu map
\[
    (U_\alpha,\varphi_\alpha),
\]
których dziedziny pokrywają całą przestrzeń:
\[
    M=\bigcup_{\alpha}U_\alpha.
\]
Taki zbiór map nazywamy \emph{atlasem}.

Jeżeli dwa obszary $U_\alpha$ i $U_\beta$ nachodzą na siebie, to ten sam punkt
rozmaitości może mieć dwa różne opisy współrzędnościowe. Przejście pomiędzy
nimi opisuje odwzorowanie
\[
    \varphi_\beta\circ\varphi_\alpha^{-1},
\]
zwane odwzorowaniem przejścia. Można je porównać do przeliczania współrzędnych
tego samego miejsca pomiędzy dwiema zachodzącymi na siebie mapami
geograficznymi.

\emph{Rozmaitość} jest zatem przestrzenią, która globalnie może być
skomplikowana lub zakrzywiona, ale w dostatecznie małym otoczeniu każdego
punktu wygląda jak zwykła przestrzeń euklidesowa $\mathbb{R}^n$.

W skrócie:
\[
\boxed{
\begin{array}{ccl}
\text{mapa}      &\longleftrightarrow& \text{lokalny układ współrzędnych},\\[2mm]
\text{atlas}     &\longleftrightarrow& \text{zbiór map opisujących całą przestrzeń},\\[2mm]
\text{rozmaitość}&\longleftrightarrow& \text{przestrzeń lokalnie wyglądająca jak }\mathbb{R}^n.
\end{array}}
\]

Przykładowo sfera $S^2$ nie jest płaską częścią $\mathbb{R}^2$, ale każdy jej
dostatecznie mały fragment można opisać za pomocą dwóch współrzędnych.
Dlatego $S^2$ jest dwuwymiarową rozmaitością.
\end{uwaga}

Od tej pory, mówiąc o rozmaitości $M$, będziemy mieli na myśli rozmaitość gładką — parę
$(M,\mathcal{A})$ — pomijając $\mathcal{A}$ w oznaczeniach, gdy nie prowadzi to do
nieporozumień.

\section{Rozmaitości z brzegiem}
\label{sec:rozmaitosci-z-brzegiem-1}
\index{brzeg rozmaitosci@Brzeg rozmaitości}

Wnętrze dysku wygląda lokalnie jak otwarty zbiór w $\R^n$,
ale punkt na jego brzegu ma tylko jedną stronę. Modelem takiego
otoczenia jest półprzestrzeń
\[
 \mathbb H^n=\{(x^1,\ldots,x^n)\in\R^n:x^n\geq0\},
 \qquad n\geq1.
\]
Otwartość podzbioru $\mathbb H^n$ rozumiemy w topologii względnej.

\begin{definicja}[Rozmaitość gładka z brzegiem]
\index{rozmaitosc gladka z brzegiem@Rozmaitość gładka z brzegiem}
\label{def:rozmaitosc-z-brzegiem-1}
\emph{Rozmaitość gładka z brzegiem} wymiaru $n\geq1$ to przestrzeń
Hausdorffa spełniająca drugi aksjomat przeliczalności, wyposażona
w atlas homeomorfizmów
$\varphi_\alpha:U_\alpha\to V_\alpha$ na zbiory
$V_\alpha$ otwarte w $\mathbb H^n$. Na częściach wspólnych
odwzorowania przejścia i ich odwrotności mają być gładkie:
w otoczeniu każdego punktu rozumie się przez to istnienie
gładkiego przedłużenia na otwarty zbiór w $\R^n$.
Maksymalny atlas zgodny z tym atlasem wyznacza strukturę gładką,
tak jak w Twierdzeniu~\ref{tw:atlas-maksymalny}.
\emph{Brzegiem rozmaitości} nazywamy zbiór
\[
 \partial M=\{p\in M:\varphi_\alpha(p)^n=0
             \text{ w pewnej mapie }(U_\alpha,\varphi_\alpha)\}.
\]
Pozostałe punkty tworzą wnętrze $M^\circ=M\setminus\partial M$.
Dla wymiaru $0$ przyjmujemy $\partial M=\varnothing$.
\end{definicja}

Warunek „w pewnej mapie” jest poprawny dzięki następującemu
faktowi; jednocześnie pokazuje, że brzeg nie zależy od wyboru
współrzędnych.

\begin{stwierdzenie}[Niezależność brzegu od mapy]
\label{prop:brzeg-niezalezny-1}
Dla $n\geq1$, jeśli $p$ ma ostatnią współrzędną równą $0$ w jednej
mapie półprzestrzennej, to ma ją równą $0$ w każdej
mapie półprzestrzennej, której dziedzina zawiera $p$.
Zbiór $\partial M$ jest zamknięty w $M$ i z map ograniczonych
do $\partial M$ otrzymuje strukturę gładkiej
rozmaitości wymiaru $n-1$ bez brzegu.
\end{stwierdzenie}
\begin{proof}
Przypuśćmy, że $\varphi(p)$ leży na płaszczyźnie
$x^n=0$, a $\psi(p)$ jest punktem wnętrza $\mathbb H^n$.
Ograniczmy mapy do wspólnego dostatecznie małego otoczenia
$p$. Wtedy $\psi$ odwzorowuje je na zbiór otwarty w $\R^n$,
a $\varphi\circ\psi^{-1}$ jest ciągłą injekcją tego zbioru
do $\R^n$. Twierdzenie Brouwera o niezmienniczości obszaru
orzeka, że obraz takiej injekcji jest otwarty w $\R^n$.
Obraz leży jednak w $\mathbb H^n$ i zawiera punkt
$\varphi(p)$ na jego płaszczyźnie brzegowej, więc nie może
być otwarty w $\R^n$. Sprzeczność dowodzi pierwszej tezy.
Dokładne sformułowanie użytego wyniku zewnętrznego podaje
\href{https://pi.math.cornell.edu/~hatcher/AT/AT.pdf}
{A. Hatcher, \emph{Algebraic Topology}, twierdzenie 2B.3}.

W każdej mapie punkty z $x^n>0$ tworzą zbiór otwarty.
Zatem $M^\circ$ jest otwarty, a $\partial M$ zamknięty.
Ograniczenie mapy $\varphi$ do
$U\cap\partial M$ przyjmuje wartości w otwartym
podzbiorze $\R^{n-1}\times\{0\}$. Przejścia między
tymi ograniczeniami są gładkie, bo są ograniczeniami
gładkich przedłużeń przejść na $\R^n$; to samo dotyczy
odwrotności. Są to więc mapy gładkiej rozmaitości
wymiaru $n-1$. Każda z nich ma obraz otwarty w całej
$\R^{n-1}$, czyli $\partial M$ nie ma własnego brzegu.
\end{proof}

W punkcie brzegu definiujemy $T_pM$ przez
wektory współrzędnych w półprzestrzennej mapie,
równoważnie przez derywacje gładkich kiełków
funkcji w $p$. Pełną konstrukcję i niezależność od przedłużeń
podamy po wprowadzeniu przestrzeni stycznej, w
Stwierdzeniu~\ref{prop:styczna-brzeg-1}. Ten model ma $n$ niezależnych
kierunków: $T_pM\cong\R^n$, natomiast
$T_p\partial M\cong\{v\in\R^n:v^n=0\}$ ma wymiar $n-1$.
Wektor o ujemnej ostatniej współrzędnej należy do $T_pM$,
choć nie jest prędkością krzywej biegnącej od $p$ do wnętrza
przy dodatnim czasie. Przykładowo $T_0[0,1]\cong\R$,
podczas gdy $T_0\partial[0,1]=0$.
Model przez krzywe określone na otwartym przedziale,
omówiony w Sekcji~\ref{sec:przestrzen-styczna},
dotyczy punktów wnętrza: w punkcie brzegu
dwustronna krzywa w $M$ nie reprezentuje
kierunku poprzecznego do brzegu.

\begin{przyklad}[Kula, okrąg i brzeg otaczającej przestrzeni]
\label{prz:brzeg-kula-okrag-1}
Domknięta kula $D^n=\{x\in\R^n:|x|\leq1\}$ ma brzeg
rozmaitościowy $S^{n-1}$. Rzeczywiście, przy $|p|=1$
funkcja $\rho(x)=1-|x|^2$ ma niezerową różniczkę
$\dd\rho_p(v)=-2\langle p,v\rangle$. Twierdzenie o funkcji
odwrotnej pozwala użyć $\rho$ jako ostatniej lokalnej
współrzędnej, zatem $D^n$ odpowiada warunkowi $\rho\geq0$.
Z kolei brzeg okręgu $S^1\subset\R^2$ jako rozmaitości
(w sensie Definicji~\ref{def:rozmaitosc-z-brzegiem-1})
jest pusty: $\partial S^1=\varnothing$. Jego brzeg topologiczny
jako \emph{podzbioru $\R^2$} jest natomiast całym $S^1$: okrąg ma
puste wnętrze w $\R^2$ i jest domknięty.
Te dwa znaczenia słowa „brzeg” trzeba odróżniać.
\end{przyklad}

Produkt dwóch rozmaitości z brzegiem może mieć
\emph{naroża}. W kwadracie $[0,1]^2$ przy wierzchołku
naturalny model pochodzący z zanurzenia to
$[0,\infty)^2$, z dwiema ścianami brzegowymi.
W rachunku uchwytów spotkamy taki model na
$\partial D^k\times\partial D^{n-k}$; naroże można
zaokrąglić w małym otoczeniu, zachowując części
odległe od niego. To operacja na gładkiej geometrii
brzegu, a nie zmiana tego, które punkty są brzegowe
topologicznie. Kołnierze i takie zaokrąglanie pojawią się
w Sekcji~\ref{sec:kolnierz-15} i przy uchwytach w
Sekcji~\ref{sec:uchwyt-morse-17}.

\section{Rozmaitości zespolone}

Analogicznie do konstrukcji z poprzedniego paragrafu, możemy modelować rozmaitość lokalnie
nie na $\R^n$, lecz na $\C^n$, żądając przy tym, by odwzorowania przejścia między mapami
były holomorficzne, a nie tylko gładkie.

\begin{definicja}
\emph{Rozmaitością zespoloną} wymiaru (zespolonego) $n$ nazywamy przestrzeń Hausdorffa $M$
spełniającą drugi aksjomat przeliczalności, wyposażoną w atlas $\{(U_\alpha,\varphi_\alpha)\}$
map $\varphi_\alpha\colon U_\alpha\to\varphi_\alpha(U_\alpha)\subseteq\C^n$ (homeomorfizmów
na zbiory otwarte w $\C^n\cong\R^{2n}$), takich że odwzorowania przejścia
\[
\varphi_\beta\circ\varphi_\alpha^{-1}\colon\varphi_\alpha(U_\alpha\cap U_\beta)\to
\varphi_\beta(U_\alpha\cap U_\beta)
\]
są \emph{holomorficzne} (jako odwzorowania między otwartymi podzbiorami $\C^n$), wraz z
maksymalnym takim atlasem (konstrukcja atlasu maksymalnego przebiega dosłownie tak samo,
jak w Sekcji~\ref{sec:atlas}, z ,,gładki'' zastąpionym przez ,,holomorficzny'' — złożenie odwzorowań
holomorficznych jest holomorficzne, więc Lemat~\ref{lem:zgodnosc-przechodnia} i
Twierdzenie~\ref{tw:atlas-maksymalny} przenoszą się bez zmian).
\end{definicja}

\begin{uwaga}
Utożsamiając $\C^n\cong\R^{2n}$ (przez część rzeczywistą i urojoną każdej współrzędnej)
oraz przypominając, że każde odwzorowanie holomorficzne jest w szczególności odwzorowaniem
gładkim (klasy $C^\infty$) w sensie rzeczywistym (ze wzoru całkowego
Cauchy'ego na małym polidysku wynika lokalne rozwinięcie w szereg potęgowy,
który można różniczkować wyraz po wyrazie), widzimy, że \emph{każda rozmaitość zespolona
wymiaru $n$ jest w naturalny sposób rzeczywistą rozmaitością gładką wymiaru $2n$}. Implikacja
odwrotna jest fałszywa: nie każda parzystowymiarowa rozmaitość gładka dopuszcza strukturę
zespoloną (a nawet jeśli dopuszcza, struktura taka zwykle nie jest jednoznaczna).
\end{uwaga}

\begin{przyklad}[Zespolona przestrzeń rzutowa]
Niech $\C_*:=\C\setminus\{0\}$ oznacza zbiór niezerowych
liczb zespolonych.
Zespolonym odpowiednikiem $\mathbb{RP}^n$ z
Przykładu~\ref{prz:RPn} jest przestrzeń ilorazowa
\[
\mathbb{CP}^n
:=
(\C^{n+1}\setminus\{0\})/\!\sim,
\]
gdzie
\[
z\sim w
\iff
\exists\,\lambda\in\C_*
\quad\text{takie, że}\quad
z=\lambda w.
\]
Przestrzeń $\mathbb{CP}^n$ wyposażamy w topologię ilorazową.
Punkty zapisujemy w postaci
\[
[z^0:z^1:\dots:z^n],
\]
przy czym
\[
[z^0:\dots:z^n]
=
[\lambda z^0:\dots:\lambda z^n]
\qquad
\text{dla każdego }\lambda\in\C_*.
\]

Dla $i=0,\dots,n$ definiujemy
\[
U_i
=
\{[z^0:\dots:z^n]\in\mathbb{CP}^n:z^i\neq0\}.
\]
Zbiory $U_i$ są otwarte i pokrywają całą przestrzeń $\mathbb{CP}^n$.
Na każdym z nich definiujemy mapę
\[
\varphi_i:U_i\longrightarrow\C^n
\]
wzorem
\[
\varphi_i([z^0:\dots:z^n])
=
\left(
\frac{z^0}{z^i},
\dots,
\frac{z^{i-1}}{z^i},
\frac{z^{i+1}}{z^i},
\dots,
\frac{z^n}{z^i}
\right).
\]
Jej odwrotność jest dana jawnie przez
\[
\varphi_i^{-1}
(w^0,\dots,w^{i-1},w^{i+1},\dots,w^n)
=
[w^0:\dots:w^{i-1}:1:w^{i+1}:\dots:w^n].
\]
Sprawdźmy także ciągłość, której nie daje sama bijektywność.
Niech $\pi\colon\C^{n+1}\setminus\{0\}\to\mathbb{CP}^n$ będzie
rzutowaniem ilorazowym. Przeciwobraz $\pi^{-1}(U_i)=\{z:z^i\ne0\}$
jest otwarty, zatem $U_i$ jest otwarty. Ograniczenie
$\pi\colon\pi^{-1}(U_i)\to U_i$ także jest odwzorowaniem ilorazowym:
otwarty w $\pi^{-1}(U_i)$ przeciwobraz podzbioru $U_i$ jest otwarty
w całej dziedzinie $\pi$, więc sam podzbiór jest otwarty w ilorazie.
Składowe $\varphi_i\circ\pi$ są ciągłymi ilorazami współrzędnych;
z własności topologii ilorazowej wynika ciągłość $\varphi_i$.
Podany wzór na $\varphi_i^{-1}$ jest złożeniem ciągłego wstawienia
współrzędnej $1$ z $\pi$. Każda $\varphi_i$ jest więc homeomorfizmem
pomiędzy $U_i$ a $\C^n$.

Na przecięciu $U_i\cap U_j$ ten sam punkt przestrzeni
$\mathbb{CP}^n$ posiada dwa różne opisy współrzędnościowe.
Przejście pomiędzy nimi jest opisane przez odwzorowanie
\[
\varphi_j\circ\varphi_i^{-1}.
\]
Jego składowe są ilorazami wielomianów zespolonych, przy czym
mianowniki nie znikają na dziedzinie odwzorowania. Odwzorowania
przejścia są zatem holomorficzne. Aksjomat Hausdorffa i istnienie
przeliczalnej bazy topologii sprawdzimy w dowodzie
twierdzenia~\ref{thm:rzutowe-zwarte-1}. Łącznie daje to strukturę
rozmaitości zespolonej wymiaru $n$ na $\mathbb{CP}^n$, a więc również
rzeczywistej rozmaitości gładkiej wymiaru $2n$.

Rozważmy teraz szczególny przypadek $n=1$. Otrzymujemy
\emph{sferę Riemanna}
\[
\mathbb{CP}^1.
\]
Mamy wówczas dwa zbiory
\[
U_0=\{[z^0:z^1]:z^0\neq0\},
\qquad
U_1=\{[z^0:z^1]:z^1\neq0\},
\]
oraz mapy
\[
\varphi_0:U_0\longrightarrow\C,
\qquad
\varphi_0([z^0:z^1])=\frac{z^1}{z^0},
\]
i
\[
\varphi_1:U_1\longrightarrow\C,
\qquad
\varphi_1([z^0:z^1])=\frac{z^0}{z^1}.
\]
Każda z tych map identyfikuje odpowiedni fragment
$\mathbb{CP}^1$ z kopią płaszczyzny zespolonej $\C$.

Na przecięciu
\[
U_0\cap U_1
=
\{[z^0:z^1]:z^0\neq0,\ z^1\neq0\}
\]
ten sam punkt ma w pierwszej mapie współrzędną
\[
w
=
\varphi_0([z^0:z^1])
=
\frac{z^1}{z^0},
\]
natomiast w drugiej mapie współrzędną
\[
v
=
\varphi_1([z^0:z^1])
=
\frac{z^0}{z^1}.
\]
Stąd
\[
v=\frac1w.
\]
Odwzorowanie przejścia ma więc postać
\[
\varphi_1\circ\varphi_0^{-1}(w)
=
\frac1w,
\qquad
w\in\C_*.
\]

Możemy zatem patrzeć na $\mathbb{CP}^1$ jako na przestrzeń otrzymaną
z dwóch rozłącznych kopii płaszczyzny zespolonej,
\[
\C_0\sqcup\C_1,
\]
przez utożsamienie punktów należących do ich niezerowych części zgodnie
z regułą
\[
w\in\C_0\setminus\{0\}
\quad\sim\quad
\frac1w\in\C_1\setminus\{0\}.
\]
W tym właśnie sensie mówi się, że dwie kopie $\C$ są
\emph{sklejone za pomocą odwzorowania przejścia}
\[
w\longmapsto\frac1w.
\]
Nie sklejamy więc całych płaszczyzn punkt po punkcie, lecz utożsamiamy
jedynie te punkty obu kopii $\C$, które odpowiadają temu samemu punktowi
przestrzeni $\mathbb{CP}^1$ na części wspólnej $U_0\cap U_1$.

Punkt $0\in\C_0$ odpowiada punktowi
\[
[1:0]\in\mathbb{CP}^1,
\]
natomiast punkt $0\in\C_1$ odpowiada punktowi
\[
[0:1]\in\mathbb{CP}^1.
\]
W szczególności, patrząc wyłącznie przez mapę $\varphi_0$, punkt
$[0:1]$ można interpretować jako dodatkowy punkt w nieskończoności.
Dlatego często zapisuje się
\[
\mathbb{CP}^1
\cong
\C\cup\{\infty\}.
\]

Przestrzeń $\mathbb{CP}^1$ jest dyfeomorficzna ze sferą $S^2$.
Trzeba jednak uważać na wybór zespolonych współrzędnych
stereograficznych. Zwykłe rzeczywiste rzuty stereograficzne z
Przykładu~\ref{prz:sfera}, po utożsamieniu $\R^2=\C$, prowadzą do
odwzorowania przejścia
\[
w\longmapsto\frac{w}{|w|^2}
=
\frac1{\overline{w}},
\]
które nie jest holomorficzne.

Aby otrzymać strukturę zespoloną zgodną ze strukturą na
$\mathbb{CP}^1$, wybieramy na $S^2$ współrzędne
\[
\zeta_N(x)
=
\frac{x^1+ix^2}{1-x^3},
\qquad
\zeta_S(x)
=
\frac{x^1-ix^2}{1+x^3}.
\]
Wtedy na części wspólnej zachodzi
\[
\zeta_S\circ\zeta_N^{-1}(w)
=
\frac1w.
\]

Jawny dyfeomorfizm
\[
H:\mathbb{CP}^1\longrightarrow S^2
\]
jest dany wzorem
\[
H([z^0:z^1])
=
\frac{
\big(
2\operatorname{Re}(\overline{z^0}z^1),
2\operatorname{Im}(\overline{z^0}z^1),
|z^1|^2-|z^0|^2
\big)
}{
|z^0|^2+|z^1|^2
}.
\]
Skalowanie obu współrzędnych przez $\lambda\neq0$ mnoży licznik
i mianownik przez $|\lambda|^2$, dlatego wzór jest dobrze określony
na klasach równoważności.

Ponadto
\[
4|z^0|^2|z^1|^2
+
\big(|z^1|^2-|z^0|^2\big)^2
=
\big(|z^0|^2+|z^1|^2\big)^2,
\]
co pokazuje, że obraz odwzorowania $H$ leży na sferze $S^2$.

W wybranych współrzędnych mamy
\[
\zeta_N\circ H=\varphi_0,
\qquad
\zeta_S\circ H=\varphi_1.
\]
Pierwsza z tych równości daje bijekcję
\[
U_0\longrightarrow S^2\setminus\{N\},
\]
natomiast pozostały punkt
\[
[0:1]
\]
przechodzi na biegun północny $N$.
W lokalnych współrzędnych odwzorowanie $H$ oraz jego odwrotność są
tożsamościami, co pokazuje, że
\[
\mathbb{CP}^1\cong S^2
\]
jako rozmaitości gładkie.
\end{przyklad}


\begin{figure}[h]
\centering
\begin{tikzpicture}[>=Stealth,scale=1]
% Sfera: jasna plama i podział równika na część widoczną oraz ukrytą.
\shade[inner color=white,outer color=manifoldblue!43] (0,0) circle (1.7);
\draw[manifoldblue!75!black,very thick] (0,0) circle (1.7);
\draw[manifoldblue!52,densely dashed,thin]
  (1.7,0) arc[start angle=0,end angle=180,x radius=1.7,y radius=.42];
\draw[manifoldblue!65!black,thin]
  (-1.7,0) arc[start angle=180,end angle=360,x radius=1.7,y radius=.42];
\draw[manifoldblue!40,densely dashed,thin] (0,1.7)
  .. controls (-.52,1.12) and (-.52,-1.12) .. (0,-1.7);
\draw[manifoldblue!65!black,thin] (0,1.7)
  .. controls (.52,1.12) and (.52,-1.12) .. (0,-1.7);
% bieguny
\fill[manifoldgold!85!black] (0,1.7) circle (2.4pt) node[above,font=\small]{$N$};
\fill[manifoldgold!85!black] (0,-1.7) circle (2.4pt) node[below,font=\small]{$S$};
\node[manifoldblue!70!black,font=\small] at (0,-2.35) {$\mathbb{CP}^1\cong S^2$};
% --- strzałki do płaszczyzn C ---
\draw[->,very thick,manifoldred] (1.85,0.6) -- (3.2,1.1);
\draw[->,very thick,manifoldred] (1.85,-0.6) -- (3.2,-1.1);
% --- dwie kopie C (zaokrąglone) ---
% górna
\shade[inner color=manifoldgreen!20, outer color=manifoldgreen!5, rounded corners=8pt]
  (3.5,0.5) rectangle (7.0,1.8);
\draw[manifoldgreen!60!black,thick,rounded corners=8pt]
  (3.5,0.5) rectangle (7.0,1.8);
\begin{scope}
  \clip[rounded corners=8pt] (3.5,0.5) rectangle (7.0,1.8);
  \draw[step=0.4,manifoldgreen!22,thin] (3.5,0.5) grid (7.0,1.8);
\end{scope}
\node[font=\small] at (5.25,1.15) {$\C$ (mapa $\zeta_N$)};
% dolna
\shade[inner color=manifoldgreen!20, outer color=manifoldgreen!5, rounded corners=8pt]
  (3.5,-1.8) rectangle (7.0,-0.5);
\draw[manifoldgreen!60!black,thick,rounded corners=8pt]
  (3.5,-1.8) rectangle (7.0,-0.5);
\begin{scope}
  \clip[rounded corners=8pt] (3.5,-1.8) rectangle (7.0,-0.5);
  \draw[step=0.4,manifoldgreen!22,thin] (3.5,-1.8) grid (7.0,-0.5);
\end{scope}
\node[font=\small] at (5.25,-1.15) {$\C$ (mapa $\zeta_S$)};
% przejście w↦1/w
\draw[<->,thick,manifoldred] (7.15,1.15) -- (7.15,-1.15);
\node[manifoldred,font=\small,right] at (7.2,0) {$w\mapsto 1/w$};
\node[font=\small] at (5.25,-2.25) {dwie kopie $\C$ sklejone przejściem $w\mapsto 1/w$};
\end{tikzpicture}
\caption{Schematyczna interpretacja $\mathbb{CP}^1$ jako sfery Riemanna:
dwie mapy zespolone pokrywają sferę; w mapie południowej odwrócono
znak drugiej współrzędnej rzeczywistej. Przejście $1/w$ jest określone dla $w\ne0$.}
\label{fig:cp1}
\end{figure}

\section{Odwzorowania gładkie między rozmaitościami}

\begin{definicja}
Niech $M$, $N$ będą rozmaitościami gładkimi wymiarów odpowiednio $m$, $n$. Odwzorowanie
$F\colon M\to N$ jest \emph{gładkie} (klasy $C^\infty$) w punkcie $p\in M$, jeżeli dla
pewnej (równoważnie: każdej) mapy dopuszczalnej $(U,\varphi)$ wokół $p$ i mapy
dopuszczalnej $(V,\psi)$ wokół $F(p)$ z $F(U)\subseteq V$, odwzorowanie
\[
\psi\circ F\circ\varphi^{-1}\colon\varphi(U)\to\psi(V)
\]
jest gładkie (klasy $C^\infty$) w pewnym otwartym otoczeniu
$\varphi(p)$. Odwzorowanie $F$ jest gładkie, jeśli jest gładkie w każdym punkcie $p\in M$.
Zbiór wszystkich gładkich funkcji $M\to\R$ oznaczamy $C^\infty(M)$.
\end{definicja}
\begin{figure}[ht]
\centering
\begin{tikzpicture}[>=Stealth,scale=.9,line cap=round,line join=round]
% Dwie pofalowane powierzchnie z siatką lokalnych współrzędnych.
\def\Mshape{(-4.75,.2) .. controls (-4.65,1.55) and (-3.65,2.0) .. (-2.7,1.56)
 .. controls (-1.65,1.12) and (-.7,1.75) .. (-.88,.36)
 .. controls (-1.03,-.78) and (-2.15,-.36) .. (-2.8,-.83)
 .. controls (-3.63,-1.25) and (-4.78,-.85) .. cycle}
\def\Nshape{(1.05,.12) .. controls (1.4,1.35) and (2.5,1.95) .. (3.45,1.58)
 .. controls (4.35,1.23) and (5.1,1.35) .. (5.02,.2)
 .. controls (4.98,-.93) and (4.1,-.87) .. (3.2,-.65)
 .. controls (2.35,-.45) and (1.42,-1.05) .. cycle}
\fill[manifoldblue!12] \Mshape;
\begin{scope}
\clip \Mshape;
\fill[manifoldblue!24,opacity=.65] (-2.65,.35) ellipse (1.02 and .55);
\end{scope}
\draw[manifoldblue!75!black,very thick] \Mshape;
\draw[manifoldblue!70!black,densely dashed] (-3.67,.18) .. controls (-3.45,.88) and (-2.28,.93) .. (-1.78,.18);
\fill[manifoldgreen!11] \Nshape;
\begin{scope}
\clip \Nshape;
\fill[manifoldgreen!25,opacity=.65] (3.03,.34) ellipse (1.03 and .53);
\end{scope}
\draw[manifoldgreen!65!black,very thick] \Nshape;
\draw[manifoldgreen!65!black,densely dashed] (1.96,.14) .. controls (2.2,.91) and (3.54,.9) .. (4.08,.12);
\node[font=\large\bfseries,manifoldblue!75!black] at (-4.22,1.38) {$M$};
\node[font=\large\bfseries,manifoldgreen!65!black] at (4.45,1.27) {$N$};
\node[font=\small,manifoldblue!75!black] at (-3.49,.7) {$U$};
\node[font=\small,manifoldgreen!65!black] at (3.8,.7) {$V$};
\fill[manifoldred] (-2.67,.34) circle (2.8pt) node[below,font=\small] {$p$};
\fill[manifoldred] (3.02,.33) circle (2.8pt) node[below,font=\small] {$F(p)$};
\draw[->,very thick,manifoldred] (-.72,.78) .. controls (-.04,1.75) and (.58,1.7) .. (1.13,.79)
  node[midway,above,font=\large\bfseries] {$F$};
% Dwie euklidesowe mapy z zaznaczonymi odpowiadającymi sobie punktami.
\fill[manifoldblue!5,rounded corners=5pt] (-4.28,-3.72) rectangle (-1.12,-2.4);
\draw[manifoldblue!65!black,thick,rounded corners=5pt] (-4.28,-3.72) rectangle (-1.12,-2.4);
\fill[manifoldgreen!5,rounded corners=5pt] (1.18,-3.72) rectangle (4.34,-2.4);
\draw[manifoldgreen!65!black,thick,rounded corners=5pt] (1.18,-3.72) rectangle (4.34,-2.4);
\draw[->,thick,manifoldblue!75!black] (-2.68,-.87)--(-2.68,-2.34)
  node[midway,left,font=\small] {$\varphi$};
\draw[->,thick,manifoldgreen!70!black] (3.03,-.84)--(3.03,-2.34)
  node[midway,right,font=\small] {$\psi$};
\fill[manifoldred] (-2.68,-3.03) circle (2.5pt) node[below,font=\scriptsize] {$u=\varphi(p)$};
\fill[manifoldred] (3.03,-3.03) circle (2.5pt) node[below,font=\scriptsize] {$f(u)=\psi(F(p))$};
\draw[->,very thick,manifoldred] (-1.0,-3.02)--(1.03,-3.02);
\node[font=\small\bfseries,manifoldred!85!black] at (0,-2.71) {$f$};
\node[font=\small] at (-2.7,-4.15) {$\varphi(U)\subset\R^m$};
\node[font=\small] at (2.76,-4.15) {$\psi(V)\subset\R^n$};
\end{tikzpicture}
\caption{Odwzorowanie $F\colon M\to N$ oraz jego lokalny zapis współrzędnościowy. Po wybraniu map $(U,\varphi)$ i $(V,\psi)$ gładkość $F$ bada się poprzez zwykłe odwzorowanie $\psi\circ F\circ\varphi^{-1}$ między otwartymi podzbiorami przestrzeni euklidesowych.}
\label{fig:odwzorowanie-gładkie}
\end{figure}

\begin{uwaga}
Definicja gładkości nie zależy od wyboru map. Niech $(U,\varphi)$ oraz
$(V,\psi)$ będą mapami wokół odpowiednio $p$ i $F(p)$, dla których
odwzorowanie
\[
\psi\circ F\circ\varphi^{-1}
\]
jest gładkie. Jeśli $(U',\varphi')$, $(V',\psi')$ są innymi mapami
dopuszczalnymi wokół $p$ i $F(p)$, to chcemy porównać oba opisy
współrzędnościowe funkcji $F$.

Najpierw zauważmy, że na odpowiednio małym otoczeniu punktu $p$ mamy
\[
F
=
\psi^{-1}\circ
(\psi\circ F\circ\varphi^{-1})
\circ\varphi.
\]
Ponieważ $\varphi$ i $\psi^{-1}$ są homeomorfizmami, a odwzorowanie
$\psi\circ F\circ\varphi^{-1}$ jest gładkie, a więc w szczególności
ciągłe, otrzymujemy, że $F$ jest ciągłe w pewnym otoczeniu punktu $p$.

Ponieważ
\[
p\in U\cap U'
\qquad\text{oraz}\qquad
F(p)\in V\cap V',
\]
a $V\cap V'$ jest zbiorem otwartym, wybierzmy otwarte otoczenie
$O$ punktu $p$, na którym $F$ jest ciągłe. Wówczas
\[
(F|_O)^{-1}(V\cap V')=O\cap F^{-1}(V\cap V')
\]
jest otwartym otoczeniem $p$ w $M$. Możemy więc zmniejszyć dziedzinę
do otwartego otoczenia
\[
W\subseteq O\cap U\cap U'\cap F^{-1}(V\cap V').
\]
Na $W$ jednocześnie obowiązują oba układy współrzędnych, a obraz
$F(W)$ jest zawarty w $V\cap V'$. Wszystkie poniższe złożenia są zatem
dobrze określone.

Na odpowiadającej zbiorowi $W$ dziedzinie mamy
\[
\psi'\circ F\circ(\varphi')^{-1}
=
(\psi'\circ\psi^{-1})
\circ
(\psi\circ F\circ\varphi^{-1})
\circ
(\varphi\circ(\varphi')^{-1}).
\]
Odwzorowania
\[
\varphi\circ(\varphi')^{-1}
\qquad\text{oraz}\qquad
\psi'\circ\psi^{-1}
\]
są odwzorowaniami przejścia, a więc są gładkie z definicji zgodności
map. Co więcej, są one dyfeomorfizmami pomiędzy odpowiednimi otwartymi
podzbiorami przestrzeni euklidesowych. Zatem złożenie z nimi nie zmienia
gładkości. W konsekwencji
\[
\psi'\circ F\circ(\varphi')^{-1}
\]
jest gładkie wtedy i tylko wtedy, gdy gładkie jest
\[
\psi\circ F\circ\varphi^{-1}.
\]
Definicja gładkości odwzorowania $F$ nie zależy więc od wyboru map
w dziedzinie i przeciwdziedzinie.
\end{uwaga}

\begin{definicja}\label{def dyfeo}
Odwzorowanie gładkie $F\colon M\to N$ jest \emph{dyfeomorfizmem}, jeśli jest bijekcją, a
$F^{-1}\colon N\to M$ jest również gładkie. Rozmaitości $M$, $N$ nazywamy
\emph{dyfeomorficznymi}, jeśli istnieje między nimi dyfeomorfizm.
\end{definicja}

\begin{stwierdzenie}
Złożenie odwzorowań gładkich jest gładkie: jeśli $F\colon M\to N$ i $G\colon N\to P$ są
gładkie, to $G\circ F\colon M\to P$ jest gładkie.
\end{stwierdzenie}

\begin{proof}
Ustalmy $p\in M$ i mapy dopuszczalne $(U,\varphi)$ wokół $p$, $(V,\psi)$ wokół $F(p)$,
$(W,\eta)$ wokół $G(F(p))$, z $F(U)\subseteq V$, $G(V)\subseteq W$ (istnienie takich map —
po ewentualnym zmniejszeniu $V$, a następnie $U$ — wynika z ciągłości $F$ i $G$, którą zapewnia gładkość).
Wówczas na $\varphi(U)$
\[
\eta\circ(G\circ F)\circ\varphi^{-1}=(\eta\circ G\circ\psi^{-1})\circ(\psi\circ F\circ
\varphi^{-1}),
\]
czyli złożenie dwóch odwzorowań gładkich (między otwartymi podzbiorami przestrzeni
euklidesowych) — a takie złożenie jest gładkie na mocy zwykłego twierdzenia analizy wielu
zmiennych o różniczkowaniu złożenia.
\end{proof}
\section{Przykłady rozmaitości}

\begin{stwierdzenie}\label{stw:otwarty-podzbior}
Jeśli $M$ jest rozmaitością gładką wymiaru $n$ i $\Omega\subseteq M$ jest zbiorem otwartym,
to $\Omega$ z topologią i strukturą gładką odziedziczoną po $M$ (mapy $(U\cap\Omega,
\varphi|_{U\cap\Omega})$ dla map dopuszczalnych $(U,\varphi)$ na $M$) jest rozmaitością
gładką wymiaru $n$.
\end{stwierdzenie}

\begin{proof}
Warunki Hausdorffa i drugiej przeliczalności dziedziczą się automatycznie na
podprzestrzenie. Rodzina $\{(U\cap\Omega,\varphi|_{U\cap\Omega})\}$, gdzie $(U,\varphi)$
przebiega atlas maksymalny na $M$, pokrywa $\Omega$ (bo mapy pokrywają $M\supseteq\Omega$)
i składa się z homeomorfizmów na zbiory otwarte $\varphi(U\cap\Omega)\subseteq\varphi(U)
\subseteq\R^n$ (obraz zbioru otwartego $U\cap\Omega$ przez homeomorfizm $\varphi$ jest
otwarty w $\varphi(U)$, a więc i w $\R^n$). Zgodność map $(U\cap\Omega,\varphi|_{U\cap
\Omega})$ i $(U'\cap\Omega,\varphi'|_{U'\cap\Omega})$ wynika natychmiast ze zgodności
$(U,\varphi)$ i $(U',\varphi')$, gdyż odwzorowania przejścia dla map obciętych są
obcięciami odwzorowań przejścia dla map oryginalnych do (otwartego) podzbioru dziedziny.
\end{proof}

\begin{przyklad}
Zbiór $\mathrm{GL}(n,\R)$ macierzy odwracalnych $n\times n$ o wyrazach rzeczywistych jest
otwartym podzbiorem przestrzeni $M_n(\R)\cong\R^{n^2}$ wszystkich macierzy $n\times n$ (bo
$\det\colon M_n(\R)\to\R$ jest funkcją ciągłą, a $\mathrm{GL}(n,\R)=\det^{-1}(\R\setminus
\{0\})$ jest przeciwobrazem zbioru otwartego). Na mocy Stwierdzenia~\ref{stw:otwarty-podzbior},
$\mathrm{GL}(n,\R)$ jest rozmaitością gładką wymiaru $n^2$.
\end{przyklad}

\begin{stwierdzenie}[Rozmaitość produktowa]
Jeśli $M$, $N$ są rozmaitościami gładkimi wymiarów odpowiednio $m$, $n$, to $M\times N$
(z topologią produktową) jest rozmaitością gładką wymiaru $m+n$, z atlasem złożonym z map
$(U\times V,\varphi\times\psi)$, gdzie $(U,\varphi)$ jest mapą na $M$, a $(V,\psi)$ mapą na
$N$, i $(\varphi\times\psi)(p,q):=(\varphi(p),\psi(q))\in\R^m\times\R^n\cong\R^{m+n}$.
\end{stwierdzenie}

\begin{proof}
Hausdorffowość i druga przeliczalność produktu dwóch przestrzeni o tych własnościach są
standardowymi faktami z topologii ogólnej. Mapy $\varphi\times\psi$ są homeomorfizmami na
zbiory otwarte jako produkty homeomorfizmów na zbiory otwarte, a ich dziedziny pokrywają
$M\times N$. Jeśli $(U_1\times V_1,\varphi_1\times\psi_1)$ i $(U_2\times V_2,\varphi_2
\times\psi_2)$ są dwiema takimi mapami, ich odwzorowanie przejścia to
\[
(\varphi_2\times\psi_2)\circ(\varphi_1\times\psi_1)^{-1}=(\varphi_2\circ\varphi_1^{-1})
\times(\psi_2\circ\psi_1^{-1}),
\]
które jest gładkie jako produkt dwóch odwzorowań gładkich. Zatem otrzymana rodzina map
jest atlasem.
\end{proof}

\begin{przyklad}[Sfera $S^n$ przez rzuty stereograficzne]\label{prz:sfera}
Niech $S^n=\{x\in\R^{n+1}:\|x\|=1\}$ z topologią podprzestrzeni $\R^{n+1}$. Oznaczmy przez
$N=(0,\dots,0,1)$ i $S=(0,\dots,0,-1)$ bieguny północny i południowy. Definiujemy rzut
stereograficzny z bieguna północnego,
\[
\varphi_N\colon S^n\setminus\{N\}\to\R^n, \qquad
\varphi_N(x^1,\dots,x^{n+1})=\frac{1}{1-x^{n+1}}(x^1,\dots,x^n),
\]
tzn.\ $\varphi_N(x)$ to punkt przecięcia prostej przechodzącej przez $N$ i $x$ z
hiperpłaszczyzną $\{x^{n+1}=0\}\cong\R^n$. Odwzorowanie odwrotne wyraża się wzorem
\[
\varphi_N^{-1}(u)=\frac1{\|u\|^2+1}\big(2u,\ \|u\|^2-1\big)\in\R^n\times\R,
\]
co sprawdzamy bezpośrednim rachunkiem: punkt $y=\varphi_N^{-1}(u)$ leży na prostej
$t\mapsto(tu,1-t)$ (parametryzacja prostej przez $N=(0,1)$ w kierunku $(u,0)-N=(u,-1)$) i
spełnia $\|y\|=1$; rozwiązując $t^2\|u\|^2+(1-t)^2=1$, tzn.\ $t\big(t(\|u\|^2+1)-2\big)=0$,
względem $t\neq0$, otrzymujemy $t=2/(\|u\|^2+1)$, skąd powyższy wzór. Analogicznie
definiujemy $\varphi_S\colon S^n\setminus\{S\}\to\R^n$, rzut z bieguna południowego, wzorem
$\varphi_S(x)=\frac1{1+x^{n+1}}(x^1,\dots,x^n)$.

Zbiory $S^n\setminus\{N\}$ i $S^n\setminus\{S\}$ pokrywają $S^n$, więc pozostaje sprawdzić
gładkość odwzorowania przejścia na $\varphi_N(S^n\setminus\{N,S\})=\R^n\setminus\{0\}$.
Podstawiając $\varphi_N^{-1}(u)$ do wzoru na $\varphi_S$, i licząc osobno mianownik
\[
1+\frac{\|u\|^2-1}{\|u\|^2+1}=\frac{2\|u\|^2}{\|u\|^2+1},
\]
otrzymujemy
\[
(\varphi_S\circ\varphi_N^{-1})(u)=\frac{2u/(\|u\|^2+1)}{2\|u\|^2/(\|u\|^2+1)}=
\frac{u}{\|u\|^2},
\]
czyli inwersja względem sfery jednostkowej. Odwzorowanie $u\mapsto u/\|u\|^2$ jest gładkie
na $\R^n\setminus\{0\}$, a jest ono swoją własną odwrotnością (dla $w=u/\|u\|^2$ mamy
$\|w\|^2=1/\|u\|^2$, więc $w/\|w\|^2=(u/\|u\|^2)\cdot\|u\|^2=u$), więc obie mapy są
zgodne. Zatem $\{(S^n\setminus\{N\},\varphi_N),(S^n\setminus\{S\},\varphi_S)\}$ jest
atlasem na $S^n$, nadającym jej strukturę rozmaitości gładkiej wymiaru $n$.
\end{przyklad}

\begin{figure}[ht]
\centering
\begin{tikzpicture}[>=Stealth,scale=1.65,line cap=round,line join=round]
\shade[inner color=white,outer color=manifoldblue!43] (0,0) circle (1);
\draw[manifoldblue!75!black,thick] (0,0) circle (1);
\draw[->,black!55] (0,-1.2)--(0,1.38) node[above,font=\small] {$x^{n+1}$};
\draw[->,manifoldgreen!65!black,thick] (-1.3,0)--(3.55,0);
\node[below,font=\small,manifoldgreen!55!black] at (2.15,-.23)
 {$\{x^{n+1}=0\}\cong\R^n$};
\coordinate (N) at (0,1);
\coordinate (X) at (.6,.8);
\coordinate (P) at (3,0);
\draw[manifoldred,very thick] (N)--(X);
\draw[manifoldred,thick,densely dashed] (X)--(P);
\foreach \name in {N,X,P}{\fill[manifoldred] (\name) circle (1.6pt);}
\node[above left,font=\small] at (N) {$N$};
\node[above right,font=\small,manifoldred] at (X) {$x$};
\node[above right,font=\small,manifoldred] at (P) {$\varphi_N(x)$};
\node[below left,font=\small] at (0,0) {$0$};
\node[manifoldblue!75!black,font=\small] at (-.73,.82) {$S^n$};
\end{tikzpicture}
\caption{Rzut stereograficzny na płaszczyznę równikową, zgodny ze wzorem
w tekście. Pokazano przekrój płaszczyzną zawierającą oś biegunową
i punkt $x$. Punkty $N,x,\varphi_N(x)$ leżą na jednej prostej;
punkt przecięcia z $\{x^{n+1}=0\}$ ma współrzędne
$\varphi_N(x)=(x^1,\ldots,x^n)/(1-x^{n+1})$.}
\label{fig:stereo}
\end{figure}


\begin{przyklad}[Rzeczywista przestrzeń rzutowa]\label{prz:RPn}
Niech $\R_*:=\R\setminus\{0\}$ oznacza zbiór niezerowych
liczb rzeczywistych.
Niech
\[
\mathbb{RP}^n
=
(\R^{n+1}\setminus\{0\})/\!\sim,
\]
gdzie
\[
x\sim y
\iff
\exists\,\lambda\in\R_*
\quad\text{takie, że}\quad
x=\lambda y.
\]
Przestrzeń $\mathbb{RP}^n$ wyposażamy w topologię ilorazową względem
rzutowania
\[
\pi:\R^{n+1}\setminus\{0\}\longrightarrow\mathbb{RP}^n,
\qquad
\pi(x)=[x].
\]
Punkty przestrzeni rzutowej zapisujemy w postaci
\[
[x^0:\dots:x^n].
\]

Dla $i=0,\dots,n$ definiujemy
\[
U_i
=
\{[x^0:\dots:x^n]\in\mathbb{RP}^n:x^i\neq0\}.
\]
Najpierw należy sprawdzić, że definicja ta nie zależy od wyboru
reprezentanta klasy $[x]$. Jeżeli $x\sim y$, to istnieje
$\lambda\neq0$ takie, że $y=\lambda x$, a więc
\[
y^i=\lambda x^i.
\]
Ponieważ $\lambda\neq0$, mamy
\[
x^i\neq0
\iff
y^i\neq0.
\]
Warunek $x^i\neq0$ zależy zatem wyłącznie od klasy $[x]$, więc zbiór
$U_i$ jest dobrze określony.

Zauważmy następnie, że
\[
\pi^{-1}(U_i)
=
\{x\in\R^{n+1}\setminus\{0\}:x^i\neq0\}.
\]
Jest to zbiór otwarty w $\R^{n+1}\setminus\{0\}$, ponieważ jest
przeciwobrazem otwartego zbioru
$\R_*$ przez ciągłą funkcję współrzędnościową
\[
x\longmapsto x^i.
\]
Z definicji topologii ilorazowej zbiór $A\subseteq\mathbb{RP}^n$
jest otwarty wtedy i tylko wtedy, gdy $\pi^{-1}(A)$ jest otwarty
w $\R^{n+1}\setminus\{0\}$. Stąd każdy $U_i$ jest otwarty.

Ponieważ każdy niezerowy wektor
$x=(x^0,\dots,x^n)$ ma przynajmniej jedną niezerową współrzędną,
zbiory $U_i$ pokrywają całą przestrzeń:
\[
\mathbb{RP}^n
=
\bigcup_{i=0}^n U_i.
\]

Na $U_i$ definiujemy odwzorowanie
\[
\varphi_i:U_i\longrightarrow\R^n
\]
wzorem
\[
\varphi_i([x^0:\dots:x^n])
=
\left(
\frac{x^0}{x^i},
\dots,
\frac{x^{i-1}}{x^i},
\frac{x^{i+1}}{x^i},
\dots,
\frac{x^n}{x^i}
\right),
\]
czyli po znormalizowaniu reprezentanta tak, aby jego $i$-ta
współrzędna była równa $1$, pomijamy tę współrzędną.

Także tutaj należy sprawdzić niezależność od wyboru reprezentanta.
Jeżeli $y=\lambda x$ dla pewnego $\lambda\neq0$, to dla każdego
$j\neq i$ mamy
\[
\frac{y^j}{y^i}
=
\frac{\lambda x^j}{\lambda x^i}
=
\frac{x^j}{x^i}.
\]
Wartość $\varphi_i([x])$ zależy więc wyłącznie od klasy $[x]$.

Odwzorowanie odwrotne jest dane jawnie przez wstawienie $1$
na $i$-te miejsce:
\[
\varphi_i^{-1}
(u^0,\dots,u^{i-1},u^{i+1},\dots,u^n)
=
[u^0:\dots:u^{i-1}:1:u^{i+1}:\dots:u^n].
\]
W szczególności $\varphi_i$ jest bijekcją.

Odwrotność $\varphi_i^{-1}$ jest ciągła, ponieważ jest złożeniem
ciągłego odwzorowania
\[
(u^0,\dots,u^{i-1},u^{i+1},\dots,u^n)
\longmapsto
(u^0,\dots,u^{i-1},1,u^{i+1},\dots,u^n)
\]
z ciągłym rzutowaniem ilorazowym $\pi$.

Aby wykazać ciągłość $\varphi_i$, zauważmy, że na zbiorze
\[
\pi^{-1}(U_i)=\{x^i\neq0\}
\]
mamy
\[
\varphi_i\circ\pi(x)
=
\left(
\frac{x^0}{x^i},
\dots,
\frac{x^{i-1}}{x^i},
\frac{x^{i+1}}{x^i},
\dots,
\frac{x^n}{x^i}
\right),
\]
a więc $\varphi_i\circ\pi$ jest ciągłe. Ponieważ obcięcie
\[
\pi|_{\pi^{-1}(U_i)}:
\pi^{-1}(U_i)\longrightarrow U_i
\]
jest odwzorowaniem ilorazowym, wynika stąd ciągłość $\varphi_i$.
Otrzymujemy zatem homeomorfizm
\[
\varphi_i:U_i\longrightarrow\R^n.
\]

Pozostaje sprawdzić zgodność map. Dla $i\neq j$ na przecięciu
$U_i\cap U_j$ odwzorowanie przejścia
\[
\varphi_j\circ\varphi_i^{-1}
\]
jest określone tam, gdzie współrzędna odpowiadająca ilorazowi
$x^j/x^i$ jest niezerowa. Każda jego składowa jest albo odwrotnością
tej współrzędnej, albo ilorazem jednej ze współrzędnych przez nią.
Są to zatem funkcje wymierne o nieznikającym na dziedzinie
mianowniku, a więc funkcje gładkie.

W konsekwencji
\[
\{(U_i,\varphi_i)\}_{i=0}^n
\]
jest atlasem gładkim. Po sprawdzeniu aksjomatów topologicznych
w twierdzeniu poniżej otrzymujemy, że $\mathbb{RP}^n$ jest
rozmaitością gładką wymiaru $n$.
\end{przyklad}


\begin{przyklad}[Dwie mapy $\mathbb{RP}^2$ krok po kroku]\label{prz:RP2-przejscie-1}
Na $U_0$ każda klasa ma jednoznacznego reprezentanta
\[
[1:u:v],
\]
natomiast na $U_1$ ma reprezentanta
\[
[a:1:b].
\]
Ich część wspólna $U_0\cap U_1$, zapisana we współrzędnych mapy
$\varphi_0$, odpowiada warunkowi $u\neq0$. Aby przejść od współrzędnych
zadanych przez $\varphi_0$ do współrzędnych zadanych przez $\varphi_1$,
dzielimy trzy współrzędne jednorodne przez $u$:
\[
[1:u:v]
=
[1/u:1:v/u].
\]
Stąd odwzorowanie zmiany współrzędnych ma postać
\[
\varphi_1\circ\varphi_0^{-1}(u,v)
=
\left(\frac1u,\frac vu\right).
\]
Oznacza to, że jeżeli punkt $p\in U_0\cap U_1$ ma w pierwszej mapie
współrzędne
\[
\varphi_0(p)=(u,v),
\]
to w drugiej mapie jego współrzędne wynoszą
\[
\varphi_1(p)=(a,b)
=
\left(\frac1u,\frac vu\right).
\]

Dziedziną odwzorowania przejścia jest
\[
\{(u,v)\in\R^2:u\neq0\},
\]
a jego macierz Jacobiego wynosi
\[
D(\varphi_1\circ\varphi_0^{-1})(u,v)
=
\begin{pmatrix}
-u^{-2} & 0\\[1mm]
-vu^{-2} & u^{-1}
\end{pmatrix}.
\]
Ponieważ
\[
\det D(\varphi_1\circ\varphi_0^{-1})(u,v)
=
-u^{-3}\neq0,
\]
a odwrotne przejście ma postać
\[
(a,b)\longmapsto
\left(\frac1a,\frac ba\right),
\]
odwzorowanie przejścia jest dyfeomorfizmem.

Zmianę współrzędnych wektora stycznego w tych mapach obliczymy
w przykładzie~\ref{prz:RP2-styczna-1}, po zdefiniowaniu przestrzeni stycznej.
\end{przyklad}


\begin{uwaga}[Geometryczna interpretacja $\mathbb{RP}^1$ i $\mathbb{RP}^2$]
Przestrzeń $\mathbb{RP}^n$ można interpretować jako zbiór prostych
przechodzących przez początek układu współrzędnych w $\R^{n+1}$.
Rzeczywiście, punkt
\[
[x^0:\dots:x^n]\in\mathbb{RP}^n
\]
odpowiada prostej
\[
\ell_x=\{\lambda x:\lambda\in\R\}.
\]
Zmiana reprezentanta $x$ na $\mu x$, gdzie $\mu\neq0$, nie zmienia tej
prostej. Współrzędne jednorodne opisują więc nie konkretny punkt
$\R^{n+1}$, lecz kierunek prostej przechodzącej przez $0$.

Każda taka prosta przecina sferę jednostkową $S^n$ w dokładnie dwóch
punktach antypodalnych $x$ i $-x$. Stąd
\[
\mathbb{RP}^n\cong S^n/\{x\sim-x\}.
\]
Jest to także identyfikacja topologii. Oznaczmy pierwotne rzutowanie
przez $\pi$ i jego ograniczenie do sfery przez $q$.
Normalizacja $\nu(x)=x/\|x\|$ jest ciągła i $\pi=q\circ\nu$.
Jeśli $q^{-1}(A)$ jest otwarty w sferze, to
$\pi^{-1}(A)=\nu^{-1}(q^{-1}(A))$ jest otwarty, zatem $A$ jest
otwarty w pierwotnej topologii ilorazowej. Odwrotna implikacja
wynika z ciągłości $q$. Stąd $q$ jest odwzorowaniem ilorazowym,
co uzasadnia powyższy homeomorfizm.

Dla $n=1$ otrzymujemy
\[
\mathbb{RP}^1\cong S^1/\{z\sim-z\}.
\]
Po utożsamieniu $S^1$ z okręgiem jednostkowym w $\C$ rozważmy
odwzorowanie
\[
F:S^1\to S^1,\qquad F(z)=z^2.
\]
Ponieważ
\[
F(z)=F(-z),
\]
odwzorowanie to jest stałe na parach antypodalnych i wyznacza
odwzorowanie
\[
\widetilde F:\mathbb{RP}^1\to S^1,
\qquad
\widetilde F([z])=z^2.
\]
Jest ono bijekcją, ponieważ
\[
z^2=w^2
\iff
z=\pm w,
\]
a każdy punkt okręgu posiada dwa przeciwne pierwiastki kwadratowe.
Ciągłość $\widetilde F$ wynika z ciągłości $F=\widetilde F\circ q$
i własności ilorazu. Ponieważ $\mathbb{RP}^1=q(S^1)$ jest zwarta,
a $S^1$ jest przestrzenią Hausdorffa, $\widetilde F$ jest homeomorfizmem.

Jest to również dyfeomorfizm. Jeśli
\[
z=e^{i\theta},
\]
to w $\mathbb{RP}^1$ kąty $\theta$ i $\theta+\pi$ opisują ten sam punkt,
natomiast
\[
\widetilde F([e^{i\theta}])=e^{2i\theta}.
\]
Sprawdźmy zgodność współrzędnych kątowych z podanym wcześniej atlasem.
Dla $[\cos\theta:\sin\theta]$, na małym przedziale z $\cos\theta\ne0$,
współrzędną afiniczną jest $\tan\theta$, o pochodnej $1/\cos^2\theta$.
Gdy $\sin\theta\ne0$, druga mapa daje $\cot\theta$, o pochodnej
$-1/\sin^2\theta$. Pochodne nie znikają, więc przejścia i ich lokalne
odwrotności są gładkie. Kąt jest zatem zgodną lokalną współrzędną.
W lokalnych współrzędnych kątowych odwzorowanie ma więc postać
\[
\theta\longmapsto 2\theta,
\]
a jego lokalna odwrotność
\[
\alpha\longmapsto \frac{\alpha}{2}.
\]
Oba odwzorowania są gładkie, zatem
\[
\mathbb{RP}^1\cong S^1
\]
jako rozmaitości gładkie.

Dla $n=2$ otrzymujemy
\[
\mathbb{RP}^2\cong S^2/\{x\sim-x\}.
\]
Możemy ograniczyć się do domkniętej półsfery północnej. Każda para
antypodalna spoza równika ma dokładnie jeden punkt w jej wnętrzu; na równiku
oba punkty pary należą do półsfery. Ponieważ domknięta półsfera jest
homeomorficzna z dyskiem $D^2$, otrzymujemy model
\[
\mathbb{RP}^2
\cong
D^2/\{x\sim-x\text{ dla }x\in\partial D^2\}.
\]
Oznacza to, że punkty wewnętrzne dysku pozostają różne, natomiast każdy
punkt brzegu utożsamiamy z punktem antypodalnym. Rzutowanie półsfery
na $\mathbb{RP}^2$ jest ciągłą surjekcją i ma dokładnie te włókna.
Wyznacza zatem ciągłą bijekcję z opisanego ilorazu dysku.
Jej dziedzina jest zwarta, a $\mathbb{RP}^2$ jest przestrzenią Hausdorffa
na mocy twierdzenia~\ref{thm:rzutowe-zwarte-1}; jest więc homeomorfizmem.

Model ten daje również intuicję, dlaczego $\mathbb{RP}^2$ nie jest
orientowalna. Wyobraźmy sobie małą strzałkę lub niewielką figurę
asymetryczną przesuwaną po dysku. Gdy dochodzimy do brzegu, przechodzimy
przez niego dzięki utożsamieniu z punktem leżącym po przeciwnej stronie
okręgu. Po takim przejściu kierunek wzdłuż brzegu pozostaje zgodny,
natomiast kierunek wskazujący do wnętrza dysku zostaje odwrócony.

Lokalnie można to zapisać schematycznie jako
\[
(s,t)\longmapsto (s,-t),
\]
gdzie $s$ oznacza kierunek wzdłuż brzegu, a $t$ kierunek poprzeczny.
Jedna z dwóch lokalnych współrzędnych zmienia więc zwrot. Po przejściu
wzdłuż odpowiedniej zamkniętej drogi wracamy do punktu początkowego
z odwróconą lokalną orientacją. Nie można zatem wybrać orientacji
w sposób zgodny na całej przestrzeni.

Jest to geometryczna intuicja nieorientowalności $\mathbb{RP}^2$;
formalny dowód podamy po wprowadzeniu pojęcia orientacji rozmaitości.
\end{uwaga}


\begin{figure}[h]
\centering
\begin{tikzpicture}[>=Stealth,scale=1]

  %-------------------- RP^1 --------------------
  \begin{scope}
    % Okrąg
    \draw[thick,manifoldblue] (0,0) circle (1.45);

    % Kilka par antypodalnych na okręgu
    \foreach \ang in {0,45,90,135}{
        \draw[dashed,gray!55]
            (\ang:1.45) -- ({\ang+180}:1.45);

        \fill[manifoldgold]
            (\ang:1.45) circle (2.2pt);
        \fill[manifoldgold]
            ({\ang+180}:1.45) circle (2.2pt);
    }

    % Wyróżniona para x i -x
    \fill[manifoldred] (20:1.45) circle (2.8pt);
    \fill[manifoldred] (200:1.45) circle (2.8pt);

    \node[manifoldred, above right] at (20:1.45) {$x$};
    \node[manifoldred, below left] at (200:1.45) {$-x$};

    % Strzałka pokazująca utożsamienie wyróżnionej pary
    \draw[<->,thick,manifoldred]
      (20:1.36) to[bend left=18] (200:1.36);

    % Opisy
    \node[manifoldblue] at (0,1.85) {$S^1$};

    \node[align=center,manifoldred] at (0,-1.95)
      {$x\sim -x$ dla każdego $x\in S^1$};

    \node at (0,-2.35)
      {$\displaystyle \mathbb{RP}^1\cong S^1/\{x\sim -x\}$};

    \node at (0,-2.75)
      {$\displaystyle \cong\, S^1$};
  \end{scope}

  %-------------------- RP^2 --------------------
  \begin{scope}[xshift=6.2cm]
    % Dysk
    \fill[manifoldblue!8] (0,0) circle (1.45);
    \draw[thick,manifoldblue] (0,0) circle (1.45);

    % Kilka przykładowych par punktów antypodalnych na brzegu
    \foreach \ang in {0,45,90,135}{
        \fill[manifoldgold]
            (\ang:1.45) circle (2.2pt);
        \fill[manifoldgold]
            ({\ang+180}:1.45) circle (2.2pt);
    }

    % Wyróżniona para x i -x
    \fill[manifoldred] (0:1.45) circle (2.8pt);
    \fill[manifoldred] (180:1.45) circle (2.8pt);

    \node[manifoldred, right] at (0:1.45) {$x$};
    \node[manifoldred, left]  at (180:1.45) {$-x$};

    % Strzałka pokazująca utożsamienie
    \draw[<->,thick,manifoldred]
        (1.36,0) to[bend left=20] (-1.36,0);
    \node[align=center,font=\scriptsize,manifoldblue] at (0,0.55)
        {wnętrze bez\\utożsamień};

    % Opisy
    \node[manifoldblue] at (0,1.85) {$D^2$};

    \node[align=center,manifoldred] at (0,-1.95)
        {$x\sim -x$ dla każdego $x\in\partial D^2$};

    \node at (0,-2.35)
        {$\displaystyle \mathbb{RP}^2
        \cong D^2/\{x\sim -x,\;x\in\partial D^2\}$};
  \end{scope}

\end{tikzpicture}
\caption{Dwa użyteczne modele przestrzeni rzutowych.
Po lewej, $\mathbb{RP}^1$ otrzymujemy z okręgu $S^1$ przez utożsamienie
par antypodalnych punktów. Po prawej, $\mathbb{RP}^2$ otrzymujemy z dysku
$D^2$ przez utożsamienie antypodalnych punktów jego brzegu. Różnica polega
na tym, że w pierwszym przypadku utożsamiamy punkty na całym okręgu,
natomiast w drugim utożsamienie zachodzi tylko na brzegu dysku.}
\label{fig:projective-models}
\end{figure}


\begin{twierdzenie}[Przestrzenie rzutowe jako zwarte rozmaitości]\label{thm:rzutowe-zwarte-1}
Przestrzeń $\mathbb{RP}^n$ jest zwartą rozmaitością gładką wymiaru $n$, a
$\mathbb{CP}^n$ jest zwartą rozmaitością zespoloną wymiaru $n$ (a więc zwartą
rozmaitością rzeczywistą wymiaru $2n$).
\end{twierdzenie}

\begin{proof}
Zaczniemy od zwartości. Niech
\[
\pi_{\R}:\R^{n+1}\setminus\{0\}\longrightarrow\mathbb{RP}^n,
\qquad
\pi_{\R}(x)=[x],
\]
będzie rzutowaniem ilorazowym. Ograniczając je do sfery jednostkowej,
otrzymujemy ciągłe odwzorowanie
\[
q_{\R}:S^n\longrightarrow\mathbb{RP}^n,
\qquad
q_{\R}(x)=[x].
\]
Jest ono surjektywne. Istotnie, jeśli
\[
[x]\in\mathbb{RP}^n,
\]
to $x\neq0$ i możemy znormalizować reprezentanta:
\[
\widehat{x}:=\frac{x}{\|x\|}\in S^n.
\]
Ponieważ $\widehat{x}$ jest niezerową wielokrotnością $x$, mamy
\[
[\widehat{x}]=[x].
\]
Każdy punkt $\mathbb{RP}^n$ należy więc do obrazu $q_{\R}$.

Sfera $S^n$ jest zwarta, a ciągły obraz przestrzeni zwartej jest zwarty.
Ponieważ
\[
q_{\R}(S^n)=\mathbb{RP}^n,
\]
otrzymujemy, że $\mathbb{RP}^n$ jest zwarta.

Dokładnie ten sam argument działa w przypadku zespolonym. Niech
\[
\pi_{\C}:\C^{n+1}\setminus\{0\}\longrightarrow\mathbb{CP}^n,
\qquad
\pi_{\C}(z)=[z].
\]
Po utożsamieniu $\C^{n+1}\cong\R^{2n+2}$ sfera jednostkowa ma postać
\[
S^{2n+1}
=
\{z\in\C^{n+1}:\|z\|=1\}.
\]
Ograniczenie
\[
q_{\C}:S^{2n+1}\longrightarrow\mathbb{CP}^n,
\qquad
q_{\C}(z)=[z],
\]
jest ciągłe i surjektywne, ponieważ dla każdego $[z]\in\mathbb{CP}^n$
wektor
\[
\widehat{z}:=\frac{z}{\|z\|}
\]
należy do $S^{2n+1}$ i reprezentuje tę samą klasę:
\[
[\widehat{z}]=[z].
\]
Ponieważ $S^{2n+1}$ jest zwarta, jej ciągły obraz
\[
q_{\C}(S^{2n+1})=\mathbb{CP}^n
\]
jest zwarty. Zatem również $\mathbb{CP}^n$ jest zwarta.

Sprawdźmy teraz aksjomat Hausdorffa bez powoływania się na ogólne
twierdzenie o ilorazach. Dla jednostkowych reprezentantów $x,y\in S^n$
dwóch klas rzeczywistych połóżmy
\[
d_{\R}([x],[y])
=
\min\{\|x-y\|,\|x+y\|\}.
\]
Dla jednostkowych reprezentantów $z,w\in S^{2n+1}$ połóżmy
\[
d_{\C}([z],[w])
=
\min_{|\lambda|=1}\|z-\lambda w\|.
\]
Oba wzory nie zależą od wyboru jednostkowych reprezentantów:
w przypadku rzeczywistym zmiana $x$ na $-x$ jedynie zamienia miejscami
dwie liczby występujące w minimum, natomiast w przypadku zespolonym
zmiana reprezentanta przez czynnik o module $1$ jedynie zmienia
parametryzację zbioru, po którym bierzemy minimum. Minima istnieją
dzięki zwartości grup $\{\pm1\}$ i $S^1$.

Nierówność trójkąta sprawdzimy jednocześnie dla obu przypadków,
pisząc $G=\{\pm1\}$ lub $G=S^1$. Dla jednostkowych $x,y,z$
wybierzmy $\lambda,\mu\in G$ realizujące minima dla par $(x,y)$ i $(y,z)$.
Ponieważ mnożenie przez $\lambda$ zachowuje normę,
\begin{align*}
d([x],[z])&\leq\|x-\lambda\mu z\|\\
&\leq\|x-\lambda y\|+\|\lambda y-\lambda\mu z\|\\
&=d([x],[y])+d([y],[z]).
\end{align*}
Symetria wynika z zastąpienia $\lambda$ przez $\lambda^{-1}$.
Ponadto osiągane minimum jest równe zeru wtedy i tylko wtedy,
gdy reprezentanty należą do tej samej orbity, czyli reprezentują
ten sam punkt przestrzeni rzutowej. Otrzymujemy więc metryki na zbiorach klas.

Pozostaje sprawdzić, że topologie wyznaczone przez te metryki są
topologiami ilorazowymi. Rzutowania
\[
q_{\R}:S^n\to(\mathbb{RP}^n,d_{\R}),
\qquad
q_{\C}:S^{2n+1}\to(\mathbb{CP}^n,d_{\C})
\]
są ciągłe, ponieważ odległość klas nie przekracza odległości ich
wybranych reprezentantów. Aby przejść od sfer do pierwotnych ilorazów,
używamy ciągłej normalizacji $\nu(v)=v/\|v\|$.
Rzutowania $\pi_{\R}$ i $\pi_{\C}$ do przestrzeni klas z metryką są ciągłe,
bo mają postać $q_{\R}\circ\nu$ i $q_{\C}\circ\nu$.
Z definicji topologii ilorazowej wynika więc ciągłość identyczności
z pierwotnego ilorazu do przestrzeni metrycznej.
Z drugiej strony mamy już wykazaną zwartość
$\mathbb{RP}^n$ i $\mathbb{CP}^n$ w ich topologiach ilorazowych.
Identyczność z przestrzeni rzutowej wyposażonej w topologię ilorazową
do tego samego zbioru klas wyposażonego w powyższą metrykę jest więc
ciągłą bijekcją ze zwartej przestrzeni do przestrzeni Hausdorffa.
Jest zatem homeomorfizmem. W szczególności pierwotna topologia
ilorazowa jest Hausdorffa.

Każda z przestrzeni rzutowych ma skończony atlas
\[
U_0,\ldots,U_n,
\]
którego obrazy są równe odpowiednio $\R^n$ lub $\C^n$. Wybierzmy
w $\R^n$ oraz $\C^n\cong\R^{2n}$ przeliczalną bazę, na przykład
złożoną z kul o wymiernych środkach i wymiernych promieniach.
Przeciwobrazy tych baz przez skończenie wiele map atlasu tworzą
przeliczalną bazę przestrzeni rzutowej.

Odwzorowania przejścia są ilorazami współrzędnych przez współrzędną
niezerową, a więc są gładkie w przypadku rzeczywistym i holomorficzne
w przypadku zespolonym. Otrzymujemy zatem strukturę rozmaitości
gładkiej wymiaru $n$ na $\mathbb{RP}^n$ oraz strukturę rozmaitości
zespolonej wymiaru $n$ na $\mathbb{CP}^n$. Ponieważ
\[
\C^n\cong\R^{2n},
\]
$\mathbb{CP}^n$ ma jako rozmaitość rzeczywista wymiar $2n$.
\end{proof}


\begin{uwaga}
Zwartość można zobaczyć bardzo bezpośrednio: każde $\mathbb{RP}^n$ jest ilorazem
sfery $S^n$ przez relację $x\sim-x$, a $\mathbb{CP}^n$ jest ilorazem sfery
jednostkowej $S^{2n+1}\subset\C^{n+1}$ przez działanie grupy $S^1$,
\[
z\sim e^{i\theta}z.
\]
W przypadku zespolonym otrzymujemy w ten sposób słynną fibrację Hopfa
\[
S^1\longrightarrow S^{2n+1}\longrightarrow\mathbb{CP}^n.
\]
Pełny opis tej konstrukcji wykracza poza zakres tego rozdziału, ale pokazuje,
że przestrzenie rzutowe są naturalnie związane zarówno z ilorazami, jak i z teorią
wiązek.
\end{uwaga}

\begin{przyklad}[Torus $T^n$]\label{prz:torus}
Niech $T^n=\R^n/\Z^n$ z topologią ilorazową względem rzutowania kanonicznego $\pi\colon\R^n
\to T^n$, $\pi(x)=x+\Z^n$. Pokażemy najpierw, że $\pi$ jest odwzorowaniem otwartym: dla
zbioru otwartego $\Omega\subseteq\R^n$, $\pi^{-1}(\pi(\Omega))=\bigcup_{k\in\Z^n}(\Omega+k)$
jest sumą zbiorów otwartych, a więc otwarty; z definicji topologii ilorazowej oznacza to,
że $\pi(\Omega)$ jest otwarty w $T^n$.

Sprawdźmy pozostałe warunki topologiczne. Jeśli $[x]\ne[y]$, to
$\delta=\inf_{k\in\Z^n}\|x-y-k\|>0$: w dowolnym ograniczonym
zbiorze jest tylko skończenie wiele punktów kraty, a $x-y\notin\Z^n$.
Otwarte otoczenia $\pi(B(x,\delta/3))$ i $\pi(B(y,\delta/3))$
są rozłączne. W przeciwnym razie istniałyby $u,v$ w tych kulach
z $u-v=k\in\Z^n$, skąd $\|x-y-k\|<2\delta/3$, sprzeczność.
Torus jest więc Hausdorffa. Obrazy przeliczalnej bazy kul w $\R^n$
stanowią bazę ilorazu: dla otwartego $O\subset T^n$ i $[x]\in O$
wybieramy kulę bazową $B$ z $x\in B\subset\pi^{-1}(O)$;
wtedy $[x]\in\pi(B)\subset O$. To daje drugi aksjomat przeliczalności.

Niech $B(x,1/2)\subseteq\R^n$ będzie kulą otwartą o promieniu $1/2$. Dla $k\in\Z^n\setminus
\{0\}$, zbiory $B(x,1/2)$ i $B(x+k,1/2)$ są rozłączne: gdyby $y$ należał do obu, to
$\|y-x\|<1/2$ i $\|y-x-k\|<1/2$, skąd z nierówności trójkąta $\|k\|\leq\|y-x\|+\|y-x-k\|<1$;
ale $\|k\|\geq1$ dla $k\in\Z^n\setminus\{0\}$ — sprzeczność. Zatem $\pi|_{B(x,1/2)}$ jest
różnowartościowe, a więc — jako odwzorowanie ciągłe, otwarte i różnowartościowe na obraz —
homeomorfizmem $B(x,1/2)\to\pi(B(x,1/2))$. Kładziemy $\varphi_x:=(\pi|_{B(x,1/2)})^{-1}$;
zbiory $\pi(B(x,1/2))$ dla $x\in\R^n$ pokrywają $T^n$.

Odwzorowania przejścia między takimi mapami są (lokalnie) postaci $y\mapsto y+k$ dla
pewnego ustalonego $k\in\Z^n$ (bo dwa punkty z tym samym obrazem przy $\pi$ różnią się
dokładnie o wektor z $\Z^n$, a $k$ jest lokalnie stałe z ciągłości) — a translacja jest
odwzorowaniem gładkim z gładką odwrotną. Zatem otrzymana rodzina map jest atlasem, i $T^n$
jest rozmaitością gładką wymiaru $n$.
\end{przyklad}
\section{Przestrzeń styczna}\label{sec:przestrzen-styczna}

W porównaniu trzech poniższych modeli zakładamy,
że $M$ nie ma brzegu. Dla punktów brzegu
przestrzeń styczną definiujemy przez derywacje
lub współrzędne przedłużane poza półprzestrzeń,
jak w Sekcji~\ref{sec:rozmaitosci-z-brzegiem-1};
model dwustronnych krzywych nie działa tam dosłownie.

Chcemy nadać sens pojęciu ,,wektora stycznego'' do rozmaitości $M$ w punkcie $p$ —
uogólnieniu intuicji wektora stycznego do powierzchni w $\R^3$, niezależnemu od
zanurzenia $M$ w jakąkolwiek przestrzeń otaczającą. Istnieje kilka naturalnych, a priori
różnych, sposobów sformalizowania tej intuicji; w tej sekcji podajemy trzy z nich i
dowodzimy, że są one kanonicznie równoważne.

\subsection*{Trzy definicje przestrzeni stycznej}

Niech $M$ będzie rozmaitością gładką wymiaru $n$, $p\in M$.

\textbf{Definicja A (poprzez krzywe).} Niech $\Gamma_p$ oznacza zbiór gładkich krzywych
$\gamma\colon(-\varepsilon,\varepsilon)\to M$ (dla pewnego $\varepsilon>0$, zależnego od
$\gamma$) z $\gamma(0)=p$. Mówimy, że $\gamma_1,\gamma_2\in\Gamma_p$ są \emph{styczne w
$p$}, $\gamma_1\sim\gamma_2$, jeśli dla pewnej (równoważnie: każdej) mapy dopuszczalnej
$(U,\varphi)$ wokół $p$,
\[
(\varphi\circ\gamma_1)'(0)=(\varphi\circ\gamma_2)'(0)\quad\text{w }\R^n.
\]

\begin{figure}[!ht]
\centering
\begin{tikzpicture}[>=Stealth,scale=1.02,line cap=round,line join=round]
% Dwie różne krzywe, których pochodne w p po przejściu do współrzędnych są równe.
\filldraw[fill=manifoldblue!9,draw=manifoldblue!75!black,very thick]
 (-5.04,-.52) .. controls (-5.0,.46) and (-4.05,1.21) .. (-2.67,1.24)
 .. controls (-1.44,1.28) and (-.63,.5) .. (-.62,-.27)
 .. controls (-.6,-1.1) and (-2.13,-1.42) .. (-3.47,-1.18)
 .. controls (-4.39,-1.02) and (-5.07,-.96) .. cycle;
\filldraw[fill=manifoldgreen!12,draw=manifoldgreen!65!black,thick]
 (-4.38,-.36) .. controls (-4.4,.4) and (-3.61,.94) .. (-2.62,.86)
 .. controls (-1.65,.79) and (-1.1,.21) .. (-1.21,-.44)
 .. controls (-1.34,-1.01) and (-2.37,-1.02) .. (-3.13,-.89)
 .. controls (-3.91,-.83) and (-4.34,-.8) .. cycle;
\node[font=\small\bfseries,manifoldblue!75!black] at (-4.68,1.37) {$M$};
\node[font=\small,manifoldgreen!65!black] at (-3.45,.55) {$U$};
\draw[manifoldred,very thick]
 (-4.08,-.68) .. controls (-3.7,-.42) and (-3.18,-.13) .. (-2.8,-.03)
 .. controls (-2.42,.07) and (-1.98,.26) .. (-1.72,.29);
\draw[manifoldgold!85!black,very thick]
 (-4.06,-.02) .. controls (-3.65,-.46) and (-3.18,-.13) .. (-2.8,-.03)
 .. controls (-2.42,.07) and (-1.86,-.15) .. (-1.52,-.55);
\node[font=\small,manifoldred] at (-4.31,-.83) {$\gamma_1$};
\node[font=\small,manifoldgold!80!black] at (-4.35,.12) {$\gamma_2$};
\fill[black!75] (-2.8,-.03) circle (2.6pt)
 node[below,font=\small,xshift=-2pt] {$p$};
\draw[->,very thick,manifoldblue!75!black] (-.33,-.08)--(1.19,-.08)
 node[midway,above,font=\small] {$\varphi$};
% Rysunkowy fragment R^n oraz obraz U.
\filldraw[fill=manifoldblue!3,draw=manifoldblue!65!black,thick]
 (1.42,-1.17)--(4.94,-1.17)--(5.43,1.08)--(1.91,1.08)--cycle;
\filldraw[fill=manifoldgreen!10,draw=manifoldgreen!65!black,thick]
 (1.87,-.41) .. controls (1.93,.27) and (2.64,.77) .. (3.51,.72)
 .. controls (4.35,.66) and (4.8,.14) .. (4.66,-.5)
 .. controls (4.52,-.9) and (3.68,-1.02) .. (2.95,-.85)
 .. controls (2.32,-.78) and (1.88,-.73) .. cycle;
\draw[manifoldred,very thick]
 (2.11,-.62) .. controls (2.39,-.4) and (2.69,-.14) .. (3.03,-.05)
 .. controls (3.37,.04) and (3.91,.21) .. (4.17,.27);
\draw[manifoldgold!85!black,very thick]
 (2.25,.03) .. controls (2.48,-.29) and (2.69,-.14) .. (3.03,-.05)
 .. controls (3.37,.04) and (3.91,-.23) .. (4.21,-.5);
\fill[black!75] (3.03,-.05) circle (2.6pt)
 node[below,font=\small,xshift=-1pt] {$\varphi(p)$};
\draw[->,very thick,manifoldblue!85!black] (3.03,-.05)--(3.85,.17);
\node[font=\small,manifoldblue!85!black] at (4.05,.34) {$v$};
\node[font=\small,manifoldgreen!65!black] at (4.23,.82) {$\varphi(U)$};
\node[font=\small\bfseries,manifoldblue!75!black] at (5.23,1.35) {$\R^n$};
\node[font=\small] at (.17,-1.73)
 {$v=(\varphi\circ\gamma_1)'(0)=(\varphi\circ\gamma_2)'(0)$};
\end{tikzpicture}
\caption{Dwie krzywe w $M$ przechodzą przez $p$ i mają ten sam kierunek
styczny. W mapie $\varphi$ ich obrazy mają wspólną pochodną $v$ w $t=0$.}
\label{fig:krzywe-styczne-w-mapie}
\end{figure}

\begin{fakt}
Powyższy warunek nie zależy od wyboru mapy $(U,\varphi)$.
\end{fakt}
\begin{proof}
Jeśli $(V,\psi)$ jest inną mapą wokół $p$, to na wspólnej dziedzinie $\psi\circ\gamma_i=
(\psi\circ\varphi^{-1})\circ(\varphi\circ\gamma_i)$ blisko $t=0$, więc z reguły łańcucha
\[
(\psi\circ\gamma_i)'(0)=D(\psi\circ\varphi^{-1})(\varphi(p))\cdot(\varphi\circ\gamma_i)'(0),
\qquad i=1,2.
\]
Ponieważ $D(\psi\circ\varphi^{-1})(\varphi(p))$ jest tym samym przekształceniem liniowym
dla $i=1$ i $i=2$, równość $(\varphi\circ\gamma_1)'(0)=(\varphi\circ\gamma_2)'(0)$
implikuje $(\psi\circ\gamma_1)'(0)=(\psi\circ\gamma_2)'(0)$, i na odwrót (bo
$D(\psi\circ\varphi^{-1})(\varphi(p))$ jest odwracalne, jako różniczka dyfeomorfizmu).
\end{proof}

Relacja $\sim$ jest relacją równoważności na $\Gamma_p$ (natychmiast z definicji, poprzez
odpowiednie własności równości w $\R^n$). Definiujemy $T_p^AM:=\Gamma_p/\!\sim$.

Aby wyposażyć $T_p^AM$ w strukturę przestrzeni wektorowej, ustalmy mapę
$(U,\varphi)$ wokół $p$ i zdefiniujmy
\[
\Theta_\varphi:T_p^AM\longrightarrow\R^n,
\qquad
\Theta_\varphi([\gamma]):=(\varphi\circ\gamma)'(0).
\]
Odwzorowanie to jest dobrze określone i bijektywne. Różnowartościowość
wynika bezpośrednio z definicji relacji równoważności, natomiast
surjektywność z faktu, że dla dowolnego $v\in\R^n$ krzywa
\[
\gamma(t)=\varphi^{-1}(\varphi(p)+tv),
\]
określona dla dostatecznie małych $|t|$, spełnia
\[
\gamma(0)=p,
\qquad
(\varphi\circ\gamma)'(0)=v.
\]

Za pomocą bijekcji $\Theta_\varphi$ przenosimy strukturę przestrzeni
wektorowej z $\R^n$ na $T_p^AM$, definiując
\[
[\gamma]+[\eta]
:=
\Theta_\varphi^{-1}
\bigl(\Theta_\varphi([\gamma])+\Theta_\varphi([\eta])\bigr)
\]
oraz
\[
a[\gamma]
:=
\Theta_\varphi^{-1}
\bigl(a\,\Theta_\varphi([\gamma])\bigr),
\qquad a\in\R.
\]
W szczególności mnożenie przez skalar można reprezentować przez
reparametryzację krzywej:
\[
a[\gamma]=[t\mapsto\gamma(at)],
\]
ponieważ
\[
(\varphi\circ\gamma(at))'_{\,t=0}
=
a(\varphi\circ\gamma)'(0).
\]
W ten sposób $\Theta_\varphi$ jest, z definicji, izomorfizmem liniowym
$T_p^AM\cong\R^n$.

\begin{fakt}
Struktura liniowa na $T_p^AM$ nie zależy od wyboru mapy $\varphi$.
\end{fakt}
\begin{proof}
Niech $(V,\psi)$ będzie inną mapą wokół $p$, $L:=D(\psi\circ\varphi^{-1})(\varphi(p))
\colon\R^n\to\R^n$. Z rachunku w poprzednim Fakcie, dla dowolnej krzywej $\gamma\in\Gamma_p$,
$\Theta_\psi([\gamma])=L(\Theta_\varphi([\gamma]))$, czyli $\Theta_\psi=L\circ\Theta_\varphi$.
Ponieważ $L$ jest izomorfizmem \emph{liniowym} (jako różniczka dyfeomorfizmu w punkcie),
struktury liniowe przeniesione na $T_p^AM$ przez $\Theta_\varphi$ i przez $\Theta_\psi=
L\circ\Theta_\varphi$ pokrywają się: liniowość $L$ gwarantuje, że dla dowolnych $[\gamma_1],
[\gamma_2]\in T_p^AM$, $a,b\in\R$, kombinacja $a[\gamma_1]+b[\gamma_2]$ wyznaczona przez
$\Theta_\varphi$ pokrywa się z kombinacją wyznaczoną przez $\Theta_\psi$. To kluczowa
obserwacja: liniowość różniczki odwzorowania przejścia gwarantuje, że mimo iż
$\Theta_\varphi\neq\Theta_\psi$, wyznaczona przez nie struktura wektorowa na $T_p^AM$ jest
jedna i ta sama.
\end{proof}

\textbf{Definicja B (klasyczna, poprzez pary mapa-wektor).} Rozważmy zbiór trójek
$(U,\varphi,v)$, gdzie $(U,\varphi)$ jest mapą dopuszczalną wokół $p$, a $v\in\R^n$.
Mówimy, że $(U,\varphi,v)\sim(V,\psi,w)$, jeśli $w=D(\psi\circ\varphi^{-1})(\varphi(p))\,v$.

\begin{fakt}
Powyższa relacja jest relacją równoważności.
\end{fakt}
\begin{proof}
Zwrotność: dla $\psi=\varphi$, $D(\id)=\id$, więc $v=v$. Symetria: jeśli $w=D(\psi\varphi^{-1})
(\varphi(p))v$, to stosując regułę łańcucha do $\varphi\circ\psi^{-1}=(\psi\circ
\varphi^{-1})^{-1}$ w punkcie $\psi(p)$, $D(\varphi\psi^{-1})(\psi(p))=\big[D(\psi\varphi^{-1})
(\varphi(p))\big]^{-1}$, skąd $v=D(\varphi\psi^{-1})(\psi(p))w$. Przechodniość: jeśli
$w=D(\psi\varphi^{-1})(\varphi(p))v$ i $u=D(\eta\psi^{-1})(\psi(p))w$, to z reguły łańcucha
zastosowanej do $\eta\circ\varphi^{-1}=(\eta\circ\psi^{-1})\circ(\psi\circ\varphi^{-1})$ w
punkcie $\varphi(p)$,
\[
D(\eta\varphi^{-1})(\varphi(p))=D(\eta\psi^{-1})(\psi(p))\cdot D(\psi\varphi^{-1})(\varphi(p)),
\]
a więc $D(\eta\varphi^{-1})(\varphi(p))v=D(\eta\psi^{-1})(\psi(p))w=u$, czyli
$(U,\varphi,v)\sim(W,\eta,u)$.
\end{proof}

Definiujemy $T_p^BM$ jako zbiór klas równoważności
\[
[U,\varphi,v].
\]
Strukturę przestrzeni wektorowej wprowadzamy wzorem
\[
a[U,\varphi,v]+b[U,\varphi,v']
:=
[U,\varphi,av+bv'].
\]
Definicja ta jest dobrze określona, to znaczy wynik nie zależy od wyboru
reprezentantów klas. Istotnie, jeżeli przejdziemy do innej mapy
$(V,\psi)$ zawierającej punkt $p$, to wektory $v,v'\in\mathbb{R}^n$
przechodzą odpowiednio w
\[
D(\psi\circ\varphi^{-1})(\varphi(p))\,v
\qquad\text{oraz}\qquad
D(\psi\circ\varphi^{-1})(\varphi(p))\,v'.
\]
Z liniowości różniczki otrzymujemy
\[
D(\psi\circ\varphi^{-1})(\varphi(p))(av+bv')
=
a\,D(\psi\circ\varphi^{-1})(\varphi(p))v
+
b\,D(\psi\circ\varphi^{-1})(\varphi(p))v',
\]
a więc oba wybory reprezentantów prowadzą do tej samej klasy
równoważności.

Ustalmy teraz mapę $(U,\varphi)$ wokół punktu $p$. Definiujemy
odwzorowanie
\[
\Phi_\varphi\colon T_p^AM\longrightarrow T_p^BM,
\qquad
\Phi_\varphi([\gamma])
=
[U,\varphi,(\varphi\circ\gamma)'(0)].
\]
Odwzorowanie to jest dobrze określone, ponieważ jeżeli
$[\gamma]=[\eta]$ w $T_p^AM$, to z definicji relacji równoważności
krzywych mamy
\[
(\varphi\circ\gamma)'(0)
=
(\varphi\circ\eta)'(0),
\]
a zatem
\[
[U,\varphi,(\varphi\circ\gamma)'(0)]
=
[U,\varphi,(\varphi\circ\eta)'(0)].
\]

Żeby sprawdzić bijektywność, rozważmy odwzorowanie
\[
j_\varphi\colon\R^n\longrightarrow T_p^BM,
\qquad j_\varphi(v)=[U,\varphi,v].
\]
Jest ono liniowe na mocy definicji działań. Jest też
różnowartościowe: jeżeli $[U,\varphi,v]=[U,\varphi,w]$,
to relacja równoważności w tej samej mapie daje $v=w$.
Jest na, gdyż dla dowolnej klasy $[V,\psi,w]$ element
\[
v=D(\varphi\circ\psi^{-1})(\psi(p))w
\]
spełnia $[V,\psi,w]=[U,\varphi,v]$. Ponieważ
$\Theta_\varphi$ jest liniową bijekcją, mamy
$\Phi_\varphi=j_\varphi\circ\Theta_\varphi$,
a więc $\Phi_\varphi$ również jest liniową bijekcją.
Zmiana mapy nie zmienia tego odwzorowania: dla $(V,\psi)$
reguła łańcucha i definicja relacji równoważności dają
$j_\psi\circ\Theta_\psi=j_\varphi\circ\Theta_\varphi$.
Otrzymujemy zatem kanoniczne utożsamienie
\[
T_p^AM\cong T_p^BM.
\]
Od tej chwili obie te przestrzenie będziemy oznaczać wspólnym symbolem
\[
T_p^{\mathrm{geom}}M.
\]
\textbf{Definicja C (algebraiczna, poprzez derywacje).} Niech $C^\infty(M)$ oznacza
przestrzeń funkcji gładkich $M\to\R$. \emph{Derywacją w punkcie $p$} nazywamy funkcjonał
liniowy $v\colon C^\infty(M)\to\R$ spełniający regułę Leibniza:
\[
v(fg)=f(p)\,v(g)+g(p)\,v(f), \qquad f,g\in C^\infty(M).
\]
Zbiór wszystkich derywacji w $p$ oznaczamy $T_p^CM$ — jest to podprzestrzeń liniowa
przestrzeni funkcjonałów liniowych na $C^\infty(M)$ (suma i wielokrotność derywacji
spełnia regułę Leibniza — sprawdzenie bezpośrednie).

W dalszej części skonstruujemy kanoniczny izomorfizm
$T_p^{\mathrm{geom}}M\to T_p^CM$. Utożsamimy w ten sposób
wszystkie trzy definicje i otrzymamy bazę przestrzeni
stycznej złożoną z pochodnych cząstkowych.

\subsection*{Lemat Hadamarda}

Poniższy lemat analityczny jest podstawowym narzędziem technicznym tej sekcji.

\begin{lemat}[Lemat Hadamarda]\label{lem:hadamard}
Niech $\Omega\subseteq\R^n$ będzie zbiorem otwartym, gwiaździstym względem $0$ (tzn.\
$tx\in\Omega$ dla wszystkich $x\in\Omega$, $t\in[0,1]$), i niech $f\in C^\infty(\Omega)$.
Wówczas istnieją funkcje $g_1,\dots,g_n\in C^\infty(\Omega)$ takie, że
\[
f(x)=f(0)+\sum_{i=1}^nx^ig_i(x), \qquad x\in\Omega,
\]
oraz $g_i(0)=\dfrac{\partial f}{\partial x^i}(0)$ dla każdego $i$.
\end{lemat}

\begin{proof}
Dla ustalonego $x\in\Omega$, funkcja $t\mapsto f(tx)$ jest dobrze określona i gładka na
$[0,1]$ (bo $\Omega$ jest gwiaździsty względem $0$), więc z twierdzenia podstawowego
rachunku różniczkowego,
\[
f(x)-f(0)=\int_0^1\frac{d}{dt}f(tx)\,dt=\int_0^1\sum_{i=1}^n\frac{\partial f}{\partial x^i}
(tx)\,x^i\,dt=\sum_{i=1}^nx^i\underbrace{\int_0^1\frac{\partial f}{\partial x^i}(tx)\,dt}
_{=:g_i(x)}.
\]
Funkcje $g_i(x):=\int_0^1\frac{\partial f}{\partial x^i}(tx)\,dt$ są gładkie na $\Omega$
(różniczkowanie pod znakiem całki, dopuszczalne, gdyż funkcja podcałkowa zależy gładko od
$(t,x)$), oraz $g_i(0)=\int_0^1\frac{\partial f}{\partial x^i}(0)\,dt=\frac{\partial f}
{\partial x^i}(0)$, co kończy dowód.
\end{proof}

\subsection*{Od krzywych do derywacji}

Zdefiniujmy odwzorowanie $\Phi\colon T_p^AM\to T_p^CM$ wzorem
\[
\Phi([\gamma])(f):=(f\circ\gamma)'(0), \qquad f\in C^\infty(M).
\]

\begin{fakt}
$\Phi$ jest dobrze określone (niezależne od wyboru reprezentanta $\gamma$), a
$\Phi([\gamma])$ jest derywacją w $p$ dla każdego $[\gamma]\in T_p^AM$.
\end{fakt}
\begin{proof}
Niezależność od reprezentanta: jeśli $\gamma_1\sim\gamma_2$, to pisząc $f\circ\gamma_i=
(f\circ\varphi^{-1})\circ(\varphi\circ\gamma_i)$ dla mapy $(U,\varphi)$ wokół $p$,
\[
(f\circ\gamma_i)'(0)=D(f\circ\varphi^{-1})(\varphi(p))\cdot(\varphi\circ\gamma_i)'(0),
\]
a prawa strona zależy od $\gamma_i$ jedynie poprzez $(\varphi\circ\gamma_i)'(0)$, które z
założenia $\gamma_1\sim\gamma_2$ jest tym samym wektorem dla $i=1,2$. Liniowość
$\Phi([\gamma])$ względem $f$ wynika z liniowości pochodnej $(\cdot)'(0)$ oraz liniowości
$D(f\circ\varphi^{-1})$ względem $f$. Reguła Leibniza: $\big((fg)\circ\gamma\big)'(0)=
\big((f\circ\gamma)(g\circ\gamma)\big)'(0)=(f\circ\gamma)'(0)\,g(p)+f(p)\,(g\circ\gamma)'(0)
=\Phi([\gamma])(f)\,g(p)+f(p)\,\Phi([\gamma])(g)$, na mocy zwykłej reguły Leibniza dla
pochodnej iloczynu funkcji jednej zmiennej $t\mapsto f(\gamma(t))g(\gamma(t))$.
\end{proof}

Odwzorowanie
\[
\Phi\colon T_p^A M\to T_p^C M
\]
jest liniowe względem struktury wektorowej na $T_p^A M$ zdefiniowanej
za pomocą izomorfizmu
\[
\Theta_\varphi\colon T_p^A M\to\mathbb R^n.
\]
Istotnie, jeśli
\[
v=\Theta_\varphi([\gamma]),
\qquad
w=\Theta_\varphi([\eta]),
\]
to dla $a,b\in\mathbb R$ mamy
\[
\Theta_\varphi(a[\gamma]+b[\eta])=av+bw.
\]
Dla dowolnego $f\in C^\infty(M)$ otrzymujemy zatem
\begin{align*}
\Phi(a[\gamma]+b[\eta])(f)
&=
D(f\circ\varphi^{-1})(\varphi(p))(av+bw)\\
&=
a\,D(f\circ\varphi^{-1})(\varphi(p))(v)
+
b\,D(f\circ\varphi^{-1})(\varphi(p))(w)\\
&=
a\,\Phi([\gamma])(f)+b\,\Phi([\eta])(f).
\end{align*}
Stąd $\Phi$ jest odwzorowaniem liniowym.

Pokażemy teraz, że $\Phi$ jest izomorfizmem.
Kluczowym krokiem będzie zastosowanie Lematu Hadamarda, które
wyjaśni również, dlaczego pochodne cząstkowe
\[
\left.\frac{\partial}{\partial x^1}\right|_p,\ldots,
\left.\frac{\partial}{\partial x^n}\right|_p
\]
tworzą bazę $T_p^C M$.

Będziemy używać gładkiej \emph{funkcji odcinającej}
$\mu\colon\R^n\to[0,1]$ (ang.\ \emph{bump function}),
równej $1$ na $\overline{B(0,1)}$ i $0$ poza $B(0,2)$.
Zbudujmy ją jawnie. Niech $h(t)=e^{-1/t}$ dla $t>0$,
a $h(t)=0$ dla $t\leq0$. Wszystkie pochodne $h$ po
prawej stronie zera są sumami wyrazów postaci
$t^{-m}e^{-1/t}$, które dążą do zera przy $t\downarrow0$;
po lewej stronie wszystkie pochodne wynoszą zero.
Funkcja $h$ jest więc gładka także w zerze. Kładziemy
\[
 \theta(t):=\frac{h(t)}{h(t)+h(1-t)},\qquad
 \mu(x):=\theta\!\left(\frac{4-\|x\|^2}{3}\right).
\]
Mianownik $\theta$ nigdzie nie znika: dla $t\leq0$
dodatni jest $h(1-t)$, dla $t\geq1$ dodatni jest $h(t)$,
a pomiędzy nimi dodatnie są oba. Mamy $\theta(t)=0$
dla $t\leq0$ i $\theta(t)=1$ dla $t\geq1$.
Jeśli $\|x\|\leq1$, argument $\theta$ jest co najmniej $1$;
jeśli $\|x\|\geq2$, jest nie większy od zera.
Cała konstrukcja jest gładka, bo $\|x\|^2$ jest
wielomianem współrzędnych.


% ============================================================
% RYSUNEK 1
% U \subset W \subset M oraz mapa \varphi
% ============================================================

\begin{figure}[ht]
\centering

\begin{tikzpicture}[
    >=Latex,
    font=\small,
    line cap=round,
    line join=round,
    scale=0.78
]

% M
\draw[
    thick,
    fill=blue!5
]
    (-5.4,0.1)
    .. controls (-5.1,1.9) and (-3.6,2.55) .. (-1.95,2.35)
    .. controls (-0.55,2.2) and (0.10,1.15) .. (-0.25,-0.10)
    .. controls (-0.55,-1.30) and (-2.0,-1.80) .. (-3.65,-1.55)
    .. controls (-4.8,-1.35) and (-5.55,-0.60) .. (-5.4,0.1);

\node[anchor=west] at (-5.55,1.75)
    {$M$};

% W
\draw[
    thick,
    fill=green!8
]
    (-4.45,0.15)
    .. controls (-4.25,1.20) and (-3.40,1.60) .. (-2.35,1.43)
    .. controls (-1.35,1.28) and (-0.95,0.62) .. (-1.20,-0.18)
    .. controls (-1.50,-0.92) and (-2.50,-1.07) .. (-3.50,-0.85)
    .. controls (-4.20,-0.70) and (-4.55,-0.25) .. (-4.45,0.15);

\node[anchor=west] at (-4.70,1.08)
    {$W$};

% U
\draw[
    thick,
    fill=yellow!18
]
    (-3.76,0.10)
    .. controls (-3.62,0.72) and (-3.05,0.96) .. (-2.42,0.83)
    .. controls (-1.88,0.71) and (-1.64,0.28) .. (-1.84,-0.20)
    .. controls (-2.05,-0.65) and (-2.70,-0.76) .. (-3.25,-0.55)
    .. controls (-3.65,-0.40) and (-3.84,-0.10) .. (-3.76,0.10);

\node[anchor=west] at (-4.05,0.67)
    {$U$};

% Punkt p
\fill (-2.63,0.04) circle (1.7pt);

\node[below=5pt] at (-2.63,0.15)
    {$p$};

% Strzałka \varphi
\draw[->,thick]
    (0.20,0.45)
    .. controls (0.90,1.00) and (1.70,1.00) ..
    (2.40,0.45);

\node[above=5pt] at (1.30,0.95)
    {$\varphi$};

% Obraz mapy
\draw[
    thick,
    fill=blue!4
]
    (3.0,-1.2)
    --
    (4.1,1.6)
    --
    (7.8,1.6)
    --
    (6.7,-1.2)
    --
    cycle;

\node[above=7pt] at (7.90,1.80)
    {$\varphi(U)\subset\mathbb{R}^n$};

% \overline{B(0,2)}
\draw[
    thick,
    fill=green!10
]
    (5.50,0.12)
    ellipse (1.40 and 0.78);

\node[above=8pt] at (5.90,1.10)
    {$\overline{B(0,2)}$};

% B(0,1)
\draw[
    thick,
    fill=yellow!25
]
    (5.50,0.12)
    ellipse (0.63 and 0.35);

\node[below=7pt] at (5.50,-0.23)
    {$B(0,1)$};

% 0 = \varphi(p)
\fill (5.50,0.12) circle (1.7pt);

\draw[thin,gray] (5.56,0.06) -- (6.45,-0.65);
\node[anchor=west] at (6.50,-0.70)
    {$0=\varphi(p)$};

\end{tikzpicture}

\caption{
Wybór mapy $(U,\varphi)$ wokół punktu $p$:
$U\subseteq W$, $\varphi(p)=0$ oraz
$\overline{B(0,2)}\subseteq\varphi(U)$.
}
\label{fig:kapelusz-mapa}

\end{figure}


% ============================================================
% RYSUNEK 2
% Wykres funkcji kapelusz
% ============================================================

\begin{figure}[ht]
\centering

\begin{tikzpicture}[
    >=Latex,
    font=\small,
    x=1.45cm,
    y=1.75cm
]

% Delikatne kolorowe obszary
\fill[blue!5]
    (-3,-0.05)
    rectangle
    (-2,1.18);

\fill[green!8]
    (-2,-0.05)
    rectangle
    (-1,1.18);

\fill[yellow!18]
    (-1,-0.05)
    rectangle
    (1,1.18);

\fill[green!8]
    (1,-0.05)
    rectangle
    (2,1.18);

\fill[blue!5]
    (2,-0.05)
    rectangle
    (3,1.18);

% Oś pozioma
\draw[->,thick]
    (-3.15,0)
    --
    (3.25,0)
    node[right] {$t$};

% Oś pionowa
\draw[->,thick]
    (0,-0.13)
    --
    (0,1.35)
    node[above] {$\mu(t)$};

% Wykres funkcji kapelusz
\draw[very thick]
    (-3,0)
    --
    (-2,0)
    .. controls (-1.75,0) and (-1.35,1) ..
    (-1,1)
    --
    (1,1)
    .. controls (1.35,1) and (1.75,0) ..
    (2,0)
    --
    (3,0);

% Linie pomocnicze
\draw[dashed,gray]
    (-1,0) -- (-1,1);

\draw[dashed,gray]
    (1,0) -- (1,1);

\draw[dashed,gray]
    (-2,0) -- (-2,0.36);

\draw[dashed,gray]
    (2,0) -- (2,0.36);

% Znaczniki na osi
\draw
    (-2,0.04) -- (-2,-0.04);
\node[below=5pt] at (-2,0)
    {$-2$};

\draw
    (-1,0.04) -- (-1,-0.04);
\node[below=5pt] at (-1,0)
    {$-1$};

\draw
    (1,0.04) -- (1,-0.04);
\node[below=5pt] at (1,0)
    {$1$};

\draw
    (2,0.04) -- (2,-0.04);
\node[below=5pt] at (2,0)
    {$2$};

% Poziom 1
\draw
    (-0.04,1) -- (0.04,1);

\node[left=6pt] at (0,1)
    {$1$};

% Opisy — odsunięte od wykresu
\node[above=4pt] at (0.72,1)
    {$\mu=1$};

\node at (-1.52,1.08)
    {$0<\mu<1$};

\node at (1.52,1.08)
    {$0<\mu<1$};

\node at (-2.52,0.20)
    {$\mu=0$};

\node at (2.52,0.20)
    {$\mu=0$};

\end{tikzpicture}

\caption{
Jednowymiarowy przekrój funkcji odcinającej $\mu$.
Funkcja jest równa $1$ dla $|t|\leq1$,
płynnie maleje dla $1<|t|<2$
i znika dla $|t|\geq2$.
}
\label{fig:kapelusz-wykres}

\end{figure}




\begin{lemat}[Lokalność derywacji]\label{lem:lokalnosc}
Niech $v\in T_p^CM$ będzie derywacją w $p$. Jeśli $h\in C^\infty(M)$ znika na pewnym
otoczeniu punktu $p$, to $v(h)=0$. W szczególności, jeśli $f,g\in C^\infty(M)$ pokrywają
się na pewnym otoczeniu $p$, to $v(f)=v(g)$; wartość $v(f)$ zależy więc jedynie od
zachowania $f$ w dowolnie małym otoczeniu $p$.
\end{lemat}
\begin{proof}
Z reguły Leibniza zastosowanej do $f=g=\mathbf1$ (funkcji stałej równej $1$),
$v(\mathbf1)=v(\mathbf1\cdot\mathbf1)=2v(\mathbf1)$, skąd $v(\mathbf1)=0$; z liniowości,
$v(c)=0$ dla każdej funkcji stałej $c$.

Niech $h\equiv0$ na otoczeniu otwartym $W\ni p$. Wybierzmy mapę $(U,\varphi)$ wokół $p$ z
$U\subseteq W$, $\varphi(p)=0$, oraz (po przeskalowaniu)
$\overline{B(0,2)}\subseteq\varphi(U)$.
Zdefiniujmy $\tilde\mu\colon M\to\R$: $\tilde\mu(q):=\mu(\varphi(q))$ dla $q\in U$, oraz
$\tilde\mu(q):=0$ dla $q\notin U$. Nośnik jest zawarty w zwartym
zbiorze $K=\varphi^{-1}(\overline{B(0,2)})\subset U$.
Ponieważ $M$ jest Hausdorffa, $K$ jest domknięty w $M$.
Na otwartym pokryciu $U, M\setminus K$ obie gładkie definicje
pokrywają się, więc $\tilde\mu\in C^\infty(M)$, równa $1$ na otoczeniu $p$ i znikająca poza
$U\subseteq W$. Wówczas $\tilde\mu\cdot h\equiv0$ na całym $M$, więc $v(\tilde\mu h)=0$.
Z reguły Leibniza, $0=v(\tilde\mu h)=\tilde\mu(p)v(h)+h(p)v(\tilde\mu)=v(h)$ (bo
$\tilde\mu(p)=1$, $h(p)=0$). Zatem $v(h)=0$. Ostatnia część lematu wynika z zastosowania
powyższego do $h:=f-g$.
\end{proof}

\begin{twierdzenie}\label{tw:baza-styczna}
Niech $(U,\varphi)$ będzie mapą dopuszczalną wokół $p$,
$\varphi=(x^1,\dots,x^n)$,
gdzie $x^i=\pi_i\circ\varphi\colon U\to\R$ są funkcjami współrzędnymi,
czyli dla każdego $q\in U$ zachodzi
$\varphi(q)=\bigl(x^1(q),\dots,x^n(q)\bigr)$,
oraz załóżmy, że $\varphi(p)=0$. Wówczas:
\begin{enumerate}[label=(\alph*)]
\item Odwzorowanie $\Phi\colon T_p^AM\to T_p^CM$ jest izomorfizmem liniowym.
\item Elementy $\frac{\partial}{\partial x^1}\big|_p,\dots,\frac{\partial}{\partial x^n}
\big|_p\in T_p^CM$, zdefiniowane przez $\frac{\partial}{\partial x^i}\big|_p(f):=
\frac{\partial(f\circ\varphi^{-1})}{\partial x^i}(\varphi(p))$, stanowią bazę $T_p^CM$.
W szczególności $\dim T_pM=n$.
\end{enumerate}
\end{twierdzenie}

\begin{proof}
\emph{(a), injektywność.} Załóżmy $\Phi([\gamma])=0$, tzn.\ $(f\circ\gamma)'(0)=0$ dla
wszystkich $f\in C^\infty(M)$. Wybierzmy funkcję odcinającą
$\tilde\mu$ o zwartym nośniku zawartym w $U$, równą $1$ blisko $p$,
jak w dowodzie Lematu~\ref{lem:lokalnosc}.
Dla $i=1,\dots,n$ niech $\tilde x^i:=\tilde\mu\,x^i
=\tilde\mu\,\varphi^i$ na $U$ (rozszerzone przez $0$ poza $U$) —
funkcja gładka na $M$, pokrywająca się z $\varphi^i$ na otoczeniu $p$. Z założenia,
$0=\Phi([\gamma])(\tilde x^i)=(\tilde x^i\circ\gamma)'(0)$; blisko $t=0$, $\gamma(t)\in U$
i tam $\tilde x^i\circ\gamma=\varphi^i\circ\gamma$, więc $(\varphi^i\circ\gamma)'(0)=0$ dla
każdego $i$, czyli $(\varphi\circ\gamma)'(0)=0\in\R^n$. Oznacza to $\gamma\sim
\gamma_{\text{stała}}$ (gdzie $\gamma_{\text{stała}}(t)\equiv p$), czyli $[\gamma]=0$ w
$T_p^AM$. Zatem $\ker\Phi=\{0\}$.

\emph{(a), surjektywność — Lemat Hadamarda.} Niech $v\in T_p^CM$. Bez straty ogólności
$\varphi(U)$ jest kulą (gwiaździstą względem $0$). Niech $\tilde f:=f\circ\varphi^{-1}\in
C^\infty(\varphi(U))$; z Lematu~\ref{lem:hadamard} istnieją $g_1,\dots,g_n$ z
\[
\tilde f(x)=\tilde f(0)+\sum_ix^ig_i(x), \qquad g_i(0)=\frac{\partial\tilde f}{\partial x^i}
(0)=\frac{\partial}{\partial x^i}\Big|_p(f).
\]
Podstawiając $x=\varphi(q)$: $f(q)=f(p)+\sum_i\varphi^i(q)\,g_i(\varphi(q))$ dla $q$
blisko $p$. Niech $\widetilde{\varphi^i}$, $\tilde g_i$ oznaczają (jak wyżej) rozszerzenia
$\varphi^i$ oraz $g_i\circ\varphi$ do funkcji gładkich na $M$, pokrywające się z
oryginałami na otoczeniu $p$. Funkcja $F:=f(p)\cdot\mathbf1+\sum_i\widetilde{\varphi^i}\,
\tilde g_i$ pokrywa się z $f$ na otoczeniu $p$, więc $v(f)=v(F)$ (Lemat~\ref{lem:lokalnosc}).
Z liniowości $v$, reguły Leibniza (korzystając z $\widetilde{\varphi^i}(p)=0$,
$v(\mathbf1)=0$):
\[
v(f)=v(F)=\sum_i\Big(\widetilde{\varphi^i}(p)\,v(\tilde g_i)+\tilde g_i(p)\,v
(\widetilde{\varphi^i})\Big)=\sum_ig_i(0)\,v(\widetilde{\varphi^i})=\sum_{i=1}^n
v(\widetilde{\varphi^i})\cdot\frac{\partial}{\partial x^i}\Big|_p(f).
\]
Kładąc $c_i:=v(\widetilde{\varphi^i})\in\R$, mamy $v=\sum_ic_i\,\frac{\partial}{\partial
x^i}\big|_p$ jako funkcjonały. Niech $\gamma(t):=\varphi^{-1}(t\cdot(c_1,\dots,c_n))$; wtedy
$\Phi([\gamma])(f)=(f\circ\gamma)'(0)=D(f\circ\varphi^{-1})(0)\cdot(c_1,\dots,c_n)=
\sum_ic_i\frac{\partial}{\partial x^i}\big|_p(f)=v(f)$ dla każdego $f$, czyli
$\Phi([\gamma])=v$. To dowodzi surjektywności, a wraz z liniowością i injektywnością —
części (a).

\emph{(b).} Powyższy rachunek pokazał, że każdy $v\in T_p^CM$ jest kombinacją liniową
$\sum_ic_i\frac{\partial}{\partial x^i}\big|_p$ — więc baza ta rozpina $T_p^CM$. Niezależność
liniowa: załóżmy $\sum_ic_i\frac{\partial}{\partial x^i}\big|_p=0$. Stosując obie strony
do $\widetilde{\varphi^j}$ i korzystając z $\frac{\partial}{\partial x^i}\big|_p
(\widetilde{\varphi^j})=\frac{\partial x^j}{\partial x^i}(0)=\delta_i^j$, otrzymujemy
$0=\sum_ic_i\delta_i^j=c_j$ dla każdego $j$. Zatem $\{\partial/\partial x^i|_p\}$ jest
bazą $T_p^CM$; w szczególności $\dim T_p^CM=n$.
\end{proof}

\begin{wniosek}
Definicje A, B i C przestrzeni stycznej są kanonicznie izomorficzne. Piszemy odtąd po
prostu $T_pM$ i mówimy o \emph{przestrzeni stycznej do $M$ w punkcie $p$} — jest to
$n$-wymiarowa przestrzeń wektorowa. W mapie $(U,\varphi)$, gdzie $\varphi=(x^1,\dots,x^n)$,
jej \emph{bazą współrzędnościową} są wektory
\[
\left.\frac{\partial}{\partial x^1}\right|_p,\ \dots,\
\left.\frac{\partial}{\partial x^n}\right|_p.
\]
\end{wniosek}

\begin{figure}[ht]
\centering
\begin{tikzpicture}[>=Stealth,scale=1.35,line cap=round,line join=round]
% S(u,v)=(u,v,.12u^2+.08v^2), rzut (x,y,z)->(x+.45y,.65y+z).
% Brzeg wypełnienia, siatka i krzywa używają tej samej parametryzacji.
\filldraw[fill=manifoldgreen!12,draw=manifoldgreen!65!black,thick]
 (-2.6525,-.9425)--(1.3475,-.9425)--(2.6525,.9425)--(-1.3475,.9425)--cycle;
\path[fill=manifoldblue!18,opacity=.65]
 (-2.65250,-0.29430)
  --(-2.48583,-0.37097)
  --(-2.31917,-0.44097)
  --(-2.15250,-0.50430)
  --(-1.98583,-0.56097)
  --(-1.81917,-0.61097)
  --(-1.65250,-0.65430)
  --(-1.48583,-0.69097)
  --(-1.31917,-0.72097)
  --(-1.15250,-0.74430)
  --(-0.98583,-0.76097)
  --(-0.81917,-0.77097)
  --(-0.65250,-0.77430)
  --(-0.48583,-0.77097)
  --(-0.31917,-0.76097)
  --(-0.15250,-0.74430)
  --(0.01417,-0.72097)
  --(0.18083,-0.69097)
  --(0.34750,-0.65430)
  --(0.51417,-0.61097)
  --(0.68083,-0.56097)
  --(0.84750,-0.50430)
  --(1.01417,-0.44097)
  --(1.18083,-0.37097)
  --(1.34750,-0.29430)
  --(1.40187,-0.24262)
  --(1.45625,-0.18861)
  --(1.51063,-0.13226)
  --(1.56500,-0.07358)
  --(1.61938,-0.01256)
  --(1.67375,0.05080)
  --(1.72812,0.11649)
  --(1.78250,0.18452)
  --(1.83688,0.25489)
  --(1.89125,0.32759)
  --(1.94562,0.40263)
  --(2.00000,0.48000)
  --(2.05437,0.55971)
  --(2.10875,0.64176)
  --(2.16312,0.72614)
  --(2.21750,0.81286)
  --(2.27188,0.90191)
  --(2.32625,0.99330)
  --(2.38063,1.08703)
  --(2.43500,1.18309)
  --(2.48937,1.28149)
  --(2.54375,1.38222)
  --(2.59813,1.48529)
  --(2.65250,1.59070)
  --(2.48583,1.51403)
  --(2.31917,1.44403)
  --(2.15250,1.38070)
  --(1.98583,1.32403)
  --(1.81917,1.27403)
  --(1.65250,1.23070)
  --(1.48583,1.19403)
  --(1.31917,1.16403)
  --(1.15250,1.14070)
  --(0.98583,1.12403)
  --(0.81917,1.11403)
  --(0.65250,1.11070)
  --(0.48583,1.11403)
  --(0.31917,1.12403)
  --(0.15250,1.14070)
  --(-0.01417,1.16403)
  --(-0.18083,1.19403)
  --(-0.34750,1.23070)
  --(-0.51417,1.27403)
  --(-0.68083,1.32403)
  --(-0.84750,1.38070)
  --(-1.01417,1.44403)
  --(-1.18083,1.51403)
  --(-1.34750,1.59070)
  --(-1.40187,1.48529)
  --(-1.45625,1.38222)
  --(-1.51063,1.28149)
  --(-1.56500,1.18309)
  --(-1.61938,1.08703)
  --(-1.67375,0.99330)
  --(-1.72812,0.90191)
  --(-1.78250,0.81286)
  --(-1.83688,0.72614)
  --(-1.89125,0.64176)
  --(-1.94562,0.55971)
  --(-2.00000,0.48000)
  --(-2.05437,0.40263)
  --(-2.10875,0.32759)
  --(-2.16312,0.25489)
  --(-2.21750,0.18452)
  --(-2.27188,0.11649)
  --(-2.32625,0.05080)
  --(-2.38063,-0.01256)
  --(-2.43500,-0.07358)
  --(-2.48937,-0.13226)
  --(-2.54375,-0.18861)
  --(-2.59813,-0.24262) --cycle;
\foreach \vv in {-1.45,-.9,-.3,.3,.9,1.45}{
 \draw[manifoldblue!48,thin] plot[domain=-2:2,samples=35]
  ({\x+.45*\vv},{.65*\vv+.12*\x*\x+.08*\vv*\vv});}
\foreach \uu in {-2,-1.5,-1,-.5,0,.5,1,1.5,2}{
 \draw[manifoldblue!48,thin] plot[domain=-1.45:1.45,samples=30]
  ({\uu+.45*\x},{.65*\x+.12*\uu*\uu+.08*\x*\x});}
\foreach \vv in {-1.45,1.45}{
 \draw[manifoldblue!75!black,thick] plot[domain=-2:2,samples=35]
  ({\x+.45*\vv},{.65*\vv+.12*\x*\x+.08*\vv*\vv});}
\foreach \uu in {-2,2}{
 \draw[manifoldblue!75!black,thick] plot[domain=-1.45:1.45,samples=30]
  ({\uu+.45*\x},{.65*\x+.12*\uu*\uu+.08*\x*\x});}
\draw[manifoldred,very thick] plot[domain=-1.75:1.75,samples=40]
 ({\x},{.12*\x*\x});
\fill[manifoldred] (0,0) circle (2.5pt) node[below left,font=\small] {$p$};
\draw[->,very thick,manifoldred] (0,0)--(1.35,0);
\node[font=\small,manifoldred] at (-1.77,.57) {$\gamma$};
\node[font=\small\bfseries,manifoldblue!75!black] at (2.92,1.65) {$M$};
\node[font=\small,manifoldgreen!55!black] at (-2.0,-1.30) {$p+T_pM$};
\draw[->,manifoldgreen!60!black,thin] (-1.9,-1.10)--(-1.5,-.75);
\node[font=\small,manifoldred] at (1.7,-1.30) {$v=\gamma'(0)$};
\draw[->,manifoldred,thin] (1.45,-1.08)--(1.08,-.10);
\end{tikzpicture}
\caption{Przestrzeń styczna: zielony równoległobok przedstawia afiniczną
płaszczyznę $p+T_pM$ przy zanurzonej powierzchni $M$. Czerwona krzywa
$\gamma$ leży na $M$, a wektor $v=\gamma'(0)$ leży w płaszczyźnie stycznej.
Siatka powierzchni odchyla się od tej płaszczyzny poza punktem $p$.}
\label{fig:tangent}
\end{figure}


\subsection*{Zmiana współrzędnych wektora stycznego}

\begin{fakt}
Niech $(U,\varphi)$ i $(V,\psi)$ będą mapami wokół $p\in M$, a
$[\gamma]\in T_pM$, gdzie
\[
\gamma:(-\varepsilon,\varepsilon)\to M,
\qquad
\gamma(0)=p.
\]
Współrzędne wektora $[\gamma]$ w mapie $\varphi$ są dane przez
\[
[\gamma]_{\varphi}:=(\varphi\circ\gamma)'(0)\in\R^n,
\]
a w mapie $\psi$ przez
\[
[\gamma]_{\psi}:=(\psi\circ\gamma)'(0).
\]
Wtedy
\[
\boxed{
[\gamma]_{\psi}
=
D(\psi\circ\varphi^{-1})(\varphi(p))\,[\gamma]_{\varphi}.
}
\]
\end{fakt}

\begin{proof}
Ponieważ
\[
\psi\circ\gamma
=
(\psi\circ\varphi^{-1})\circ(\varphi\circ\gamma),
\]
z reguły łańcuchowej otrzymujemy
\[
(\psi\circ\gamma)'(0)
=
D(\psi\circ\varphi^{-1})(\varphi(p))
\,(\varphi\circ\gamma)'(0),
\]
czyli dokładnie
\[
[\gamma]_{\psi}
=
D(\psi\circ\varphi^{-1})(\varphi(p))\,[\gamma]_{\varphi}.
\]
\end{proof}


\begin{przyklad}[Zmiana współrzędnych wektora stycznego na $\mathbb{RP}^2$]
\label{prz:RP2-styczna-1}
Wróćmy do map z przykładu~\ref{prz:RP2-przejscie-1}:
punkt $[1:u:v]$, gdzie $u\ne0$, ma w drugiej mapie współrzędne
$(a,b)=(1/u,v/u)$. Zastosujemy do tego przejścia właśnie udowodnioną
regułę zmiany współrzędnych wektora stycznego.

Odwzorowanie $\varphi_1\circ\varphi_0^{-1}$ opisuje zmianę
\emph{współrzędnych punktów} rozmaitości. Jego różniczka opisuje natomiast
odpowiadającą jej zmianę \emph{współrzędnych wektorów stycznych}.

Niech $p\in U_0\cap U_1$ oraz
\[
\varphi_0(p)=(u,v).
\]
Mapy $\varphi_0$ i $\varphi_1$ wyznaczają dwie bazy przestrzeni stycznej
$T_p\mathbb{RP}^2$:
\[
\left\{
\left.\frac{\partial}{\partial u}\right|_p,
\left.\frac{\partial}{\partial v}\right|_p
\right\},
\qquad
\left\{
\left.\frac{\partial}{\partial a}\right|_p,
\left.\frac{\partial}{\partial b}\right|_p
\right\}.
\]
Jeżeli wektor $X\in T_p\mathbb{RP}^2$ ma w bazie indukowanej przez
$\varphi_0$ współrzędne
\[
X
=
\dot u\left.\frac{\partial}{\partial u}\right|_p
+
\dot v\left.\frac{\partial}{\partial v}\right|_p,
\]
to jego współrzędne w bazie indukowanej przez $\varphi_1$ otrzymujemy,
stosując różniczkę odwzorowania przejścia:
\[
\begin{pmatrix}
\dot a\\[1mm]
\dot b
\end{pmatrix}
=
D(\varphi_1\circ\varphi_0^{-1})(u,v)
\begin{pmatrix}
\dot u\\[1mm]
\dot v
\end{pmatrix}.
\]
A zatem
\[
\dot a=-\frac{\dot u}{u^2},
\qquad
\dot b
=
-\frac{v\dot u}{u^2}
+
\frac{\dot v}{u}.
\]
Ten sam wektor styczny możemy więc zapisać jako
\[
X
=
\dot u\left.\frac{\partial}{\partial u}\right|_p
+
\dot v\left.\frac{\partial}{\partial v}\right|_p
=
\dot a\left.\frac{\partial}{\partial a}\right|_p
+
\dot b\left.\frac{\partial}{\partial b}\right|_p.
\]

Innymi słowy, odwzorowanie przejścia
$\varphi_1\circ\varphi_0^{-1}$ zmienia współrzędne punktu na rozmaitości,
natomiast jego różniczka
\[
d(\varphi_1\circ\varphi_0^{-1})_{\varphi_0(p)}
\]
zmienia reprezentację współrzędnościową tego samego wektora
$X\in T_p\mathbb{RP}^2$ względem baz przestrzeni stycznej indukowanych
przez obie mapy.
\end{przyklad}

\subsection*{Przestrzeń styczna w punkcie brzegu}

Niech teraz $p\in\partial M$. \emph{Kiełek}\index{kielek funkcji@Kiełek funkcji} funkcji gładkiej w $p$
to jej klasa względem równości w pewnym otoczeniu punktu $p$
otwartym w $M$;
funkcje są tu gładkie w sensie przedłużalności z półprzestrzeni.
Ich algebra to $C^\infty_p(M)$. Definiujemy $T_pM$ jako przestrzeń
$\R$-liniowych derywacji $D:C^\infty_p(M)\to\R$ spełniających
\[
 D(fg)=f(p)D(g)+g(p)D(f).
\]
Wartość kiełka w $p$ jest dobrze określona. To lokalna wersja
Definicji C, która nie wymaga krzywej przechodzącej przez $p$
w obu kierunkach czasu.

Użycie kiełków jest tutaj istotne. Gładkie przedłużenie funkcji
$f$ poza brzeg, do otoczenia $\varphi(p)$ w $\R^n$, nie jest
jednoznaczne i stanowi jedynie pomocnicze narzędzie służące do
zdefiniowania gładkości. Przestrzeń styczna powinna natomiast zależeć
wyłącznie od lokalnego zachowania funkcji na samej rozmaitości $M$,
a nie od arbitralnego wyboru jej wartości poza $M$.

Dokładniej, dwa różne gładkie przedłużenia
$\widetilde f_1,\widetilde f_2$ mogą być różnymi funkcjami na
otoczeniu $\varphi(p)$ w $\R^n$, mimo że ich ograniczenia do
$\mathbb H^n$ wyznaczają tę samą funkcję na $M$. Nie chcemy więc
traktować samych przedłużeń jako obiektów, na których działa derywacja.
Obiektem właściwym jest kiełek funkcji na $M$, natomiast przedłużenie
wybieramy dopiero wtedy, gdy chcemy wykonać rachunek we współrzędnych.

Co więcej, derywacja w punkcie $p$ jest obiektem lokalnym: jej wartość
na funkcji zależy jedynie od zachowania tej funkcji w dowolnie małym
otoczeniu $p$. Dlatego naturalną dziedziną derywacji jest algebra
kiełków $C^\infty_p(M)$.

\begin{stwierdzenie}[Pełny wymiar przestrzeni stycznej przy brzegu]
\label{prop:styczna-brzeg-1}
W mapie półprzestrzennej $\varphi$ z $\varphi(p)=0$ każda taka
derywacja ma jednoznaczną postać
\[
D(f)=\sum_{i=1}^n v^i\partial_i\widetilde f(0),
\qquad v\in\R^n,
\]
gdzie $\widetilde f$ jest dowolnym gładkim przedłużeniem
$f\circ\varphi^{-1}$ na otoczenie zera w $\R^n$.
Wzór nie zależy od wyboru przedłużenia. Przy zmianie mapy współrzędne
$v$ zmieniają się przez odwracalną macierz Jacobiego przejścia.
W szczególności $\dim T_pM=n$, a inkluzja brzegu identyfikuje
$T_p\partial M$ z podprzestrzenią $v^n=0$.
\end{stwierdzenie}

\begin{proof}
Najpierw sprawdźmy, że prawa strona wzoru nie zależy od wyboru
przedłużenia $\widetilde f$. Jeżeli $\widetilde f_1$ i
$\widetilde f_2$ są dwoma gładkimi przedłużeniami reprezentującymi
ten sam kiełek funkcji $f\circ\varphi^{-1}$, to ich różnica znika
na pewnej małej półkuli
\[
B_\varepsilon(0)\cap\mathbb H^n.
\]
Na jej części wewnętrznej, gdzie $x^n>0$, znikają więc wszystkie
pierwsze pochodne tej różnicy, a przez ciągłość znikają one również
w punkcie $0$. W konsekwencji
\[
\partial_i\widetilde f_1(0)
=
\partial_i\widetilde f_2(0),
\qquad i=1,\ldots,n.
\]
Zatem dla każdego $v\in\R^n$ wzór
\[
D_v(f):=\sum_{i=1}^n v^i\partial_i\widetilde f(0)
\]
jest dobrze określony na kiełkach. Jest on liniowy i spełnia regułę
Leibniza, a więc definiuje derywację w punkcie $p$.

Pokażmy teraz odwrotnie, że każda derywacja $D$ ma taką postać.
Najpierw
\[
D(1)=D(1\cdot1)=2D(1),
\]
zatem $D(1)=0$, a przez liniowość $D$ znika na wszystkich funkcjach
stałych.

Wybierzmy gładkie przedłużenie $\widetilde f$ funkcji
$f\circ\varphi^{-1}$ na małą kulę wokół zera. Lemat Hadamarda daje
\[
\widetilde f(x)
=
\widetilde f(0)+\sum_{i=1}^n x^i g_i(x),
\qquad
g_i(0)=\partial_i\widetilde f(0).
\]
Po ograniczeniu do półkuli otrzymujemy równość funkcji
reprezentujących kiełki w $0$, a więc
\[
[f\circ\varphi^{-1}]_0
=
\widetilde f(0)[1]_0
+
\sum_{i=1}^n [x^i]_0[g_i]_0.
\]

Niech również $x^i=\pi_i\circ\varphi$ oznacza $i$-tą funkcję
współrzędną na $U$. Stosując derywację $D$ do powyższej równości
kiełków oraz korzystając z $D(1)=0$ i $x^i(p)=0$, otrzymujemy
\[
D(f)
=
\sum_{i=1}^n D(x^i)\,g_i(0).
\]
Ponieważ
\[
g_i(0)=\partial_i\widetilde f(0),
\]
mamy
\[
D(f)
=
\sum_{i=1}^n D(x^i)\,\partial_i\widetilde f(0).
\]
Definiując
\[
v^i:=D(x^i),
\qquad i=1,\ldots,n,
\]
otrzymujemy żądaną postać
\[
D(f)
=
\sum_{i=1}^n v^i\partial_i\widetilde f(0).
\]

Współczynniki $v^i$ są jednoznaczne. Istotnie, stosując powyższy
wzór do funkcji współrzędnej $x^j$, otrzymujemy
\[
D(x^j)
=
\sum_{i=1}^n v^i\partial_i x^j(0)
=
\sum_{i=1}^n v^i\delta_i^j
=
v^j.
\]

Rozważmy dwie mapy półprzestrzenne $\varphi$ i $\psi$.
Odwzorowania przejścia
\[
\varphi\circ\psi^{-1}
\qquad\text{oraz}\qquad
\psi\circ\varphi^{-1}
\]
są wzajemnie odwrotne. Na punktach odpowiadających wnętrzu $M$ mamy zatem
\[
(\psi\circ\varphi^{-1})\circ(\varphi\circ\psi^{-1})
=\operatorname{id}
\]
oraz
\[
(\varphi\circ\psi^{-1})\circ(\psi\circ\varphi^{-1})
=\operatorname{id}.
\]
Stosując regułę łańcuchową do gładkich przedłużeń odwzorowań
przejścia, otrzymujemy na punktach wnętrza
\[
D(\psi\circ\varphi^{-1})(\varphi(q))\,
D(\varphi\circ\psi^{-1})(\psi(q))
=I.
\]
Przechodząc z $q\in\operatorname{int}M$ do $p\in\partial M$
i korzystając z ciągłości różniczek, otrzymujemy
\[
D(\psi\circ\varphi^{-1})(\varphi(p))\,
D(\varphi\circ\psi^{-1})(\psi(p))
=I.
\]
Analogicznie iloczyn w odwrotnej kolejności jest równy $I$.
Zatem macierze Jacobiego odwzorowań przejścia w punkcie brzegowym
są wzajemnie odwrotne.

Wreszcie ograniczenie mapy $\varphi$ do brzegu ma współrzędne
$(x^1,\ldots,x^{n-1})$. Różniczka inkluzji
$\partial M\hookrightarrow M$ dodaje zerową ostatnią współrzędną,
a zatem identyfikuje $T_p\partial M$ z podprzestrzenią
\[
\{v\in\R^n:v^n=0\}\subset T_pM.
\]
\end{proof}

Można też użyć derywacji funkcji globalnych: lokalność dowodzimy
jak wcześniej, a każdy kiełek ma globalnego reprezentanta po
pomnożeniu przez funkcję odcinającą równą $1$ blisko $p$,
o nośniku zawartym w dziedzinie reprezentanta, i przedłużeniu przez zero.
Taką funkcję otrzymujemy przez ograniczenie
euklidesowej funkcji odcinającej do półprzestrzeni.

Natomiast model przestrzeni stycznej za pomocą dwustronnych krzywych
nie daje przy brzegu całej przestrzeni $T_pM$. Istotnie, niech
\[
\gamma:(-\varepsilon,\varepsilon)\to M,
\qquad
\gamma(0)=p\in\partial M,
\]
i wybierzmy mapę półprzestrzenną $\varphi$ taką, że
$\varphi(p)=0$ oraz
\[
\varphi(U)\subset\mathbb H^n
=
\{x\in\R^n:x^n\geq0\}.
\]
Ponieważ $\gamma(t)\in M$, dla dostatecznie małych $t$ mamy
\[
(x^n\circ\gamma)(t)\geq0,
\qquad
(x^n\circ\gamma)(0)=0.
\]
Funkcja $x^n\circ\gamma$ ma więc w $t=0$ minimum lokalne, a zatem
\[
(x^n\circ\gamma)'(0)=0.
\]
W konsekwencji prędkość krzywej w mapie ma postać
\[
(\varphi\circ\gamma)'(0)
=
(v^1,\ldots,v^{n-1},0),
\]
czyli należy do podprzestrzeni odpowiadającej
$T_p\partial M$.

Odwrotnie, każdy wektor
\[
v=(v^1,\ldots,v^{n-1},0)\in T_p\partial M
\]
jest realizowany przez dwustronną krzywą
\[
\gamma(t)
=
\varphi^{-1}(tv^1,\ldots,tv^{n-1},0)
\]
dla dostatecznie małych $|t|$. Dwustronne krzywe przechodzące przez
punkt brzegowy opisują więc dokładnie $T_p\partial M$, a nie pełną
przestrzeń
\[
T_pM\simeq\R^n.
\]
W szczególności nie mogą realizować wektorów o niezerowej składowej
normalnej $v^n$. To wyjaśnia, dlaczego w punkcie brzegu model
dwustronnych krzywych nie może służyć jako definicja pełnej przestrzeni
stycznej.

\section{Przestrzeń kostyczna}\label{sec:przestrzen-kostyczna}

\begin{definicja}
\emph{Przestrzenią kostyczną} do $M$ w punkcie $p$ nazywamy przestrzeń dualną
$T_p^*M:=(T_pM)^*$, tzn.\ przestrzeń funkcjonałów liniowych $T_pM\to\R$. Jej elementy
nazywamy \emph{kowektorami} w $p$.
\end{definicja}

Dla $f\in C^\infty(M)$ definiujemy \emph{różniczkę} $f$ w punkcie $p$, $df_p\in T_p^*M$,
wzorem
\[
df_p(v):=v(f), \qquad v\in T_pM
\]
(korzystając z utożsamienia $T_pM=T_p^CM$ — wektor styczny $v$ jest derywacją, a $v(f)\in
\R$ jest dobrze określoną liczbą). Odwzorowanie $df_p$ jest liniowe względem $v$ (z
definicji przestrzeni wektorowej derywacji), a więc istotnie $df_p\in T_p^*M$.

Niech $(U,\varphi)$, gdzie $\varphi=(x^1,\dots,x^n)$, będzie mapą wokół $p$. Rozpatrzmy różniczki funkcji
współrzędnościowych, $dx^i_p\in T_p^*M$ (traktując lokalnie zdefiniowaną funkcję
$x^i=\varphi^i$ jako element $C^\infty$ w otoczeniu $p$; z Lematu~\ref{lem:lokalnosc} o
lokalności derywacji, wartość $v(x^i)$ obliczamy na dowolnej funkcji
gładkiej na $M$, która zgadza się z $x^i$ blisko $p$; funkcję taką
otrzymujemy przez odcięcie wewnątrz $U$ i przedłużenie przez zero).

\begin{fakt}\label{fakt:dual-baza}
$dx^i_p\Big(\dfrac{\partial}{\partial x^j}\Big|_p\Big)=\delta^i_j$.
\end{fakt}
\begin{proof}
Z definicji, $dx^i_p\big(\partial/\partial x^j|_p\big)=\frac{\partial}{\partial x^j}
\Big|_p(x^i)=\frac{\partial(x^i\circ\varphi^{-1})}{\partial x^j}(\varphi(p))$. Ale
$x^i\circ\varphi^{-1}\colon\varphi(U)\to\R$ to rzutowanie na $i$-tą współrzędną,
$y\mapsto y^i$, którego pochodna cząstkowa po $j$-tej zmiennej wynosi $\delta^i_j$ w
każdym punkcie.
\end{proof}

Przypomnijmy ogólny fakt z algebry liniowej.

\begin{fakt}[Baza dualna]\label{fakt:baza-dualna}
Jeśli $e_1,\dots,e_n$ jest bazą przestrzeni wektorowej $V$ (wymiaru $n$), a $e^1,\dots,
e^n\in V^*$ są funkcjonałami spełniającymi $e^i(e_j)=\delta^i_j$, to $e^1,\dots,e^n$ jest
bazą $V^*$ (,,bazą dualną'' do $e_1,\dots,e_n$).
\end{fakt}
\begin{proof}
Rozpinanie: dla dowolnego $\omega\in V^*$, kładziemy $c_i:=\omega(e_i)$ i sprawdzamy, że
$\omega=\sum_ic_ie^i$ — obie strony, jako funkcjonały liniowe, pokrywają się na bazie
$e_1,\dots,e_n$: $\big(\sum_ic_ie^i\big)(e_j)=\sum_ic_i\delta^i_j=c_j=\omega(e_j)$;
funkcjonały liniowe pokrywające się na bazie są równe. Niezależność liniowa: jeśli
$\sum_ic_ie^i=0$, to podstawiając $e_j$ dla każdego $j$: $0=\sum_ic_i\delta^i_j=c_j$.
Zatem układ $e^1,\dots,e^n$ rozpina $V^*$ i jest liniowo niezależny, czyli jest bazą (w
szczególności $\dim V^*=\dim V=n$).
\end{proof}

Łącząc Fakt~\ref{fakt:dual-baza} z Faktem~\ref{fakt:baza-dualna} (dla $V=T_pM$, $e_i=
\partial/\partial x^i|_p$, $e^i=dx^i_p$), otrzymujemy odpowiedź na pytanie postawione we
wstępie: \emph{kowektory $dx^1_p,\dots,dx^n_p$ stanowią bazę $T_p^*M$ dokładnie dlatego, że
są bazą dualną do bazy $\{\partial/\partial x^i|_p\}$ przestrzeni stycznej} — żaden
dodatkowy argument nie jest potrzebny, wystarczy ogólny fakt algebry liniowej o istnieniu
i jednoznaczności bazy dualnej.

\begin{stwierdzenie}
Dla $f\in C^\infty(M)$,
\[
df_p=\sum_{i=1}^n\frac{\partial f}{\partial x^i}\Big|_p\,dx^i_p.
\]
\end{stwierdzenie}
\begin{proof}
Obie strony są elementami $T_p^*M$; wystarczy sprawdzić, że działają tak samo na każdym
wektorze bazy $\partial/\partial x^j|_p$. Lewa strona: $df_p(\partial/\partial x^j|_p)=
\frac{\partial}{\partial x^j}\big|_p(f)$ z definicji $df_p$. Prawa strona, na mocy
Faktu~\ref{fakt:dual-baza}: $\sum_i\frac{\partial f}{\partial x^i}\big|_p\,dx^i_p
(\partial/\partial x^j|_p)=\sum_i\frac{\partial f}{\partial x^i}\Big|_p\delta^i_j=
\frac{\partial f}{\partial x^j}\Big|_p$. Obie strony są równe.
\end{proof}
\section{Wiązki wektorowe}\label{sec:wiazki-wektorowe}

\begin{definicja}
\emph{Wiązką wektorową} rzędu $k$ nad rozmaitością gładką $M$ (wymiaru $n$) nazywamy
rozmaitość gładką $E$ (wymiaru $n+k$) wraz z odwzorowaniem gładkim $\pi\colon E\to M$
(\emph{rzutowaniem}) takim, że:
\begin{enumerate}[label=(\roman*)]
\item dla każdego $p\in M$, włókno $E_p:=\pi^{-1}(p)$ jest wyposażone w strukturę
$k$-wymiarowej rzeczywistej przestrzeni wektorowej;
\item (lokalna trywializacja) każdy punkt $p\in M$ posiada otoczenie otwarte $U$ oraz
dyfeomorfizm $\Psi\colon\pi^{-1}(U)\to U\times\R^k$ z $\mathrm{pr}_1\circ\Psi=\pi$ na
$\pi^{-1}(U)$, taki że dla każdego $q\in U$, obcięcie $\Psi|_{E_q}\colon E_q\to\{q\}\times
\R^k\cong\R^k$ jest izomorfizmem liniowym.
\end{enumerate}
\end{definicja}

\begin{figure}[ht]
\centering
\begin{tikzpicture}[>=Stealth,scale=1.05,line cap=round]
\fill[manifoldblue!12] (-3.1,1.0)
 .. controls (-2.1,.92) and (-1.1,1.15) .. (0,1.04)
 .. controls (1.1,.94) and (2.1,1.18) .. (3.1,1.08)
 --(3.1,-.39)
 .. controls (2.1,-.30) and (1.1,-.52) .. (0,-.43)
 .. controls (-1.1,-.35) and (-2.1,-.55) .. (-3.1,-.47) --cycle;
\draw[manifoldblue!75!black,very thick]
 plot[domain=-3.1:3.1,samples=35] ({\x},{1.04+.10*sin(50*\x)});
\draw[manifoldblue!75!black,thick]
 plot[domain=-3.1:3.1,samples=35] ({\x},{-.43+.10*sin(50*\x)});
\foreach \xx in {-2.65,-1.95,-1.3,-.65,.65,1.3,1.95,2.65}
 \draw[manifoldblue!52,thin] (\xx,{-.43+.10*sin(50*\xx)})
 --(\xx,{1.04+.10*sin(50*\xx)});
\draw[manifoldgreen!65!black,very thick] (0,-.43)--(0,1.04);
\draw[manifoldgreen!65!black,very thick]
 plot[domain=-3.15:3.15,samples=40] ({\x},{-2.04+.10*sin(50*\x)});
\fill[manifoldred] (0,.42) circle (2.6pt) node[right,font=\small] {$e\in E_p$};
\fill[manifoldred] (0,-2.04) circle (2.6pt) node[below,font=\small] {$p=\pi(e)$};
\draw[->,thick,manifoldred] (0,-.52)--(0,-1.9)
 node[midway,right,font=\small] {$\pi$};
\node[manifoldblue!75!black,font=\large\bfseries] at (-2.9,1.4) {$E$};
\node[manifoldgreen!60!black,font=\large\bfseries] at (3.45,-2.03) {$M$};
\end{tikzpicture}
\caption{Lokalny obraz wiązki $\pi\colon E\to M$. Wyróżniono wektor $e$ i jego rzut $\pi(e)=p$. Włókna narysowano jako fragmenty prostych; ogólnie mają wymiar $k$.}
\label{fig:diagram-wiazki}
\end{figure}







\begin{definicja}[Trywializacja wiązki]
\index{trywializacja wiazki@Trywializacja wiązki}
Niech $\pi\colon E\to M$ będzie wiązką wektorową rzędu $k$. \emph{Trywializacją}
nad otwartym zbiorem $U\subseteq M$ nazywamy dyfeomorfizm
\[
\Phi\colon\pi^{-1}(U)\longrightarrow U\times\mathbb R^k
\]
spełniający $\operatorname{pr}_1\circ\Phi=\pi$, taki że dla każdego $p\in U$
obcięcie $\Phi|_{E_p}$ jest izomorfizmem liniowym włókna $E_p$ na
$\{p\}\times\mathbb R^k\cong\mathbb R^k$.
\end{definicja}

\begin{figure}[ht]
\centering
\begin{tikzpicture}[>=Stealth,scale=.95,line cap=round]
\fill[manifoldblue!13] (-4.4,1.16)
 .. controls (-3.55,1.22) and (-2.2,.93) .. (-1.35,.98)
 --(-1.35,-.46)
 .. controls (-2.2,-.51) and (-3.55,-.23) .. (-4.4,-.28) --cycle;
\draw[manifoldblue!75!black,thick]
 plot[domain=-4.4:-1.35,samples=25] ({\x},{1.08+.11*sin(52*\x)});
\draw[manifoldblue!75!black,thick]
 plot[domain=-4.4:-1.35,samples=25] ({\x},{-.36+.11*sin(52*\x)});
\foreach \xx in {-3.9,-3.35,-2.8,-2.25,-1.7}
 \draw[manifoldblue!50,thin] (\xx,{-.36+.11*sin(52*\xx)})
 --(\xx,{1.08+.11*sin(52*\xx)});
\fill[manifoldgreen!13] (1.35,-.36) rectangle (4.4,1.08);
\draw[manifoldgreen!65!black,thick] (1.35,-.36) rectangle (4.4,1.08);
\foreach \xx in {1.85,2.4,2.95,3.5,4.05}
 \draw[manifoldgreen!55,thin] (\xx,-.36)--(\xx,1.08);
\fill[manifoldred] (-2.8,.6) circle (2.4pt)
 node[below,font=\small] {$e\in E_p$};
\fill[manifoldred] (2.95,.6) circle (2.4pt)
 node[below,font=\small] {$\Phi(e)=(p,a)$};
\draw[->,very thick,manifoldred] (-1.14,.39)--(1.12,.39)
 node[midway,above,font=\small] {$\Phi$};
\draw[manifoldblue!75!black,thick]
 plot[domain=-4.4:-1.35,samples=25] ({\x},{-1.6+.07*sin(52*\x)});
\draw[manifoldgreen!65!black,thick] (1.35,-1.6)--(4.4,-1.6);
\fill[manifoldred] (-2.8,-1.6) circle (2.2pt) node[below,font=\small] {$p$};
\fill[manifoldred] (2.95,-1.6) circle (2.2pt) node[below,font=\small] {$p$};
\draw[->,manifoldblue!70!black,thick] (-2.8,-.47)--(-2.8,-1.4)
 node[midway,left,font=\small] {$\pi$};
\draw[->,manifoldgreen!65!black,thick] (2.95,-.47)--(2.95,-1.4)
 node[midway,right,font=\small] {$\operatorname{pr}_1$};
\draw[->,manifoldgold!75!black,thick] (-1.12,-1.6)--(1.12,-1.6)
 node[midway,above,font=\small] {$\mathrm{id}_U$};
\node[manifoldblue!75!black,font=\small\bfseries] at (-2.85,1.5) {$E|_U$};
\node[manifoldgreen!65!black,font=\small\bfseries] at (2.9,1.5) {$U\times\R^k$};
\end{tikzpicture}
\caption{Trywializacja $\Phi\colon E|_U\to U\times\R^k$ jest liniowa na włóknach i zachowuje punkt bazy: $\operatorname{pr}_1\circ\Phi=\pi$. Linie przedstawiają włókna schematycznie.}
\label{fig:trywializacja}
\end{figure}

\begin{definicja}[Sekcja wiązki wektorowej]
\index{sekcja wiazki wektorowej@Sekcja wiązki wektorowej}
Niech $\pi\colon E\to M$ będzie wiązką wektorową. \emph{Sekcją}
(inaczej: \emph{przekrojem}) wiązki $E$ nazywamy
gładkie odwzorowanie $s\colon M\to E$ spełniające
\[
\pi\circ s=\id_M.
\]
Zatem dla każdego $p\in M$ mamy $s(p)\in E_p$. Sekcja wybiera dokładnie jeden wektor
z każdego włókna i robi to w sposób gładki.
\end{definicja}

\begin{figure}[ht]
\centering
\begin{tikzpicture}[>=Stealth,scale=1.06,line cap=round]
\fill[manifoldblue!11] (-3.05,1.07)
 .. controls (-1.8,.99) and (-.8,1.22) .. (0,1.13)
 .. controls (1.6,1.05) and (2.2,1.27) .. (3.05,1.19)
 --(3.05,-.32)
 .. controls (2.2,-.26) and (1.6,-.46) .. (0,-.38)
 .. controls (-.8,-.29) and (-1.8,-.52) .. (-3.05,-.44) --cycle;
\draw[manifoldblue!75!black,thick]
 plot[domain=-3.05:3.05,samples=34] ({\x},{1.13+.11*sin(48*\x)});
\draw[manifoldblue!75!black,thick]
 plot[domain=-3.05:3.05,samples=34] ({\x},{-.38+.11*sin(48*\x)});
\foreach \xx in {-2.65,-1.95,-1.3,-.65,0,.65,1.3,1.95,2.65}
 \draw[manifoldblue!50,thin] (\xx,{-.38+.11*sin(48*\xx)})
 --(\xx,{1.13+.11*sin(48*\xx)});
\draw[manifoldgreen!65!black,very thick]
 plot[domain=-3.12:3.12,samples=38] ({\x},{-1.85+.09*sin(48*\x)});
\draw[manifoldred,very thick]
 plot[domain=-3.05:3.05,samples=40] ({\x},{.34+.13*sin(48*\x)});
\fill[manifoldred] (0,.34) circle (2.6pt) node[above left,font=\small] {$s(p)$};
\fill[manifoldred] (0,-1.85) circle (2.6pt) node[below,font=\small] {$p$};
\draw[->,manifoldred,thick] (0,-.5)--(0,-1.7)
 node[midway,right,font=\small] {$\pi$};
\node[font=\small\bfseries,manifoldred!85!black] at (3.7,.72) {$s(U)$};
\node[font=\small\bfseries,manifoldgreen!65!black] at (3.43,-2.22) {$U\subset M$};
\node[font=\small,manifoldblue!75!black] at (-3.5,1.35) {$E|_U$};
\end{tikzpicture}
\caption{Sekcja $s$ wybiera dokładnie jeden wektor $s(p)$ w każdym włóknie nad $p$. Czerwona krzywa przedstawia jej obraz nad pokazanym fragmentem $U$.}
\label{fig:sekcja}
\end{figure}

\begin{przyklad}[Sekcja zerowa]
Każda wiązka wektorowa posiada naturalną \emph{sekcję zerową}
\[
s_0\colon M\to E,\qquad s_0(p)=0_p,
\]
gdzie $0_p$ jest wektorem zerowym we włóknie $E_p$.
Gładkość wynika z lokalnego zapisu $p\mapsto(p,0)$.
\end{przyklad}


\begin{przyklad}[Wiązka trywialna]
Najprostszym przykładem jest $E=M\times\R^k$ z $\pi=\mathrm{pr}_1$ i $\Psi=\id$ — wiązka ta
nosi nazwę \emph{trywialnej}.
\end{przyklad}

Izomorfizmem wiązek nad $M$ nazywamy dyfeomorfizm przestrzeni
całkowitych zachowujący punkt bazy i liniowy na każdym włóknie.
Nie każda wiązka wektorowa jest (globalnie) izomorficzna z wiązką trywialną — mimo że
lokalnie, z definicji, każda wiązka ,,wygląda jak'' trywialna. Poniższy klasyczny przykład
ilustruje to zjawisko w najprostszym możliwym przypadku.

\begin{przyklad}[Wstęga Möbiusa jako nietrywialna wiązka rzędu 1]\label{prz:mobius}
Niech $M=S^1=\R/2\pi\Z$. Wstęgę Möbiusa można zdefiniować jako przestrzeń
ilorazową
\[
E:=\big([0,2\pi]\times\R\big)\big/\!\sim,
\qquad
(0,t)\sim(2\pi,-t),
\]
z rzutowaniem
\[
\pi\colon E\longrightarrow S^1,\qquad
\pi([\theta,t])=[\theta].
\]
Włókno nad punktem $[\theta]\in S^1$ jest kopią $\R$:
\[
E_{[\theta]}=\pi^{-1}([\theta])\cong\R.
\]
W szczególności $E$ ma wymiar $2$, natomiast baza $S^1$ ma wymiar $1$.

\medskip
\emph{Dlaczego jest to wiązka wektorowa?}
Jeżeli $I\subset S^1$ jest otwartym łukiem, który nie zawiera punktu
odpowiadającego jednocześnie $\theta=0$ i $\theta=2\pi$, można wybrać jednoznacznego
reprezentanta kąta $\theta$ w przedziale odpowiadającym $I$. Wtedy
\[
\Psi_I\colon \pi^{-1}(I)\longrightarrow I\times\R,
\qquad
[\theta,t]\longmapsto([\theta],t)
\]
jest lokalną trywializacją i na każdym włóknie jest izomorfizmem liniowym.
Także punkt sklejenia ma trywializację. Na łuku
$I_0=\{[a]:-\varepsilon<a<\varepsilon\}$, gdzie $0<\varepsilon<\pi$,
jej odwrotność jest dana przez
\[
 ([a],v)\longmapsto
 \begin{cases}
 [a,v],&a\geq0,\\
 [a+2\pi,-v],&a<0.
 \end{cases}
\]
W punkcie $a=0$ zgodność zapewnia identyfikacja końców.
Jest to homeomorfizm z $I_0\times\R$ na $\pi^{-1}(I_0)$:
w topologii ilorazowej otoczenie sklejonego punktu obejmuje
oba przyległe półpasy. Przejścia z poprzednią trywializacją
są $(p,v)\mapsto(p,v)$ po jednej stronie łuku i
$(p,v)\mapsto(p,-v)$ po drugiej, więc są gładkie.
Te mapy określają strukturę gładką na $E$. Punkty o różnych
podstawach rozdzielamy otoczeniami w $S^1$, a punkty tego samego
włókna w jednej trywializacji; zatem $E$ jest Hausdorffa.
Dwie trywializacje nad łukami dają przeliczalną bazę $E$.
Po obejściu całej bazy współrzędna włókna zmienia znak:
\[
t\longmapsto -t.
\]
Jest to istotna informacja globalna: lokalnie wiązka jest nierozróżnialna od
$S^1\times\R$, ale globalnie nie można wybrać w niej spójnie jednego kierunku
w każdym włóknie.

W języku funkcji przejścia można to zapisać jeszcze prościej. Dla wiązki rzędu $1$
funkcje przejścia przyjmują wartości w
$\mathrm{GL}(1,\R)=\R_*$. W przypadku wstęgi Möbiusa pojawia się
przejście
\[
g=-1,
\]
podczas gdy dla wiązki trywialnej można wybrać $g=+1$. Znak $-1$ jest właśnie
algebraicznym zapisem ``odwrócenia'' włókna przy jednym pełnym okrążeniu bazy.
\end{przyklad}

\begin{figure}[h]
\centering
\begin{minipage}{0.43\textwidth}
\centering
\begin{tikzpicture}[>=Stealth,scale=0.82]
  \fill[manifoldblue!8] (0,-1.25) rectangle (6.3,1.25);
  \draw[thick,manifoldblue] (0,-1.25) rectangle (6.3,1.25);
  \foreach \x in {0.9,1.8,2.7,3.6,4.5,5.4}
    \draw[manifoldblue!35] (\x,-1.25)--(\x,1.25);
  \draw[very thick,->,manifoldred] (-0.05,1.05) -- (-0.05,-1.05);
  \draw[very thick,->,manifoldred] (6.35,-1.05) -- (6.35,1.05);
  \node[above] at (0,1.28) {$0$};
  \node[above] at (6.3,1.28) {$2\pi$};
  \node[manifoldred] at (3.15,-1.65) {$t\mapsto-t$};
\end{tikzpicture}

\smallskip
\small (a) Prostokąt przed sklejeniem
\end{minipage}
\hfill
\begin{minipage}{0.53\textwidth}
\centering
\begin{tikzpicture}
\begin{axis}[
  width=\linewidth,
  height=5.4cm,
  view={55}{25},
  axis lines=none,
  ticks=none,
  z buffer=sort,
  clip=false
]
\addplot3[
  surf,
  shader=faceted,
  faceted color=manifoldgreen!55!black,
  draw=manifoldgreen!70!black,
  domain=0:360,
  y domain=-0.55:0.55,
  samples=34,
  samples y=8
]
({(1+y*cos(x/2))*cos(x)},
 {(1+y*cos(x/2))*sin(x)},
 {y*sin(x/2)});
\end{axis}
\end{tikzpicture}

\smallskip
\small (b) Model geometryczny w $\R^3$
\end{minipage}
\caption{Wstęga Möbiusa: (a) konstrukcja przez sklejenie boków prostokąta z
odwróceniem współrzędnej włókna, (b) odpowiadająca jej klasyczna powierzchnia
zanurzona w $\R^3$. Parametryzacja w części (b) ma postać
$x=(1+t\cos(\theta/2))\cos\theta$,
$y=(1+t\cos(\theta/2))\sin\theta$,
$z=t\sin(\theta/2)$. Rysunek pokazuje tylko pas $|t|\leq0.55$;
w całej wiązce wektorowej $t$ przebiega $\R$.}
\label{fig:mobius}
\end{figure}

\begin{stwierdzenie}[Nietrywialność wstęgi Möbiusa]
Wiązka $E$ z Przykładu~\ref{prz:mobius} nie jest izomorficzna jako wiązka wektorowa
z wiązką trywialną $S^1\times\R$.
\end{stwierdzenie}

\begin{proof}
Niech $q\colon[0,2\pi]\times\R\to E$ oznacza rzutowanie ilorazowe. Każdej ciągłej
funkcji $f\colon[0,2\pi]\to\R$ spełniającej
\[
f(2\pi)=-f(0)
\]
odpowiada ciągły przekrój
\[
s_f\colon S^1\to E,\qquad s_f([\theta])=q(\theta,f(\theta)).
\]
Warunek brzegowy jest dokładnie warunkiem dobrze określoności w punkcie sklejenia.
Odwrotnie, każdy ciągły przekrój ma tę postać: nad przedziałem
$[0,2\pi]$ wybieramy współrzędną włóknową $t$, a ciągłość w
trywializacji przy sklejeniu daje ciągłość $f$ i podany warunek końcowy.
Przekrój jest gładki dokładnie wtedy, gdy $f$ jest gładka na
przedziale (z jednostronnymi pochodnymi na końcach) i
$f^{(j)}(2\pi)=-f^{(j)}(0)$ dla każdego $j\geq0$.
Sam warunek na wartości zapewnia ciągłość, lecz nie gładkość.

Załóżmy, że $s_f$ jest nigdzie zerowa. Wtedy $f$ nie przyjmuje wartości $0$.
Ponieważ $[0,2\pi]$ jest spójny, a $f$ jest ciągła, z twierdzenia o wartości
pośredniej wynika, że $f$ ma stały znak. W szczególności $f(0)$ i $f(2\pi)$
mają ten sam znak. Jest to niemożliwe, gdyż
$f(2\pi)=-f(0)$ i $f(0)\neq0$.

Zatem $E$ nie posiada nigdzie niezerowej sekcji. Natomiast wiązka trywialna
$S^1\times\R$ posiada taką sekcję, mianowicie
\[
s(\theta)=(\theta,1).
\]
Gdyby istniał izomorfizm wiązek
$\Psi\colon S^1\times\R\to E$, sekcja $\Psi\circ s$ byłaby nigdzie niezerową
sekcją $E$, co daje sprzeczność. Zatem
\[
E\not\cong S^1\times\R.
\]
\end{proof}

\begin{uwaga}[Interpretacja topologiczna]
Wstęga Möbiusa jest przykładem wiązki wektorowej, która jest lokalnie trywialna,
ale nie jest globalnie trywialna. To rozróżnienie jest fundamentalne:
\[
\boxed{\text{lokalnie trywialna}\ \not\Rightarrow\ \text{globalnie trywialna}.}
\]
Przestrzeń całkowita $E$, z włóknem $\R$, jest dwuwymiarową
rozmaitością \emph{bez brzegu}; bywa nazywana otwartą wstęgą Möbiusa.
Jej środkowy przekrój $\theta\mapsto[\theta,0]$ jest przekrojem zerowym.
Dopiero ograniczenie $|t|\leq c$ dla $c>0$ daje zwartą wstęgę
Möbiusa z brzegiem, widoczną na rysunku. Jej włókna są odcinkami,
więc ten ograniczony pas nie jest wiązką wektorową.
Odcinki $t=c$ i $t=-c$ sklejają się w jedną składową brzegu.
\end{uwaga}

\begin{przyklad}[Wiązka tautologiczna nad $\mathbb{RP}^1$]
Przykład wstęgi Möbiusa ma bardzo naturalny związek z przestrzenią rzutową.
Niech
\[
\gamma^1:=\{([x],v)\in\mathbb{RP}^1\times\R^2:\ v\in\operatorname{span}(x)\},
\]
gdzie $x\in\R^2\setminus\{0\}$ jest dowolnym reprezentantem klasy $[x]$.
Włókno nad $[x]$ jest zatem prostą
\[
\gamma^1_{[x]}=\operatorname{span}(x)\subset\R^2.
\]
Jest to \emph{wiązka tautologiczna} (kanoniczna wiązka liniowa) nad $\mathbb{RP}^1$.

Ponieważ $\mathbb{RP}^1$ jest dyfeomorficzna z okręgiem $S^1$, a przejście
z kierunku $x$ do tego samego kierunku po jednym pełnym okrążeniu wymusza zmianę
reprezentanta z $x$ na $-x$, wiązka $\gamma^1$ jest izomorficzna z wstęgą Möbiusa.
Jest to więc drugi, geometryczny sposób zobaczenia tej samej nietrywialnej wiązki:
w każdym punkcie przestrzeni rzutowej wybieramy \emph{linię} w $\R^2$, a nie
wyróżniony wektor tej linii.
Jawny izomorfizm nad dyfeomorfizmem baz
$[\theta]\mapsto[\cos(\theta/2):\sin(\theta/2)]$ ma postać
\[
 [\theta,t]\longmapsto
 \bigl([\cos(\theta/2):\sin(\theta/2)],
       t(\cos(\theta/2),\sin(\theta/2))\bigr).
\]
Przy zwiększeniu $\theta$ o $2\pi$ wektor jednostkowy zmienia znak,
który kompensuje identyfikacja $t\mapsto-t$. Wzór jest więc dobrze
określony. W każdym łuku wybieramy gładko ten wektor jednostkowy;
wtedy odwrotność odczytuje współczynnik $t$, co dowodzi gładkości
obu kierunków i liniowości na włóknach.
\end{przyklad}

\subsection*{Wiązka styczna}

Pokażemy, że zbiór $TM:=\bigsqcup_{p\in M}T_pM$ można w naturalny sposób wyposażyć w
strukturę rozmaitości gładkiej wymiaru $2n$, tak by rzutowanie $\pi\colon TM\to M$,
$\pi(v)=p$ dla $v\in T_pM$, było wiązką wektorową rzędu $n$.

Niech $\{(U_\alpha,\varphi_\alpha)\}$ będzie atlasem na $M$, $\varphi_\alpha=(x^1_\alpha,
\dots,x^n_\alpha)$. Definiujemy $\tilde U_\alpha:=\pi^{-1}(U_\alpha)\subseteq TM$ oraz
\[
\tilde\varphi_\alpha\colon\tilde U_\alpha\to\varphi_\alpha(U_\alpha)\times\R^n\subseteq
\R^{2n}, \qquad \tilde\varphi_\alpha\Big(\sum_{i=1}^nv^i\,\frac{\partial}{\partial
x^i_\alpha}\Big|_p\Big):=(\varphi_\alpha(p),\,v^1,\dots,v^n).
\]
Każde $\tilde\varphi_\alpha$ jest bijekcją na $\varphi_\alpha(U_\alpha)\times\R^n$ (dla
ustalonego $p$, współrzędne $v^i$ względem bazy $\{\partial/\partial x^i_\alpha|_p\}$
wyznaczają bijekcję liniową $T_pM\to\R^n$).

Obliczmy odwzorowanie przejścia $\tilde\varphi_\beta\circ\tilde\varphi_\alpha^{-1}$. Ze
wzoru na zmianę bazy pochodnych cząstkowych (reguła łańcucha zastosowana do funkcji
przejścia $\varphi_\beta\circ\varphi_\alpha^{-1}$):
\[
\frac{\partial}{\partial x^i_\alpha}\bigg|_p=\sum_{j=1}^n\frac{\partial x^j_\beta}
{\partial x^i_\alpha}(\varphi_\alpha(p))\;\frac{\partial}{\partial x^j_\beta}\bigg|_p
\]
(równość ta jest bezpośrednią konsekwencją definicji $\partial/\partial x^i_\alpha|_p$ jako
funkcjonału $f\mapsto\partial(f\circ\varphi_\alpha^{-1})/\partial x^i(\varphi_\alpha(p))$
oraz reguły łańcucha dla $f\circ\varphi_\alpha^{-1}=(f\circ\varphi_\beta^{-1})\circ
(\varphi_\beta\circ\varphi_\alpha^{-1})$). Zatem wektor $v=\sum_iv^i\,\partial/\partial
x^i_\alpha|_p$ ma we współrzędnych $\beta$ postać $w^j=\sum_i\frac{\partial x^j_\beta}
{\partial x^i_\alpha}(\varphi_\alpha(p))\,v^i$, czyli $w=D(\varphi_\beta\circ
\varphi_\alpha^{-1})(\varphi_\alpha(p))\cdot v$. Stąd
\[
\big(\tilde\varphi_\beta\circ\tilde\varphi_\alpha^{-1}\big)(x,v)=\Big(
(\varphi_\beta\circ\varphi_\alpha^{-1})(x),\ D(\varphi_\beta\circ\varphi_\alpha^{-1})(x)
\cdot v\Big).
\]
Najpierw określamy topologię na $TM$: zbiór $A\subset TM$
jest otwarty wtedy i tylko wtedy, gdy każdy
$\tilde\varphi_\alpha(A\cap\tilde U_\alpha)$ jest otwarty.
Zmiany współrzędnych powyżej są homeomorfizmami, a ich dziedziny
$\varphi_\alpha(U_\alpha\cap U_\beta)\times\R^n$ są otwarte.
Wynika stąd, że $\tilde U_\alpha$ są otwarte, a
$\tilde\varphi_\alpha$ są homeomorfizmami: obraz dowolnego otwartego
zbioru z jednej mapy jest otwarty także w każdej nakładającej się mapie.
Rzutowanie $\pi$ jest ciągłe, bo lokalnie ma postać $(x,v)\mapsto x$.
Powyższa zmiana współrzędnych jest gładka: pierwsza składowa jest gładka jako odwzorowanie przejścia
atlasu na $M$; druga jest liniowa w $v$ ze współczynnikami będącymi pochodnymi cząstkowymi
$\varphi_\beta\circ\varphi_\alpha^{-1}$, a więc gładkimi funkcjami $x$. Zatem $\{(\tilde
U_\alpha,\tilde\varphi_\alpha)\}$ jest atlasem na $TM$. Punkty o różnych
podstawach rozdzielają przeciwobrazy rozłącznych otoczeń w $M$;
dwa różne wektory w jednym włóknie rozdzielają kule w jednej mapie
$\tilde\varphi_\alpha$. Stąd $TM$ jest Hausdorffa. Z przeliczalności
bazy $M$ wybieramy przeliczalną rodzinę dziedzin map pokrywającą $M$.
Obrazy tych map w $TM$ mają przeliczalne bazy euklidesowe, których
przeciwobrazy tworzą przeliczalną bazę $TM$.

Tak skonstruowana rozmaitość $TM$, wraz z rzutowaniem $\pi$ i lokalnymi trywializacjami
$\Psi_\alpha=(\varphi_\alpha^{-1}\times\id)\circ\tilde\varphi_\alpha$, jest wiązką
wektorową rzędu $n$ nad $M$ — \emph{wiązką styczną}.

Analogicznie definiuje się \emph{wiązkę kostyczną} $T^*M:=\bigsqcup_{p\in M}T_p^*M$, z
mapami $\tilde\varphi_\alpha^*\big(\sum_i\xi_i\,dx^i_\alpha|_p\big):=(\varphi_\alpha(p),
\xi_1,\dots,\xi_n)$. Odwzorowanie przejścia oblicza się analogicznie — współczynniki
$\xi_i$ transformują się \emph{kowariantnie}: jeśli $\omega=\sum_i\xi_i^\alpha
dx^i_\alpha|_p=\sum_j\xi_j^\beta dx^j_\beta|_p$, to z reguły łańcucha $dx^i_\alpha|_p=
\sum_j\frac{\partial x^i_\alpha}{\partial x^j_\beta}(\varphi_\beta(p))\,dx^j_\beta|_p$,
skąd $\xi_j^\beta=\sum_i\xi_i^\alpha\frac{\partial x^i_\alpha}{\partial x^j_\beta}
(\varphi_\beta(p))$ — co ponownie daje gładkie odwzorowania przejścia. Dowód, że $T^*M$
jest wiązką wektorową rzędu $n$ nad $M$, dopełnia to samo sprawdzenie:
punkty o różnych podstawach rozdzielają przeciwobrazy otoczeń bazy,
punkty jednego włókna rozdzielają kule współrzędnych dualnych; bazy
euklidesowe w przeliczalnym pokryciu map dają drugi aksjomat przeliczalności.

\begin{definicja}
\emph{Polem wektorowym} na $M$ nazywamy gładką sekcję wiązki stycznej, tzn.\ gładkie
odwzorowanie $X\colon M\to TM$ z $\pi\circ X=\id_M$ (a więc $X(p)\in T_pM$ dla każdego
$p$). \emph{Polem kowektorowym} na $M$ nazywamy gładką sekcję wiązki kostycznej $T^*M$.
\end{definicja}
\section{Funkcje przejścia i operacje na wiązkach}

Lokalna trywialność pozwala opisywać wiązkę macierzami. Ta postać jest
szczególnie użyteczna przy konstruowaniu wiązek normalnych podrozmaitości.
Niech $\Psi_i:E|_{U_i}\to U_i\times\R^r$ będą trywializacjami i zapiszmy
\[
\Psi_j\Psi_i^{-1}(p,v)=(p,g_{ji}(p)v),\qquad
g_{ji}:U_i\cap U_j\longrightarrow\mathrm{GL}(r,\R).
\]
Macierze $g_{ji}$ są gładkie; na przecięciach spełniają
\[
g_{ii}=I,\qquad g_{ij}=g_{ji}^{-1},\qquad
g_{ki}=g_{kj}g_{ji}\quad\text{na }U_i\cap U_j\cap U_k.
\tag{1.1}
\]
Ostatnia równość porównuje bezpośrednie przejście z $i$ do $k$
z przejściem przez $j$; nazywamy ją warunkiem kocyklu.

\begin{twierdzenie}[Sklejanie wiązki z funkcji przejścia]
Niech $\{U_i\}$ będzie otwartym pokryciem rozmaitości $M$, a gładkie funkcje
$g_{ji}:U_i\cap U_j\to\mathrm{GL}(r,\R)$ spełniają (1.1). Wtedy istnieje
wiązka wektorowa rzędu $r$ o tych funkcjach przejścia.
\end{twierdzenie}
\begin{proof}
W sumie rozłącznej $\bigsqcup_i(U_i\times\R^r)$ utożsamiamy
$(p,v)_i$ z $(p,g_{ji}(p)v)_j$ dla $p\in U_i\cap U_j$. To rzeczywiście
relacja równoważności: pierwsza i druga tożsamość w (1.1) dają zwrotność
i symetrię, trzecia daje przechodniość. Niech $E$ będzie zbiorem klas i
$\pi([p,v]_i)=p$. Dla każdego $i$ odwzorowanie
\[
\Psi_i:\pi^{-1}(U_i)\longrightarrow U_i\times\R^r,\qquad
\Psi_i([p,v]_j)=(p,g_{ij}(p)v)
\]
jest dobrze określone: jeżeli $(p,v)_j\sim(p,w)_k$, to $w=g_{kj}v$,
a $g_{ik}w=g_{ik}g_{kj}v=g_{ij}v$. Odwrotność wysyła $(p,v)$ na
$[p,v]_i$. Nadajemy $E$ topologię, deklarując $A\subset E$ otwartym,
gdy każde $\Psi_i(A\cap\pi^{-1}(U_i))$ jest otwarte w $U_i\times\R^r$.
Zmiany trywializacji są homeomorfizmami między otwartymi podzbiorami,
więc każda $\pi^{-1}(U_i)$ jest otwarta, a $\Psi_i$ jest homeomorfizmem.
Łącząc $\Psi_i$ z mapami bazy $U_i$, otrzymujemy mapy do $\R^{\dim M+r}$.
Na przecięciach zmiana współrzędnych to dokładnie
$(p,v)\mapsto(p,g_{ji}(p)v)$, więc jest gładka i ma gładką odwrotność.
Włókna dostają działania liniowe przeniesione z $\R^r$; są niezależne
od wyboru trywializacji, bo $g_{ji}(p)$ jest liniowa. Otrzymujemy żądaną wiązkę.
Uzasadnijmy jeszcze dwie własności topologiczne. Punkty o różnych obrazach
w $M$ rozdzielamy przeciwobrazami rozłącznych otwartych otoczeń bazy;
punkty jednego włókna rozdzielamy we wspólnej trywializacji, wybierając
rozłączne otwarte kule wokół ich współrzędnych włóknowych. Stąd $E$ jest
Hausdorffa. Z drugiej przeliczalności $M$ wynika istnienie przeliczalnego
podpokrycia spośród $U_i$: dla każdego elementu przeliczalnej bazy,
który mieści się w jakimś $U_i$, wybieramy jeden taki $U_i$.
Te wybrane zbiory pokrywają $M$. Każde $U_i\times\R^r$
z podpokrycia ma
przeliczalną bazę, której obrazy przez $\Psi_i^{-1}$ razem dają
przeliczalną bazę $E$.
\end{proof}

\begin{figure}[ht]
\centering
\begin{tikzpicture}[>=Stealth,scale=1.02,line cap=round]
\node[font=\small] at (0,1.75) {ten sam $e\in E_p$, gdzie $p\in U_i\cap U_j$};
\fill[manifoldblue!12] (-3,0) ellipse (1.27 and .83);
\draw[manifoldblue!70!black,thick] (-3,0) ellipse (1.27 and .83);
\draw[->,manifoldblue!45!black,thin] (-4.08,0)--(-1.96,0);
\draw[->,manifoldblue!45!black,thin] (-3,-.62)--(-3,.63);
\draw[->,manifoldred,very thick] (-3,0)--(-2.35,.47)
 node[right,font=\small] {$v_i$};
\fill[manifoldgreen!12] (3,0) ellipse (1.27 and .83);
\draw[manifoldgreen!65!black,thick] (3,0) ellipse (1.27 and .83);
\draw[->,manifoldgreen!45!black,thin] (1.92,0)--(4.04,0);
\draw[->,manifoldgreen!45!black,thin] (3,-.62)--(3,.63);
\draw[->,manifoldred,very thick] (3,0)--(3.79,.29)
 node[right,font=\small] {$v_j$};
\draw[->,very thick,manifoldgold!75!black] (-1.45,0)--(1.45,0)
 node[midway,above,font=\small] {$g_{ji}(p)$};
\node[font=\small,manifoldblue!75!black] at (-3,-1.2)
 {$\Psi_i(E_p)=\{p\}\times\R^r$};
\node[font=\small,manifoldgreen!65!black] at (3,-1.2)
 {$\Psi_j(E_p)=\{p\}\times\R^r$};
\node[font=\small\bfseries,manifoldred!85!black] at (0,-1.85)
 {$v_j=g_{ji}(p)v_i$};
\end{tikzpicture}
\caption{Dwie trywializacje tego samego włókna $E_p$. Wektor $e$ ma współrzędne $v_i$ i $v_j$, związane macierzą przejścia $g_{ji}(p)$.}
\end{figure}

\begin{uwaga}[Grupa strukturalna]
Wybór baz włókien daje funkcje o wartościach w $\mathrm{GL}(r,\R)$; jest to
naturalna grupa strukturalna wiązki rzędu $r$. Jeżeli można dobrać
trywializacje tak, aby wszystkie macierze przejścia należały do podgrupy
$G\subseteq\mathrm{GL}(r,\R)$, mówimy o redukcji grupy strukturalnej do $G$.
Na przykład lokalne bazy ortonormalne wiązki z iloczynem skalarnym mają
macierze przejścia w $\mathrm{O}(r)$. Dla wiązki liniowej przejście $-1$
po okrążeniu okręgu wyraża skręcenie wstęgi Möbiusa.
\end{uwaga}

\begin{przyklad}[Sekcje w zapisie lokalnym]
Sekcja $s$ zapisuje się w trywializacji nad $U_i$ jako $s_i:U_i\to\R^r$. Na
przecięciu $s_j(p)=g_{ji}(p)s_i(p)$. Odwrotnie, rodzina gładkich funkcji
$s_i$ spełniająca te równości wyznacza jedną sekcję $s(p)=[p,s_i(p)]_i$;
niezależność od indeksu to dokładnie warunek zgodności. Dla wstęgi
Möbiusa można wybrać współrzędną na przedziale $[0,2\pi]$ i otrzymać
warunek $s(2\pi)=-s(0)$. Ciągła funkcja rzeczywista spełniająca ten
warunek musi mieć zero: jeśli wartość w początku jest niezerowa,
na końcu ma przeciwny znak, więc stosujemy twierdzenie o wartości
pośredniej; jeśli jest zerowa, teza jest natychmiastowa.
\end{przyklad}

\begin{definicja}[Suma i wiązka dualna]
\index{suma i wiazka dualna@Suma i wiązka dualna}
Dla wiązek $E,F$ nad tą samą bazą definiujemy włókna
$(E\oplus F)_p=E_p\oplus F_p$ oraz
$(E^*)_p=\mathrm{Hom}(E_p,\R)$. Przy macierzach przejścia $g_{ji}$
i $h_{ji}$ odpowiednie macierze to
\[
g_{ji}\oplus h_{ji},\qquad(g_{ji}^{-1})^{\mathsf T}.
\]
Ostatni wzór wymusza zachowanie wartości funkcjonału:
jeżeli $v_j=g_{ji}v_i$ oraz $\xi_i(v_i)=\xi_j(v_j)$, to
$\xi_j=(g_{ji}^{-1})^{\mathsf T}\xi_i$. Obie rodziny spełniają
warunek kocyklu, więc twierdzenie o sklejaniu tworzy z nich wiązki.
\end{definicja}



\section{Różniczka, immersje, submersje i zanurzenia}

\begin{definicja}
Niech $F\colon M\to N$ będzie odwzorowaniem gładkim, $p\in M$. \emph{Różniczką} $F$ w
punkcie $p$ nazywamy odwzorowanie liniowe $dF_p\colon T_pM\to T_{F(p)}N$ zdefiniowane (w
języku derywacji) wzorem
\[
(dF_p\,v)(g):=v(g\circ F), \qquad v\in T_pM,\ g\in C^\infty(N).
\]
\end{definicja}

\begin{fakt}
$dF_pv$ jest istotnie derywacją w $F(p)$, a $dF_p$ jest odwzorowaniem liniowym.
\end{fakt}
\begin{proof}
Liniowość względem $g$: oczywista z liniowości $v$ i liniowości $g\mapsto g\circ F$.
Reguła Leibniza: $(dF_pv)(gh)=v((gh)\circ F)=v((g\circ F)(h\circ F))=(g\circ F)(p)\,
v(h\circ F)+(h\circ F)(p)\,v(g\circ F)=g(F(p))(dF_pv)(h)+h(F(p))(dF_pv)(g)$, na mocy reguły
Leibniza dla $v$ zastosowanej do $g\circ F,h\circ F\in C^\infty(M)$. Zatem $dF_pv\in
T_{F(p)}N$. Liniowość $dF_p$ względem $v$ jest natychmiastowa z liniowości $v\mapsto
v(g\circ F)$ w $v$, dla każdego ustalonego $g$.
\end{proof}

\begin{twierdzenie}[Reguła łańcucha]
Jeśli $F\colon M\to N$ i $G\colon N\to P$ są gładkie, $p\in M$, to
\[
d(G\circ F)_p=dG_{F(p)}\circ dF_p\colon T_pM\to T_{G(F(p))}P.
\]
\end{twierdzenie}
\begin{proof}
Dla $v\in T_pM$ i $h\in C^\infty(P)$, z definicji różniczki:
\[
\big(d(G\circ F)_p\,v\big)(h)=v\big(h\circ(G\circ F)\big)=v\big((h\circ G)\circ F\big)=
(dF_pv)(h\circ G)=\big(dG_{F(p)}(dF_pv)\big)(h),
\]
gdzie w trzeciej równości zastosowaliśmy definicję $dF_p$ do $h\circ G\in C^\infty(N)$, a
w czwartej — definicję $dG_{F(p)}$ do $dF_pv\in T_{F(p)}N$. Zatem $d(G\circ F)_pv=
dG_{F(p)}(dF_pv)$ dla każdego $v$ — co dowodzi tezy.
\end{proof}

\begin{uwaga}[Skąd bierze się macierz Jacobiego?]
W mapach $(U,\varphi)$ o współrzędnych $(x^1,\ldots,x^m)$ wokół $p$ i
$(V,\psi)$ o współrzędnych $(y^1,\ldots,y^n)$ wokół $F(p)$ niech
$\widehat F=\psi\circ F\circ\varphi^{-1}$.
Aby znaleźć współczynniki $dF_p(\partial/\partial x^i|_p)$
w bazie $\partial/\partial y^a|_{F(p)}$, przykładamy ten
wektor do funkcji współrzędnej $y^a$ (lokalnie; w razie
potrzeby przedłużamy ją funkcją odcinającą). Z definicji
różniczki otrzymujemy
\[
 \bigl(dF_p(\partial/\partial x^i|_p)\bigr)(y^a)
 =\left.\frac{\partial \widehat F^{\,a}}{\partial x^i}
 \right|_{\varphi(p)}.
\]
Współczynniki te tworzą macierz Jacobiego
$D\widehat F(\varphi(p))$. \emph{Rzędem} $F$ w punkcie $p$
nazywamy rząd $dF_p$, równy rzędowi tej macierzy.
\end{uwaga}

\begin{definicja}
Niech $F\colon M\to N$ będzie gładkie, $\dim M=m$, $\dim N=n$. Mówimy, że $F$ jest w
punkcie $p$ \emph{immersją}, jeśli $dF_p$ jest injektywne, oraz \emph{submersją}, jeśli
$dF_p$ jest surjektywne. $F$ jest immersją (submersją), jeśli jest nią w każdym punkcie
$p\in M$. \emph{Zanurzeniem} (ang.\ \emph{embedding}) nazywamy immersję $F\colon M\to N$,
która dodatkowo jest homeomorfizmem na obraz $F(M)$ (z topologią podprzestrzeni $N$).
\end{definicja}

Poniższy przykład pokazuje, że sama różnowartościowa immersja nie musi być zanurzeniem —
warunek ,,homeomorfizm na obraz'' jest istotnie mocniejszy.

\begin{przyklad}[Krzywa ósemkowa]\label{prz:osemka}
Niech $\beta\colon(-\pi,\pi)\to\R^2$, $\beta(t):=(\sin2t,\ \sin t)$.

\emph{$\beta$ jest immersją.} $\beta'(t)=(2\cos2t,\cos t)$. Gdyby $\beta'(t_0)=0$, to
$\cos t_0=0$ (więc $t_0=\pm\pi/2$) oraz $\cos2t_0=0$; ale $\cos(\pm\pi)=-1\neq0$ —
sprzeczność. Zatem $\beta'(t)\neq0$ dla wszystkich $t$, czyli $d\beta_t$ jest injektywne w
każdym punkcie.

\emph{$\beta$ jest różnowartościowa.} Przypuśćmy $\beta(t_1)=\beta(t_2)$, tzn.\ $\sin t_1=
\sin t_2=:s$ oraz $\sin2t_1=\sin2t_2$. Jeśli $s=0$: na $(-\pi,\pi)$ równość $\sin t=0$
zachodzi jedynie dla $t=0$, więc $t_1=t_2=0$. Jeśli $s\neq0$: korzystając z $\sin2t_i=
2\sin t_i\cos t_i=2s\cos t_i$, równość $\sin2t_1=\sin2t_2$ daje (dzieląc przez $2s\neq0$)
$\cos t_1=\cos t_2$. Zatem $\sin t_1=\sin t_2$ i $\cos t_1=\cos t_2$ jednocześnie — co
oznacza $t_1-t_2\in2\pi\Z$; ponieważ $t_1,t_2\in(-\pi,\pi)$, jedyna możliwość to $t_1=t_2$.
W obu przypadkach $\beta$ jest różnowartościowa.

\emph{$\beta$ nie jest zanurzeniem.} $\lim_{t\to\pi^-}\beta(t)=(\sin2\pi,\sin\pi)=(0,0)$
oraz $\lim_{t\to-\pi^+}\beta(t)=(0,0)$ — oba końce krzywej zbiegają do tego samego punktu
$(0,0)=\beta(0)$. Odwzorowanie odwrotne $\beta^{-1}$ nie jest więc ciągłe w $(0,0)$: dla
ciągu $t_k\to\pi^-$, $\beta(t_k)\to(0,0)=\beta(0)$, ale $\beta^{-1}(\beta(t_k))=t_k\to\pi
\neq0=\beta^{-1}(0,0)$. Zatem $\beta$ jest różnowartościową immersją, która nie jest
homeomorfizmem na obraz — nie jest zanurzeniem.
\end{przyklad}

\begin{figure}[h]
\centering
\begin{tikzpicture}[>=Stealth,scale=1.4]
\draw[->] (-1.4,0) -- (1.4,0) node[right]{$x$};
\draw[->] (0,-1.3) -- (0,1.3) node[above]{$y$};
\draw[thick,blue,domain=-179:179,samples=150,smooth]
plot ({sin(2*\x)},{sin(\x)});
\fill (0,0) circle (1pt) node[below right]{$0$};
\end{tikzpicture}
\caption{Krzywa ósemkowa $\beta(t)=(\sin2t,\sin t)$: różnowartościowa immersja, która nie
jest zanurzeniem — oba końce ($t\to\pm\pi$) zbiegają do początku układu współrzędnych.}
\label{fig:osemka}
\end{figure}

\begin{uwaga}
Podobny, choć bardziej subtelny, przykład daje \emph{nawinięcie niewymierne} torusa
$T^2=\R^2/\Z^2$: krzywa $\gamma(t)=(t,\alpha t)+\Z^2$
dla $\alpha\in\R\setminus\Q$ jest injektywną immersją.
Rzeczywiście $\gamma(t)=\gamma(s)$ wymaga $t-s\in\Z$ i
$\alpha(t-s)\in\Z$, więc $t=s$; lokalna pochodna to $(1,\alpha)\ne0$.
Pokażmy wprost brak ciągłości odwrotności. Dzieląc $[0,1)$ na
$N$ równych przedziałów i rozpatrując części ułamkowe
$0,\alpha,\ldots,N\alpha$, z zasady szufladkowej otrzymujemy
$1\leq q_N\leq N$ oraz $p_N\in\Z$ takie, że
$|q_N\alpha-p_N|<1/N$. Możemy wybrać podciąg z $q_N\to\infty$:
dla skończenie wielu dodatnich $q$ odległości $q\alpha$ od $\Z$
mają dodatnie minimum. Wtedy $\gamma(q_N)\to\gamma(0)$,
ale $q_N$ nie dąży do zera, więc odwrotność nie jest ciągła.
Obraz jest również gęsty. Istotnie, dla dowolnie małego $\varepsilon>0$
istnieje niezerowe $\delta=q\alpha-p$ z $|\delta|<\varepsilon$.
Wielokrotności $\delta$ modulo $1$ tworzą zbiór
$\varepsilon$-gęsty w okręgu (kolejne punkty są oddalone o
$|\delta|$, a ostatnia luka jest nie większa).
Są to części ułamkowe całkowitych wielokrotności $\alpha$.
Zatem dla dowolnego $(a,b)+\Z^2$ można dobrać $k\in\Z$,
aby $\gamma(a+k)=(a,\alpha a+\alpha k)+\Z^2$ było dowolnie
blisko tego punktu.
\end{uwaga}

\subsection*{Twierdzenie o rzędzie stałym}

\begin{lemat}[Półciągłość rzędu z dołu]\label{lem:polciaglosc}
Niech $F\colon M\to N$ będzie gładkie. Zbiór $\{p\in M:\mathrm{rz}(dF_p)\geq r\}$ jest
otwarty w $M$, dla każdego $r$.
\end{lemat}
\begin{proof}
Warunek $\mathrm{rz}(dF_p)\geq r$ oznacza, w mapach lokalnych, że macierz Jacobiego $DF$ w
$\varphi(p)$ ma pewien minor stopnia $r$ różny od zera. Wyznacznik danego minora jest
funkcją ciągłą (wielomianem od wyrazów macierzy Jacobiego, gładkich funkcji punktu), więc
zbiór punktów, w których ten minor jest niezerowy, jest otwarty. Zbiór $\{\mathrm{rz}
(dF_p)\geq r\}$ jest sumą (po wszystkich wyborach minora stopnia $r$) takich zbiorów
otwartych, a więc otwarty.
\end{proof}

\begin{wniosek}\label{wn:stalyrzad}
Jeśli $F\colon M\to N$ jest immersją (submersją) w punkcie $p$, to $F$ jest immersją
(submersją) na pewnym otoczeniu $p$, o \emph{stałym} rzędzie na tym otoczeniu.
\end{wniosek}
\begin{proof}
Niech $m=\dim M$, $n=\dim N$. Jeśli $F$ jest immersją w $p$, $\mathrm{rz}(dF_p)=m$; ale
rząd $dF_q\colon\R^m\to\R^n$ jest zawsze $\leq m$, więc z Lematu~\ref{lem:polciaglosc}
($r=m$), $\{\mathrm{rz}(dF_q)\geq m\}=\{\mathrm{rz}(dF_q)=m\}$ jest otoczeniem otwartym
$p$, na którym $F$ ma stały rząd $m$. Przypadek submersji jest analogiczny z $r=n$.
\end{proof}

\begin{twierdzenie}[Twierdzenie o rzędzie stałym]\label{tw:rzad-staly}
Niech $F\colon M\to N$ będzie odwzorowaniem gładkim o stałym rzędzie $k$ na otoczeniu
punktu $p\in M$ ($\dim M=m$, $\dim N=n$). Wówczas istnieją mapy $(U,\varphi)$ wokół $p$
($\varphi(p)=0$) i $(V,\psi)$ wokół $F(p)$ ($\psi(F(p))=0$, $F(U)\subseteq V$), w których
\[
(\psi\circ F\circ\varphi^{-1})(u^1,\dots,u^m)=(u^1,\dots,u^k,0,\dots,0).
\]
\end{twierdzenie}

\begin{proof}
Pracujemy w mapach startowych, utożsamiając otoczenia $p$, $F(p)$ z otwartymi podzbiorami
$\R^m$, $\R^n$ (z $p=0$, $F(p)=0$), $F$ z gładkim odwzorowaniem $\hat F$ o stałym rzędzie
$k$ na całej (małej) dziedzinie.
Jeśli $k=0$, wszystkie pochodne $D\hat F$ znikają na małej
kuli wypukłej. Całkując wzdłuż odcinków od $0$, dostajemy
$\hat F(u)=\hat F(0)=0$, czyli żądaną postać. Dalej zakładamy $k>0$.

\emph{Krok 1: dyfeomorfizm lokalny przez przenumerowanie współrzędnych.} Ponieważ
$\mathrm{rz}(D\hat F(0))=k$, pewien minor stopnia $k$ macierzy $D\hat F(0)$ jest niezerowy;
przenumerowując współrzędne w $\R^m$, $\R^n$, możemy założyć, że jest to minor złożony z
pierwszych $k$ wierszy i kolumn — tzn., pisząc $\hat F=(\hat F^1,\dots,\hat F^n)$, że
$A:=\big(\partial\hat F^i/\partial u^j(0)\big)_{1\leq i,j\leq k}$ jest odwracalna.

Definiujemy $\Phi(u):=\big(\hat F^1(u),\dots,\hat F^k(u),u^{k+1},\dots,u^m\big)$. Macierz
$D\Phi(0)=\begin{pmatrix}A&B\\0&I_{m-k}\end{pmatrix}$ (blok $0$ w lewym dolnym rogu, bo
współrzędne $u^{k+1},\dots,u^m$ funkcji $\Phi$ nie zależą od $u^1,\dots,u^k$) ma
$\det D\Phi(0)=\det A\neq0$. Na mocy twierdzenia o odwzorowaniu odwrotnym (klasyczny fakt
analizy wielu zmiennych, przyjmowany tu bez dowodu), $\Phi$ jest dyfeomorfizmem pewnego
otoczenia $0$ na swój obraz.

\emph{Krok 2: przepisanie $\hat F$ przez $\Phi^{-1}$.} Niech $G:=\hat F\circ\Phi^{-1}$. Z
definicji $\Phi$, dla $w=\Phi(u)$: $G^i(w)=w^i$ dla $i=1,\dots,k$. Oznaczmy $h^i(w):=
G^i(w)$ dla $i=k+1,\dots,n$.

\emph{Krok 3: $h^i$ zależą tylko od $w^1,\dots,w^k$.} $G$ ma, jako złożenie $\hat F$ z
dyfeomorfizmem $\Phi^{-1}$, ten sam stały rząd $k$. Macierz $DG(w)=\begin{pmatrix}I_k&0\\
{*}&(\partial h^i/\partial w^j)_{i,j>k}\end{pmatrix}$. Gdyby $\partial h^{i_0}/\partial
w^{j_0}(w_0)\neq0$ dla pewnych $i_0,j_0>k$, minor $DG(w_0)$ złożony z wierszy $\{1,\dots,
k,i_0\}$ i kolumn $\{1,\dots,k,j_0\}$ miałby wyznacznik $\partial h^{i_0}/\partial
w^{j_0}(w_0)\neq0$ (macierz blokowo-trójkątna z blokiem $I_k$) — więc $\mathrm{rz}
(DG(w_0))\geq k+1$, sprzeczność ze stałością rzędu. Zatem $\partial h^i/\partial w^j\equiv0$
dla $i,j>k$. Zmniejszamy dziedzinę do produktu kul
$A\times B\subset\R^k\times\R^{m-k}$ o środku $0$.
Dla $(a,b)\in A\times B$ mamy
\[
 h^i(a,b)-h^i(a,0)
 =\int_0^1\sum_{j=k+1}^m b^j
       \frac{\partial h^i}{\partial w^j}(a,tb)\,dt=0.
\]
Zatem $h^i$ zależy tylko od $a=(w^1,\ldots,w^k)$;
dalej przez $h^i(a)$ rozumiemy $h^i(a,0)$.

\emph{Krok 4: druga zamiana współrzędnych, w $N$.} Definiujemy
\[
\Psi(y^1,\dots,y^n):=\big(y^1,\dots,y^k,\,y^{k+1}-h^{k+1}(y^1,\dots,y^k),\,\dots,\,
y^n-h^n(y^1,\dots,y^k)\big).
\]
Odwzorowanie $\Psi$ ma jawną odwrotną $\Psi^{-1}(z)=(z^1,\dots,z^k,z^{k+1}+h^{k+1}(z^1,
\dots,z^k),\dots,z^n+h^n(z^1,\dots,z^k))$ (co sprawdza się bezpośrednim podstawieniem), a
więc jest dyfeomorfizmem swojej dziedziny na obraz.

Dla $w$ blisko $0$:
\[
(\Psi\circ G)(w)=\Psi\big(w^1,\dots,w^k,h^{k+1}(w^{\leq k}),\dots,h^n(w^{\leq k})\big)=
(w^1,\dots,w^k,0,\dots,0),
\]
bo $h^i(w^{\leq k})-h^i(w^{\leq k})=0$. Podstawiając $\psi:=\Psi$, $\varphi:=\Phi$ (złożone
z przenumerowaniem z Kroku 1), otrzymujemy $\psi\circ F\circ\varphi^{-1}=\Psi\circ G=
(u^1,\dots,u^k,0,\dots,0)$, co kończy dowód.
\end{proof}

\begin{wniosek}[Postać lokalna immersji]\label{wn:immersja-lokalnie}
Jeśli $F\colon M\to N$ jest immersją w $p$ ($m\leq n$), to istnieją mapy wokół $p$, $F(p)$,
w których $F(u^1,\dots,u^m)=(u^1,\dots,u^m,0,\dots,0)$.
\end{wniosek}
\begin{proof}
Na mocy Wniosku~\ref{wn:stalyrzad}, $F$ ma stały rząd $m$ na otoczeniu $p$; stosujemy
Twierdzenie~\ref{tw:rzad-staly} z $k=m$.
\end{proof}

\begin{wniosek}[Postać lokalna submersji]\label{wn:submersja-lokalnie}
Jeśli $F\colon M\to N$ jest submersją w $p$ ($m\geq n$), to istnieją mapy wokół $p$, $F(p)$,
w których $F(u^1,\dots,u^m)=(u^1,\dots,u^n)$.
\end{wniosek}
\begin{proof}
Na mocy Wniosku~\ref{wn:stalyrzad}, $F$ ma stały rząd $n$ na otoczeniu $p$; stosujemy
Twierdzenie~\ref{tw:rzad-staly} z $k=n$ (teza redukuje się do $(u^1,\dots,u^n,0,\dots,0)\in
\R^n$, czyli po prostu $(u^1,\dots,u^n)$).
\end{proof}

\begin{przyklad}
Rzutowanie kanoniczne $\mathrm{pr}_1\colon M\times N\to M$ jest submersją: w mapach
produktowych, $\mathrm{pr}_1$ przyjmuje postać rzutowania liniowego $\R^m\times\R^n\to
\R^m$, którego różniczka (ono samo) jest surjektywna. Podobnie, rzutowanie wiązki
stycznej $\pi\colon TM\to M$ jest submersją — w mapach $\tilde\varphi_\alpha$ z Sekcji~\ref{sec:wiazki-wektorowe},
$\pi$ przyjmuje postać $(x,v)\mapsto x$.
\end{przyklad}
\section{Podrozmaitości}

\begin{definicja}
Niech $M$ będzie rozmaitością gładką wymiaru $n$, $S\subseteq M$ podzbiorem, $k\leq n$.
Mapa $(U,\varphi)$ na $M$ jest \emph{mapą plasterkową} (dostosowaną do $S$) rzędu $k$,
jeżeli
\[
\varphi(U\cap S)=\varphi(U)\cap(\R^k\times\{0\})\subseteq\R^k\times\R^{n-k}\cong\R^n.
\]
Mówimy, że $S$ jest \emph{podrozmaitością zanurzoną} wymiaru $k$ w $M$, jeśli każdy punkt
$p\in S$ posiada mapę plasterkową rzędu $k$ wokół $p$.
\end{definicja}

\begin{figure}[h]
\centering
\begin{tikzpicture}[>=Stealth,scale=1]
\draw[thick] (0,0) rectangle (5,3.5);
\node at (2.5,3.8) {$\varphi(U)\subseteq\R^k\times\R^{n-k}$};
\draw[thick,blue] (0,1.75) -- (5,1.75);
\node[blue] at (6.05,1.75) {$\R^k\times\{0\}$};
\draw[->,thick] (2.5,-0.3) -- (2.5,-1.6) node[midway,right]{$\varphi^{-1}$};
\node at (2.5,-2.1) {$U\cap S=\varphi^{-1}\big(\varphi(U)\cap(\R^k\times\{0\})\big)$};
\end{tikzpicture}
\caption{Mapa plasterkowa: w mapie $\varphi$, przecięcie $U\cap S$ odpowiada dokładnie
przecięciu $\varphi(U)$ z podprzestrzenią $\R^k\times\{0\}$.}
\label{fig:plasterek}
\end{figure}

\begin{stwierdzenie}\label{stw:podrozm-struktura}
Jeśli $S\subseteq M$ jest podrozmaitością zanurzoną wymiaru $k$, to $S$ (z topologią
podprzestrzeni) jest rozmaitością topologiczną wymiaru $k$, posiadającą strukturę gładką,
względem której inkluzja $\iota\colon S\hookrightarrow M$ jest immersją; co więcej, $\iota$
jest zanurzeniem.
\end{stwierdzenie}
\begin{proof}
Wyposażamy $S$ w topologię podprzestrzeni. Ponieważ podprzestrzeń
przestrzeni Hausdorffa jest Hausdorffa, $S$ spełnia pierwszy warunek.
Jeśli $\mathcal B$ jest przeliczalną bazą $M$, to
$\{B\cap S:B\in\mathcal B\}$ jest przeliczalną bazą $S$:
dla otwartego $O\cap S$ i $p\in O\cap S$ znajdziemy
$B\in\mathcal B$ z $p\in B\subset O$, a wtedy
$p\in B\cap S\subset O\cap S$.

Ustalmy $p\in S$ i mapę plasterkową $(U,\varphi)$ przy $p$.
Ponieważ $U$ jest otwarte w $M$, zbiór $U_S=U\cap S$ jest otwarty
w $S$. Oznaczmy $W=\varphi(U)$, otwarty w $\R^n$.
Odwzorowanie $\varphi_S=\operatorname{pr}_1\circ\varphi|_{U_S}$
ma obraz $W_S=\{u\in\R^k:(u,0)\in W\}$, otwarty w $\R^k$ jako
przeciwobraz $W$ przez ciągłe zanurzenie $u\mapsto(u,0)$.
Odwrotność $u\mapsto\varphi^{-1}(u,0)$ jest ciągła; zatem
$\varphi_S:U_S\to W_S$ jest homeomorfizmem. To dowodzi lokalnej
euklidesowości $S$.

Niech teraz $(V,\psi)$ będzie drugą mapą plasterkową. Na
$\varphi_S(U\cap V\cap S)$ przejście ma wzór
\[
\psi_S\circ\varphi_S^{-1}(u)
=\operatorname{pr}_1\bigl((\psi\circ\varphi^{-1})(u,0)\bigr).
\]
Zanurzenie $u\mapsto(u,0)$, przejście $\psi\circ\varphi^{-1}$
i rzutowanie $\operatorname{pr}_1$ są gładkie, więc złożenie jest
gładkie. Zamieniając $\varphi$ i $\psi$, otrzymujemy gładką
odwrotność. Mapy $\varphi_S$ tworzą więc atlas gładki na $S$.

W mapach $\varphi_S$ i $\varphi$ inkluzja $\iota:S\to M$ ma postać
$u\mapsto(u,0)$, a jej różniczka wysyła $v$ na $(v,0)$ i jest
injektywna. Inkluzja jest więc immersją. Jako odwzorowanie
$S\to\iota(S)$ z \emph{topologią podprzestrzeni}, $\iota$ jest
homeomorfizmem już z wyboru topologii na $S$ (tożsamość na
zbiorze punktów). Zatem jest zanurzeniem.
\end{proof}

\begin{przyklad}[Wiązka normalna jako iloraz]
Niech $S\subset M$ będzie podrozmaitością zanurzoną. Poprzez różniczkę
inkluzji utożsamiamy $T_pS$ z podprzestrzenią $T_pM$. Wiązka $TS$ jest
podwiązką $TM|_S$. Włókna wiązki normalnej to ilorazy
\[
\nu_pS=T_pM/T_pS,\qquad \nu S=TM|_S/TS.
\]
Jej rząd wynosi $\dim M-\dim S$. W mapie plasterkowej
o współrzędnych $(x,y)$ klasa wektora stycznego
ma jednoznaczny reprezentant w kierunkach $\partial/\partial y^a$.
Sprawdźmy zmiany współrzędnych. Przejście między dwiema mapami
plasterkowymi ma postać $(x,y)\mapsto(a(x,y),b(x,y))$, przy czym
$b(x,0)=0$. Różniczkując tę tożsamość po $x$, dostajemy
$D_xb(x,0)=0$, więc różniczka przejścia na $S$ ma postać blokową
\[
 \begin{pmatrix} A&B\\0&C\end{pmatrix}.
\]
Cała macierz jest odwracalna, zatem kwadratowy blok $C$ jest
odwracalny. Klasa wektora o składowej poprzecznej $v$ zmienia
współrzędne na $Cv$, niezależnie od składowej stycznej.
Macierze $C$ są gładkie i spełniają warunek kocyklu, bo różniczki
przejść się składają. Twierdzenie o sklejaniu daje więc wiązkę
o zadanych włóknach ilorazowych. Konstrukcja jest kanoniczna bez metryki; po wybraniu
metryki można identyfikować klasy z wektorami prostopadłymi do $T_pS$.
\end{przyklad}

Poniższe twierdzenie i jego następstwa pokazują, w jaki sposób podrozmaitości powstają ,,w
praktyce'' — jako obrazy zanurzeń oraz jako przeciwobrazy wartości regularnych.

\begin{twierdzenie}[Twierdzenie o obrazie]\label{tw:obraz}
Jeśli $F\colon S\to M$ jest zanurzeniem ($\dim S=k\leq n=\dim M$), to $F(S)$ jest
podrozmaitością zanurzoną wymiaru $k$ w $M$, a $F\colon S\to F(S)$ (z topologią
podprzestrzeni) jest dyfeomorfizmem.
\end{twierdzenie}
\begin{proof}
Ustalmy $q=F(p_0)\in F(S)$. Ponieważ $F$ jest immersją w $p_0$, na mocy
Wniosku~\ref{wn:immersja-lokalnie} istnieją mapy $(U,\varphi)$ wokół $p_0$ w $S$ i
$(V,\psi)$ wokół $q$ w $M$, w których $F$ przyjmuje postać $u\mapsto(u,0)$.
Uzasadnijmy potrzebne dopasowanie dziedzin: wybieramy małe otwarte
kule $A\subset\R^k$, $B\subset\R^{n-k}$ o środku $0$, tak aby
$A\subset\varphi(U)$ i $A\times B\subset\psi(V)$.
Zastępujemy $U$ przez $\varphi^{-1}(A)$, a $V$ przez
$\psi^{-1}(A\times B)$. Wtedy rzeczywiście
$\psi(F(U))=A\times\{0\}=\psi(V)\cap(\R^k\times\{0\})$.
Zatem $(V,\psi)$ jest ,,prawie'' mapą
plasterkową dla $F(S)$, z tą różnicą, że warunek dotyczy $F(U)$, a nie od razu $F(S)\cap V$.

W tym miejscu wykorzystujemy założenie, że $F$ jest \emph{zanurzeniem}, nie tylko
różnowartościową immersją (por.\ Przykład~\ref{prz:osemka}): ponieważ $F\colon S\to F(S)$
jest homeomorfizmem na obraz, zbiór $F(U)$ jest otwarty w $F(S)$, tzn.\ $F(U)=F(S)\cap W$
dla pewnego otwartego $W\subseteq M$. Zastępując $V$ przez $V':=V\cap W$, otrzymujemy
$F(S)\cap V'=F(U)\cap V'$, a
\begin{align*}
\psi(F(S)\cap V')
 &=\psi(F(U)\cap V')\\
 &=\psi(F(U))\cap\psi(V')\\
 &=\big(\psi(V)\cap(\R^k\times\{0\})\big)\cap\psi(V')\\
 &=\psi(V')\cap(\R^k\times\{0\}).
\end{align*}
Zatem $(V',\psi|_{V'})$ jest mapą plasterkową rzędu $k$ dla $F(S)$ wokół $q$. Ponieważ
$q\in F(S)$ był dowolny, $F(S)$ jest podrozmaitością zanurzoną wymiaru $k$.

Wreszcie, $F\colon S\to F(S)$ jest bijekcją gładką (w mapach plasterkowych powyżej, $F$
przyjmuje postać identyczności), której odwrotność $F^{-1}\colon F(S)\to S$ jest gładka z
tego samego powodu. Zatem $F\colon S\to F(S)$ jest dyfeomorfizmem.
\end{proof}

\begin{definicja}
Niech $F\colon M\to N$ będzie gładkie, $q\in N$. Punkt $q$ jest \emph{wartością regularną}
$F$, jeśli dla każdego $p\in F^{-1}(q)$, $F$ jest submersją w $p$ (jeśli $F^{-1}(q)=
\emptyset$, $q$ jest — pusto — wartością regularną).
\end{definicja}

\begin{twierdzenie}[Twierdzenie o przeciwobrazie wartości regularnej]\label{tw:przeciwobraz}
Niech $F\colon M\to N$ będzie gładkie, $\dim M=m$, $\dim N=n$, i niech $q\in N$ będzie
wartością regularną $F$ z $F^{-1}(q)\neq\emptyset$. Wówczas $F^{-1}(q)$ jest
podrozmaitością zanurzoną wymiaru $m-n$ w $M$, a dla każdego $p\in F^{-1}(q)$,
\[
T_p\big(F^{-1}(q)\big)=\ker(dF_p)\subseteq T_pM.
\]
\end{twierdzenie}
\begin{proof}
Niech $p\in F^{-1}(q)$. Ponieważ $F$ jest submersją w $p$, na mocy
Wniosku~\ref{wn:submersja-lokalnie} istnieją mapy $(U,\varphi)$ wokół $p$ ($\varphi(p)=0$)
i $(V,\psi)$ wokół $q$ ($\psi(q)=0$), w których $F$ przyjmuje postać $\psi\circ F\circ
\varphi^{-1}(u^1,\dots,u^m)=(u^1,\dots,u^n)$. W tych mapach,
\[
\varphi(U\cap F^{-1}(q))=\{u\in\varphi(U):(u^1,\dots,u^n)=0\}=\varphi(U)\cap(\{0\}\times
\R^{m-n}).
\]
Po przenumerowaniu współrzędnych w $\R^m$ (zamiana miejscami pierwszych $n$ z pozostałymi
$m-n$), $(U,\varphi)$ staje się mapą plasterkową rzędu $m-n$ dla $F^{-1}(q)$ wokół $p$.
Ponieważ $p\in F^{-1}(q)$ był dowolny, $F^{-1}(q)$ jest podrozmaitością zanurzoną wymiaru
$m-n$.

W drugiej części niech $\iota\colon F^{-1}(q)\hookrightarrow M$
będzie inkluzją. Jest ona immersją na mocy
Stwierdzenia~\ref{stw:podrozm-struktura}. Ponieważ $F\circ\iota$ jest
odwzorowaniem stale równym $q$, jego różniczka znika: $dF_p\circ d\iota_p=0$ (reguła
łańcucha). Zatem $d\iota_p(T_pF^{-1}(q))\subseteq\ker(dF_p)$; utożsamiając (poprzez
injektywną immersję $d\iota_p$) $T_pF^{-1}(q)$ z jego obrazem w $T_pM$, mamy $T_pF^{-1}(q)
\subseteq\ker(dF_p)$. Z rachunku wymiarów: $\dim T_pF^{-1}(q)=m-n$ (z pierwszej części
dowodu), a $\dim\ker(dF_p)=m-\mathrm{rz}(dF_p)=m-n$ (bo $dF_p$ jest surjektywne na $T_qN$
wymiaru $n$, z definicji submersji). Zatem obie podprzestrzenie $T_pM$ mają ten sam,
skończony wymiar $m-n$, a jedna zawiera się w drugiej — więc są równe.
\end{proof}

\begin{wniosek}[Twierdzenie o poziomicy]
Jeśli $\dim M=n$, $f\colon M\to\R$ jest gładka, $c\in\R$ jest wartością regularną $f$ (tzn.\
$df_p\neq0$ dla każdego $p\in f^{-1}(c)$) i $f^{-1}(c)\neq\emptyset$, to $f^{-1}(c)$ jest
podrozmaitością zanurzoną wymiaru $n-1$ w $M$ (,,hiperpowierzchnią''), z $T_pf^{-1}(c)=
\ker(df_p)$.
\end{wniosek}
\begin{proof}
Szczególny przypadek Twierdzenia~\ref{tw:przeciwobraz} z $N=\R$: submersywność $f$ w
$p\in f^{-1}(c)$ to dokładnie warunek $df_p\neq0$.
\end{proof}

\begin{przyklad}
Niech $f\colon\R^{n+1}\to\R$, $f(x)=\|x\|^2$. Wtedy $df_x(v)=2\langle x,v\rangle$, niezerowe
dla $x\neq0$. Ponieważ $S^n=f^{-1}(1)$ i $0\notin f^{-1}(1)$, $1$ jest wartością regularną
— na mocy Wniosku powyżej, $S^n$ jest podrozmaitością zanurzoną wymiaru $n$ w $\R^{n+1}$
(zgodnie z tym, co wiedzieliśmy już z Przykładu~\ref{prz:sfera}; twierdzenie o poziomicy
daje jednak krótszy dowód gładkości struktury na $S^n$, choć nie podaje jawnych map).
\end{przyklad}

\begin{przyklad}[Grupa ortogonalna]
Niech $F\colon M_n(\R)\to\mathrm{Sym}(n)$, $F(A):=A^TA$, gdzie $\mathrm{Sym}(n)\cong
\R^{n(n+1)/2}$ oznacza przestrzeń macierzy symetrycznych. Różniczka: dla $H\in M_n(\R)$,
\[
dF_A(H)=\frac{d}{dt}\Big|_{t=0}(A+tH)^T(A+tH)=A^TH+H^TA.
\]
Dla $A\in\mathrm{O}(n):=F^{-1}(I)=\{A:A^TA=I\}$, $dF_A$ jest surjektywne na
$\mathrm{Sym}(n)$: dla symetrycznej $S$, kładąc $H:=\tfrac12AS$,
\[
dF_A(H)=A^T\big(\tfrac12AS\big)+\big(\tfrac12AS\big)^TA=\tfrac12(A^TA)S+\tfrac12S^TA^TA=
\tfrac12S+\tfrac12S=S
\]
(korzystając z $A^TA=I$, $S^T=S$). Zatem $I$ jest wartością regularną $F$, więc
$\mathrm{O}(n)=F^{-1}(I)$ jest podrozmaitością zanurzoną $M_n(\R)\cong\R^{n^2}$ wymiaru
$n^2-n(n+1)/2=n(n-1)/2$. Co więcej, $T_I\mathrm{O}(n)=\ker(dF_I)=\{H:H+H^T=0\}=:
\mathfrak{so}(n)$ — przestrzeń macierzy antysymetrycznych, wymiaru $n(n-1)/2$ (zgodnie z
powyższym rachunkiem wymiaru).
\end{przyklad}

\begin{uwaga}
Powyższy przykład jest pierwszym spotkaniem z fundamentalnym zjawiskiem: $\mathrm{O}(n)$
jest nie tylko rozmaitością, ale i grupą (względem mnożenia macierzy), w której mnożenie i
branie odwrotności są gładkie — tzw.\ \emph{grupą Liego} — a $T_I\mathrm{O}(n)=
\mathfrak{so}(n)$ jest jej \emph{algebrą Liego}.
\end{uwaga}

\begin{przyklad}[Torus jako regularna poziomica]
Niech $R>r>0$, $\rho=\sqrt{x^2+y^2}$ oraz
$\Omega=\{(x,y,z)\in\R^3:\rho>0\}$. Na tym otwartym zbiorze
$f(x,y,z)=(\rho-R)^2+z^2$ jest gładka. Poziom $f=r^2$
leży w $\Omega$, gdyż $|\rho-R|\leq r<R$ daje $\rho\geq R-r>0$.
Różniczka w punkcie $\Omega$ ma postać
\[
df=2(\rho-R)\left(\frac{x}{\rho}\,dx+
\frac{y}{\rho}\,dy\right)+2z\,dz.
\]
Jeżeli $df=0$, to $z=0$ oraz $\rho=R$: ponieważ $(x,y)\ne(0,0)$,
nie mogą jednocześnie znikać współczynniki przy $dx,dy$, chyba że
$\rho-R=0$. Lecz wtedy $f=0$, a więc taki punkt nie leży na
poziomie $r^2>0$. Zatem $r^2$ jest wartością regularną i torus
obrotowy $f^{-1}(r^2)$ jest gładką powierzchnią. Ponadto
$T_pT^2=\ker df_p$: wektor $(a,b,c)$ jest styczny wtedy i tylko
wtedy, gdy $(\rho-R)(xa+yb)/\rho+zc=0$.
\end{przyklad}
\clearpage
